\documentclass[12pt,oneside]{book}

% ============================================================
% DERECHOS REALES — MANUAL DE ESTUDIO INTEGRADO
% Documento autónomo: no depende de plantillas ni archivos externos.
% ============================================================

% ------------------------------------------------------------
% METADATOS (únicamente datos presentes en los apuntes originales)
% ------------------------------------------------------------
\newcommand{\universidad}{Universidad de Buenos Aires (UBA)}
\newcommand{\facultad}{Facultad de Derecho}
\newcommand{\carrera}{Derecho Civil}
\newcommand{\materia}{Derechos Reales}
\newcommand{\titulo}{Derechos Reales}
\newcommand{\subtitulo}{Manual de estudio integrado}
\newcommand{\catedra}{Cátedra Prof. Farina}
\newcommand{\autor}{Isaac Edgar Camacho Ocampo}
\newcommand{\anio}{2026}

% ------------------------------------------------------------
% PAQUETES
% ------------------------------------------------------------
\usepackage[utf8]{inputenc}
\usepackage[T1]{fontenc}
\usepackage[spanish,es-nodecimaldot]{babel}

\usepackage[
    top=2.5cm,
    bottom=2.5cm,
    left=3cm,
    right=2.5cm,
    headheight=15pt
]{geometry}

\usepackage{amsmath}
\usepackage{amssymb}
\usepackage{mathtools}

\usepackage{graphicx}
\usepackage{float}
\usepackage{caption}
\usepackage{subcaption}

\usepackage{booktabs}
\usepackage{array}
\usepackage{longtable}
\usepackage{tabularx}

\usepackage{enumitem}

\usepackage{xcolor}
\usepackage[most]{tcolorbox}

\usepackage{fancyhdr}
\usepackage{ragged2e}

\usepackage[
    colorlinks=true,
    linkcolor=blue!50!black,
    citecolor=green!40!black,
    urlcolor=blue!60!black,
    pdftitle={Derechos Reales. Manual de estudio integrado},
    pdfauthor={Isaac Edgar Camacho Ocampo},
    pdfsubject={Derechos Reales},
    bookmarks=true
]{hyperref}

% ------------------------------------------------------------
% CONFIGURACIÓN GENERAL
% ------------------------------------------------------------
\setcounter{tocdepth}{2}
\setcounter{secnumdepth}{2}

\setlength{\parindent}{0pt}
\setlength{\parskip}{0.5em}

\setlist[itemize]{topsep=0.3em, itemsep=0.2em, parsep=0pt}
\setlist[enumerate]{topsep=0.3em, itemsep=0.2em, parsep=0pt}

\clubpenalty=10000
\widowpenalty=10000
\raggedbottom

\renewcommand{\arraystretch}{1.25}
\captionsetup{font=small, labelfont=bf}

% ------------------------------------------------------------
% CAJAS ACADÉMICAS (uso semántico)
% ------------------------------------------------------------
\newtcolorbox{definicion}[1]{
    enhanced, breakable,
    title={Definición: #1}, fonttitle=\bfseries\small,
    colback=blue!3, colframe=blue!50!black,
    boxrule=0.7pt, arc=1.5mm, left=6pt, right=6pt, top=4pt, bottom=4pt
}
\newtcolorbox{conceptoclave}[1]{
    enhanced, breakable,
    title={Concepto clave: #1}, fonttitle=\bfseries\small,
    colback=teal!4, colframe=teal!50!black,
    boxrule=0.7pt, arc=1.5mm, left=6pt, right=6pt, top=4pt, bottom=4pt
}
\newtcolorbox{importante}{
    enhanced, breakable,
    title={Importante}, fonttitle=\bfseries\small,
    colback=yellow!8, colframe=orange!70!black,
    boxrule=0.7pt, arc=1.5mm, left=6pt, right=6pt, top=4pt, bottom=4pt
}
\newtcolorbox{ejemplo}[1]{
    enhanced, breakable,
    title={Ejemplo: #1}, fonttitle=\bfseries\small,
    colback=green!4, colframe=green!40!black,
    boxrule=0.7pt, arc=1.5mm, left=6pt, right=6pt, top=4pt, bottom=4pt
}
\newtcolorbox{regla}[1]{
    enhanced, breakable,
    title={Regla / texto normativo: #1}, fonttitle=\bfseries\small,
    colback=purple!3, colframe=purple!50!black,
    boxrule=0.7pt, arc=1.5mm, left=6pt, right=6pt, top=4pt, bottom=4pt
}
\newtcolorbox{ejercicio}[1]{
    enhanced, breakable,
    title={Ejercicio: #1}, fonttitle=\bfseries\small,
    colback=gray!5, colframe=gray!60!black,
    boxrule=0.7pt, arc=1.5mm, left=6pt, right=6pt, top=4pt, bottom=4pt
}
\newtcolorbox{nota}{
    enhanced, breakable,
    title={Nota}, fonttitle=\bfseries\small,
    colback=white, colframe=gray!55,
    boxrule=0.5pt, arc=1.5mm, left=6pt, right=6pt, top=4pt, bottom=4pt
}
\newtcolorbox{advertencia}{
    enhanced, breakable,
    title={Advertencia}, fonttitle=\bfseries\small,
    colback=red!4, colframe=red!55!black,
    boxrule=0.7pt, arc=1.5mm, left=6pt, right=6pt, top=4pt, bottom=4pt
}

% ------------------------------------------------------------
% CUADRO CORNELL (tabla real con tabularx)
%   \cornell{título}{filas \cq{pregunta}{respuesta}...}{resumen}
% ------------------------------------------------------------
\newcommand{\cq}[2]{\textbf{#1} & #2 \\[0.45em]}
\newcommand{\cornell}[3]{%
\par\bigskip\noindent\textbf{\large Cuadro Cornell: #1}\par\smallskip
\begin{center}
\small
\renewcommand{\arraystretch}{1.15}
\begin{tabularx}{\textwidth}{>{\raggedright\arraybackslash}p{0.27\textwidth}>{\raggedright\arraybackslash}X}
\toprule
\textbf{Preguntas / palabras clave} & \textbf{Notas / respuestas} \\
\midrule
#2
\midrule
\multicolumn{2}{p{\dimexpr\textwidth-2\tabcolsep\relax}}{\textbf{Resumen.} #3} \\
\bottomrule
\end{tabularx}
\end{center}
}

% ------------------------------------------------------------
% ENCABEZADOS Y PIES
% ------------------------------------------------------------
\pagestyle{fancy}
\fancyhf{}
\fancyhead[L]{\small\nouppercase{\leftmark}}
\fancyhead[R]{\small\materia}
\fancyfoot[C]{\thepage}
\renewcommand{\headrulewidth}{0.4pt}
\renewcommand{\footrulewidth}{0pt}

\fancypagestyle{plain}{
    \fancyhf{}
    \fancyfoot[C]{\thepage}
    \renewcommand{\headrulewidth}{0pt}
}

% ============================================================
% DOCUMENTO
% ============================================================
\begin{document}
\frontmatter

% ------------------------------------------------------------
% PORTADA
% ------------------------------------------------------------
\begin{titlepage}
\thispagestyle{empty}
\begin{center}
\vspace*{1cm}
{\large \textsc{\universidad}}\\[0.2cm]
{\large \textsc{\facultad}}\\[0.2cm]
{\large \carrera}

\vspace{3cm}

{\Huge \textbf{\materia}}

\vspace{1cm}

{\LARGE \subtitulo}

\vspace{0.4cm}

{\normalsize \catedra}

\vspace{1.2cm}

\rule{0.70\textwidth}{0.5pt}

\vspace{0.4cm}

{\large Teoría general, relaciones de poder, dominio, condominio,\\ propiedad horizontal y nuevas formas de propiedad}

\vspace{0.4cm}

\rule{0.70\textwidth}{0.5pt}

\vfill

{\large \autor}

\vspace{0.3cm}

{\large \anio}
\end{center}

\vfill

\begin{flushleft}
\small\textit{Este es un modesto aporte a los estudiantes de «Derechos Reales»; no pretende sustituir el material oficial de la materia.}
\end{flushleft}
\end{titlepage}

% ------------------------------------------------------------
% ÍNDICE ÚNICO
% ------------------------------------------------------------
\tableofcontents

\mainmatter

% ------------------------------------------------------------
% INTRODUCCIÓN GENERAL
% ------------------------------------------------------------
\chapter*{Introducción general}
\addcontentsline{toc}{chapter}{Introducción general}
\markboth{Introducción general}{Introducción general}

Este manual reúne, en un único desarrollo, el estudio de los derechos reales en el Código Civil y Comercial (CCyC). Los derechos reales organizan el modo en que una sociedad reconoce y protege la relación de poder de las personas con las cosas: quién puede usarlas, gozarlas y disponer de ellas, frente a quiénes puede hacer valer ese poder y con qué límites. El recorrido va de lo general a lo particular.

La \textbf{Parte I} establece los fundamentos: qué distingue a un derecho real de un derecho personal, cómo se ubican los derechos reales dentro del Código, cuáles son sus principios estructurales y cuáles son los catorce derechos reales que la ley reconoce.

La \textbf{Parte II} estudia cómo nacen, se hacen conocer, se transmiten y se extinguen los derechos reales, con especial atención a la teoría del título y modo.

La \textbf{Parte III} se ocupa de las relaciones de poder de hecho entre la persona y la cosa (posesión y tenencia), que explican la tradición, la publicidad posesoria y la adquisición legal de derechos sobre cosas muebles.

A partir de la \textbf{Parte IV} se aplican esas bases a cada derecho real en particular: el dominio y sus límites (Parte IV); el condominio y la medianería (Parte V); la propiedad horizontal (Parte VI); y los conjuntos inmobiliarios, el tiempo compartido y los cementerios privados (Parte VII). Al final se incluyen las referencias mencionadas en las fuentes y un apéndice de ejercicios de evaluación.

Cada capítulo cierra con un cuadro Cornell para la recuperación activa de los conceptos centrales.

\begin{nota}
Siempre que los apuntes originales citan artículos de forma inconsistente entre sí, el texto conserva ambas referencias y lo señala mediante un comentario \texttt{TODO} en el código fuente, sin resolverlo con información externa. Las abreviaturas utilizadas son CCyC (Código Civil y Comercial de la Nación) y CN (Constitución Nacional).
\end{nota}

% ============================================================
\part{Los derechos reales en el sistema jurídico}
% ============================================================

% ------------------------------------------------------------
\chapter[Derechos personales y derechos reales]{Derechos personales y derechos reales}\label{cap:dp-dr}
% ------------------------------------------------------------

El punto de partida es una distinción que organiza todo el derecho patrimonial: la que separa los derechos personales de los derechos reales. Comprenderla con precisión permite entender por qué los derechos reales reciben una regulación propia, de orden público, y por qué el resto de este manual dedica tanto espacio a sus principios comunes.

\section{La teoría clásica o del vínculo}

El Código Civil y Comercial vigente adopta la distinción entre \textbf{derechos reales} y \textbf{derechos personales}, conocida como \textbf{teoría clásica o del vínculo}, que es la teoría aplicable en la mayoría de los sistemas occidentales de derecho. Los derechos personales ---también llamados \textbf{derechos obligacionales}--- y los derechos reales ---también llamados \textbf{derechos de las cosas}--- presentan una profunda distinción y reciben una regulación específica cada uno.

La teoría los distingue según el vínculo que establecen:

\begin{itemize}
    \item \textbf{derechos personales:} vínculo entre sujetos (entre personas);
    \item \textbf{derechos reales:} vínculo directo entre una persona y una cosa.
\end{itemize}

Ambas categorías integran el \textbf{patrimonio} de una persona. Según el apunte, el Código adopta esta teoría porque organiza el sistema patrimonial occidental moderno y permite dar respuestas claras acerca de quién puede exigir qué y sobre qué objeto.

\section{El derecho personal u obligacional}

\begin{definicion}{Obligación (art.~724 CCyC)}
«La obligación es una relación jurídica en virtud de la cual el acreedor tiene el derecho de exigir del deudor una prestación destinada a satisfacer un interés lícito y ante el incumplimiento obtener forzadamente la satisfacción de dicho interés.»
\end{definicion}

\subsection{Elementos}

El derecho personal tiene tres elementos esenciales:

\begin{itemize}
    \item \textbf{sujeto activo o acreedor:} quien tiene el derecho de exigir la prestación;
    \item \textbf{sujeto pasivo o deudor:} quien está obligado a cumplirla;
    \item \textbf{prestación:} la conducta debida, que puede ser de dar, de hacer o de no hacer (omisión).
\end{itemize}

El vínculo se representa como una \textbf{ola o vínculo circular} que encierra dentro de sí al acreedor, al deudor y a la prestación. Es un vínculo \textbf{cerrado}, lo que explica su carácter relativo.

\subsection{Caracteres}

\begin{itemize}
    \item \textbf{Relativos:} el vínculo existe solo entre las partes; el acreedor solo puede exigir el cumplimiento a su propio deudor.
    \item \textbf{Autonomía de la voluntad:} los particulares pueden crear, modificar y extinguir obligaciones; la ley es supletoria.
    \item \textbf{Objeto:} una conducta humana (dar, hacer, no hacer).
    \item \textbf{Fuente:} nacen de hechos o actos jurídicos, esencialmente el contrato, y también de la obligación de no dañar (responsabilidad civil).
    \item \textbf{Extinción:} se extinguen por cumplimiento, por prescripción liberatoria u otras formas de extinción de las obligaciones. Son siempre temporales.
    \item \textbf{No se adquieren} por el mero transcurso del tiempo.
\end{itemize}

\begin{ejemplo}{obligación alimentaria, crédito y préstamo}
En la \textbf{obligación alimentaria}, el hijo (acreedor alimentario) solo puede exigir la prestación a quienes la ley identifica como obligados (padres, abuelos, etc.), y no a un tercero. Del mismo modo, el acreedor de un \textbf{crédito} (por ejemplo, el banco) solo puede reclamar el pago a su deudor.

Si se presta una suma de dinero a un amigo, nace una obligación: el prestamista es el acreedor, el amigo es el deudor y la prestación es dar, es decir, devolver el dinero. Solo puede reclamarse al amigo, y no a un tercero (su pareja o su vecino): por eso el derecho personal es \textbf{relativo}. Si el deudor no paga, el acreedor puede acudir a la justicia para obligarlo; pero cuando cobra, el vínculo se termina. Y si el acreedor no reclama durante años, el derecho se extingue por prescripción: el tiempo juega en su contra.
\end{ejemplo}

\section{El derecho real}

\begin{definicion}{Derecho real (art.~1882 CCyC)}
«El derecho real es el poder jurídico de estructura legal que se ejerce directamente sobre su objeto en forma autónoma y que le atribuye a su titular las facultades de persecución y preferencia y las demás previstas en este código.»
\end{definicion}

El derecho real implica un vínculo \textbf{directo e inmediato} entre el \textbf{titular} y la \textbf{cosa}; estos son sus dos elementos. Ese vínculo es tan intenso que permite al titular:

\begin{itemize}
    \item \textbf{disponer de la cosa u obtener algún beneficio de ella} en forma directa, sin intermediarios;
    \item \textbf{perseguir la cosa} en poder de quien se encuentre;
    \item \textbf{ser preferido} a cualquier otro que pudiera tener sobre ella algún derecho real o personal de fecha posterior.
\end{itemize}

\subsection{Caracteres}

\begin{itemize}
    \item \textbf{Absolutos o \textit{erga omnes}:} son oponibles a toda la comunidad, que debe respetar el derecho del titular. El sujeto pasivo es, por tanto, toda la comunidad, con una obligación general de no perturbar al titular.
    \item \textbf{Estructura legal y orden público:} son creados y regulados exclusivamente por la ley (\textit{numerus clausus}, art.~1884). Los particulares no pueden crear ni modificar derechos reales, y todo acto contrario es nulo.
    \item \textbf{Objeto:} la cosa, es decir, un bien material. Excepcionalmente, bienes inmateriales taxativamente previstos.
    \item \textbf{Nacimiento:} título suficiente y modo (tradición), ley o prescripción adquisitiva.
    \item \textbf{Duración:} pueden ser perpetuos (como el dominio) o temporales.
    \item \textbf{Facultades:} persecución (\textit{ius persequendi}) y preferencia.
\end{itemize}

\begin{ejemplo}{la compra de un departamento}
Quien compra un departamento, firma la escritura y recibe la entrega (tradición) adquiere el dominio, que es un derecho real. Su relación no es con una persona específica, sino directamente con la cosa: es dueño y no necesita permiso de nadie para usarla, disfrutarla o venderla. Todos deben respetar esa propiedad, no solo el vendedor; por eso el derecho es \textit{erga omnes}. Si la cosa fuera robada u ocupada ilegalmente y transmitida a un tercero, el titular puede perseguirla: el derecho real «sigue» a la cosa sin importar en qué manos se encuentre (\textit{ius persequendi}).
\end{ejemplo}

\section{Cuadro comparativo}

\begin{table}[H]
\centering
\small
\begin{tabularx}{\textwidth}{>{\raggedright\arraybackslash}p{0.19\textwidth}>{\raggedright\arraybackslash}X>{\raggedright\arraybackslash}X}
\toprule
\textbf{Aspecto} & \textbf{Derecho personal} & \textbf{Derecho real} \\
\midrule
Vínculo & Entre sujetos (acreedor y deudor) & Entre titular y cosa \\
Oponibilidad & Relativo: solo entre las partes & Absoluto: \textit{erga omnes} \\
Sujeto pasivo & El deudor determinado & Toda la comunidad \\
Objeto & Conducta humana (dar, hacer, no hacer) & La cosa (bien material) \\
Creación y regulación & Autonomía de la voluntad; la ley es supletoria & Exclusivamente la ley; normas de orden público \\
Nacimiento & Hechos o actos jurídicos: contrato, daño & Título y modo (tradición), ley, prescripción adquisitiva \\
Duración & Siempre temporal & Perpetuo o temporal \\
Prescripción & Liberatoria: extingue el derecho & Adquisitiva: permite adquirirlo \\
Facultades de persecución y preferencia & No hay persecución & Sí \\
Ejemplo & Contrato de compraventa (obligación de entregar) & Dominio sobre un inmueble \\
\bottomrule
\end{tabularx}
\caption{Derechos personales y derechos reales}
\end{table}

\begin{conceptoclave}{la diferencia esencial}
El derecho personal es un vínculo entre personas: relativo y temporal. El derecho real es un vínculo entre una persona y una cosa: absoluto, \textit{erga omnes}, regulado solo por la ley y defendido mediante la persecución de la cosa donde quiera que esté.
\end{conceptoclave}

\cornell{derechos personales y derechos reales}{
\cq{¿Qué teoría adopta el CCyC y qué criterio usa?}{La teoría clásica o del vínculo: los derechos personales vinculan sujetos; los reales vinculan a una persona con una cosa. Ambos integran el patrimonio.}
\cq{Art.~724: ¿qué es una obligación y cuáles son sus elementos?}{Relación jurídica por la que el acreedor puede exigir al deudor una prestación que satisface un interés lícito y, ante el incumplimiento, obtener forzadamente su satisfacción. Elementos: acreedor, deudor y prestación (dar, hacer, no hacer).}
\cq{¿Por qué el derecho personal es relativo?}{Porque el vínculo es cerrado: existe solo entre las partes y el acreedor solo puede exigir a su deudor.}
\cq{Art.~1882: ¿qué es un derecho real y qué facultades confiere?}{Poder jurídico de estructura legal ejercido directamente sobre su objeto, en forma autónoma. Confiere persecución y preferencia, además de las previstas en el Código.}
\cq{¿Por qué el derecho real es absoluto?}{Porque es oponible \textit{erga omnes}: toda la comunidad debe respetar al titular.}
\cq{Compare la prescripción en ambas esferas.}{En los personales es liberatoria (extingue el derecho); en los reales, adquisitiva (permite adquirirlo). Los personales son siempre temporales; los reales pueden ser perpetuos.}
}{El Código distingue los derechos personales (vínculo entre acreedor y deudor, relativo, temporal, regido por la autonomía de la voluntad) de los reales (vínculo directo con la cosa, absoluto, de estructura legal y orden público, con persecución y preferencia). Ambos integran el patrimonio.}

% ------------------------------------------------------------
\chapter[Los derechos reales en el Código Civil y Comercial]{Los derechos reales en el Código Civil y Comercial}\label{cap:ccyc}
% ------------------------------------------------------------

Una vez delimitado el concepto, corresponde ubicar a los derechos reales dentro del Código y precisar sobre qué objeto recaen. Para ello son centrales los artículos 15, 16 y 17.

\section{Estructura metodológica del Código}

El CCyC tiene una \textbf{metodología novedosa} respecto del Código de Vélez Sársfield. Su estructura es la siguiente.

\begin{table}[H]
\centering
\small
\begin{tabularx}{\textwidth}{>{\raggedright\arraybackslash}p{0.22\textwidth}>{\raggedright\arraybackslash}p{0.27\textwidth}>{\raggedright\arraybackslash}X}
\toprule
\textbf{Ubicación} & \textbf{Denominación} & \textbf{Contenido relevante} \\
\midrule
Título Preliminar & Principios rectores & Buena fe, abuso del derecho, abuso de posición dominante, preeminencia del orden público, tratados de derechos humanos como fuente y guía de interpretación; incluye el Capítulo 4, sobre bienes y cosas \\
Libro Primero & Parte General & Persona, capacidad \\
Libro Segundo & Relaciones de Familia & Derecho de familia \\
Libro Tercero & Derechos Personales & Obligaciones, contratos \\
Libro Cuarto & Derechos Reales & Teoría general y derechos reales en particular (arts.~1882 a 2276) \\
Libro Quinto & Transmisión de derechos por causa de muerte & Sucesiones \\
Libro Sexto & Disposiciones comunes & A derechos reales y personales \\
\bottomrule
\end{tabularx}
\caption{Estructura del Código Civil y Comercial}
\end{table}
% TODO: verificar consistencia entre fuentes originales (una versión denomina «Título Sexto» y otra «Libro VI» a las disposiciones comunes).

\begin{importante}
El estudio de los derechos reales \textbf{excede el Libro Cuarto}: deben integrarse el Título Preliminar, las disposiciones comunes del Libro Sexto y el resto del ordenamiento.
\end{importante}

\section{Patrimonio, bienes y cosas: artículos 15, 16 y 17}

\begin{regla}{art.~15 CCyC, titularidad}
«Las personas son titulares de los derechos individuales sobre los bienes que integran su patrimonio conforme con lo que se establece en este Código.»
\end{regla}
% TODO: verificar referencia presente en las fuentes originales (una versión transcribe el art.~15 sin la expresión final «conforme con lo que se establece en este Código»).

El patrimonio está integrado por bienes, y estos comprenden tanto los derechos personales como los derechos reales y las cosas materiales. Ambas categorías conviven en el patrimonio. El artículo 15 expresa, además, el sistema de propiedad privada (véase el capítulo~\ref{cap:legalidad}).

\begin{regla}{art.~16 CCyC, bienes y cosas}
«Los derechos pueden recaer sobre bienes susceptibles de valor económico. Cuando son materiales se llaman cosas. Las disposiciones sobre las cosas se aplican a la energía y a las fuerzas naturales susceptibles de ser puestas al servicio del hombre.»
\end{regla}

Esta norma es esencial para determinar el \textbf{objeto} de los derechos reales. Los bienes son los derechos susceptibles de valor económico; cuando son materiales (tangibles) se llaman cosas. Toda cosa es un bien, pero no todo bien es una cosa: por ejemplo, un crédito es un bien pero no una cosa. La energía y las fuerzas naturales se asimilan a las cosas.

\begin{regla}{art.~17 CCyC, derechos sobre el cuerpo humano}
Los derechos sobre el cuerpo humano o sus partes no tienen valor comercial, sino afectivo, científico, terapéutico, humanitario o social. Su disposición está sujeta a lo dispuesto en las leyes especiales (por ejemplo, órganos y genes).
\end{regla}

Consecuencia directa: el cuerpo humano y sus partes \textbf{no son objeto de derechos reales}, porque el artículo 16 exige que el bien tenga valor económico para ser considerado cosa.

\section{Los principios comunes: artículos 1882 a 1891}

Los artículos 1882 a 1891 regulan los principios comunes a todos los derechos reales. Esta parte general es la base teórica de toda la materia: muchos de sus principios no se repiten al estudiar cada derecho real en particular. Por eso, cada vez que se aborda un derecho real específico conviene «exponer en acto» lo aprendido sobre las disposiciones comunes, aplicándolas al derecho de que se trate (así se hace, por ejemplo, con el dominio en el capítulo~\ref{cap:dominio}).

\cornell{los derechos reales en el Código}{
\cq{¿Dónde se regulan los derechos reales y basta con el Libro Cuarto?}{En el Libro Cuarto (arts.~1882 a 2276), pero su estudio lo excede: se integran el Título Preliminar, las disposiciones comunes y el resto del ordenamiento.}
\cq{¿Qué contiene el Título Preliminar?}{Los principios rectores de todo el Código: buena fe, abuso del derecho, abuso de posición dominante, orden público y tratados de derechos humanos. Incluye el capítulo sobre bienes y cosas.}
\cq{¿Qué integra el patrimonio (art.~15)?}{Los bienes: derechos personales, derechos reales y cosas materiales. Las personas son titulares de derechos individuales sobre ellos.}
\cq{¿Cuándo un bien es una cosa (art.~16)? ¿Todo bien es cosa?}{Cuando es material. Toda cosa es un bien, pero no todo bien es cosa (el crédito es un bien, no una cosa). La energía y las fuerzas naturales se asimilan a las cosas.}
\cq{¿Por qué el cuerpo humano no es objeto de derechos reales (art.~17)?}{Porque sus partes tienen valor afectivo, científico, terapéutico, humanitario o social, pero no valor comercial ni económico, y se rigen por leyes especiales.}
}{Los derechos reales se regulan en el Libro Cuarto, cuyos arts.~1882 a 1891 contienen los principios comunes. Los arts.~15, 16 y 17 fijan la titularidad sobre los bienes del patrimonio, el concepto de cosa como bien material con valor económico y la exclusión del cuerpo humano como objeto de derechos reales.}

% ------------------------------------------------------------
\chapter[Legalidad, orden público y fundamento constitucional]{Legalidad, orden público y fundamento constitucional}\label{cap:legalidad}
% ------------------------------------------------------------

Los caracteres de los derechos reales no son una enumeración arbitraria: responden a una idea común. Como el derecho real debe ser respetado por toda la comunidad, solo la ley puede crearlo y regularlo, y su ejercicio queda sometido a los intereses de la sociedad. Este capítulo desarrolla esa idea y los principios que la expresan.

\section{\textit{Numerus clausus} y ley federal}

\begin{regla}{art.~1884 CCyC, estructura de los derechos reales}
«La regulación de los derechos reales en cuanto a sus elementos, contenido, adquisición, constitución, modificación, transmisión, duración y extinción está establecida solo por la ley. Es nula la configuración de un derecho real no previsto en la ley o la modificación de su estructura.»
\end{regla}

Este principio de \textit{numerus clausus} (número cerrado) significa que los particulares \textbf{no pueden crear nuevos derechos reales}: solo la ley los crea y regula. Los derechos reales son \textbf{catorce y solo catorce}, y hoy se encuentran todos codificados. La sanción por apartarse de la regulación legal es la \textbf{nulidad} del acto. El derecho real es, además, un señorío de la voluntad que se ejerce en forma autónoma e independiente de otra voluntad.

La competencia para crearlos es exclusivamente federal: se trata de \textbf{ley federal}, porque en materia de derechos reales solo la Nación es competente para legislar.

\begin{ejemplo}{la lista cerrada}
Los derechos reales son como una lista cerrada escrita por el Congreso, o como los planetas del sistema solar: solo existen los que la ley nombra. Si dos personas firman un contrato en el que dicen «creamos un nuevo derecho real que se llama semi-dominio» (o «semi-usufructo», o «super-dominio compartido especial»), ese acuerdo es nulo en cuanto al derecho real: no existe. En cambio, en los derechos personales las partes pueden crear contratos atípicos: allí hay libertad. En los derechos reales, no, porque afectan a toda la sociedad.
\end{ejemplo}

\section{Normas de orden público}

Los particulares deben ceñirse estrictamente a la regulación legal de cada derecho real, porque las normas que los rigen son de \textbf{orden público}: expresan la voluntad de la Nación y no pueden ser modificadas por los particulares.

\begin{regla}{art.~12 CCyC}
Las convenciones particulares no pueden dejar sin efecto las leyes en cuya observancia está interesado el orden público.
\end{regla}

La razón de fondo es que el derecho real de cualquier país responde a un \textbf{esquema de sociedad único para toda la Nación}. Tiene que ver con tres cuestiones: el modo en que una nación organiza su sistema de propiedad y su sistema de defensa de las relaciones de poder respecto de las cosas; el modo en que define el alcance de las facultades que está dispuesta a reconocer a cada titular; y las limitaciones que está dispuesta a aceptar en función de los intereses de la comunidad.

\section{Constitucionalización e internacionalización del derecho privado}

El CCyC evidencia un cambio de paradigma respecto de la legislación anterior, esencialmente porque incorpora los tratados internacionales de derechos humanos como ley vigente. Por eso puede hablarse de \textbf{constitucionalización} e \textbf{internacionalización del derecho privado}. Los artículos 1 y 2 son fundamentos del Código y, a la vez, una regla de interpretación: todo el sistema se mira bajo esa perspectiva.

En ese marco conviven en tensión varias normas.

\begin{table}[H]
\centering
\small
\begin{tabularx}{\textwidth}{>{\raggedright\arraybackslash}p{0.15\textwidth}>{\raggedright\arraybackslash}X>{\raggedright\arraybackslash}p{0.27\textwidth}}
\toprule
\textbf{Artículo} & \textbf{Contenido} & \textbf{Efecto} \\
\midrule
Art.~15 & Las personas son titulares de los derechos individuales sobre los bienes que integran su patrimonio. & Consagra el sistema de propiedad privada. \\
Art.~14 & La ley no ampara el ejercicio abusivo de los derechos individuales cuando pueda afectar al ambiente y a los derechos de incidencia colectiva. & Limita los derechos individuales frente al interés colectivo. \\
Art.~240 & Establece un límite a los derechos individuales cuando resulte necesario compatibilizarlos con los derechos de incidencia colectiva. & Preeminencia del interés colectivo sobre el individual. \\
\bottomrule
\end{tabularx}
\caption{Propiedad privada y sus límites}
\end{table}

El derecho al ambiente y el derecho a respirar son ejemplos de derechos de incidencia colectiva que pueden entrar en tensión con el ejercicio de los derechos reales. En el sistema jurídico existe una clara \textbf{preeminencia del interés colectivo por sobre el interés individual}; por eso, en ocasiones (por ejemplo, en una situación de contingencia), resulta necesario modelar el alcance del ejercicio de un derecho real en función de los intereses de la comunidad.

El derecho real se vincula así con grandes temas del derecho: el ejercicio del derecho de propiedad, los límites al dominio, el reconocimiento de la propiedad como derecho humano y la función social. Son temas susceptibles de enfrentar ideológicamente a distintos sectores de la doctrina nacional e internacional (se desarrollan en el capítulo~\ref{cap:limites}).

\section{\textit{Nemo plus iuris} y convalidación}

\begin{definicion}{\textit{nemo plus iuris} (art.~399 CCyC)}
Nadie puede transmitir a otro un derecho mejor o más extenso que el que tiene, y nadie puede adquirir sobre un objeto un derecho mejor o más extenso que el que tenía aquel de quien lo adquiere.
\end{definicion}

Es un principio de origen romano que limita la transmisión de derechos y es el límite infranqueable de la transmisión. Se aplica a los derechos reales y a los personales y, si bien admite excepciones previstas en la ley, constituye la regla general. Quien transmite un derecho que no tiene realiza un acto que nace nulo.

\begin{regla}{art.~1885 CCyC, convalidación}
«Si alguien constituye o transmite un derecho real que no tiene, y lo adquiere posteriormente, la constitución o transmisión queda convalidada.»
\end{regla}
% TODO: verificar consistencia entre fuentes originales (el art.~1885 se transcribe con redacciones distintas y, en el desarrollo del dominio, se lo describe con un alcance diferente: ver capítulo sobre adquisición del dominio).

La convalidación es la \textbf{excepción} al \textit{nemo plus iuris}: si después de la transmisión nula el transmitente adquiere el derecho, el acto queda automáticamente saneado, como si hubiera sido válido desde el principio. El apunte señala que existe para proteger a quien adquirió de buena fe y dar seguridad jurídica al sistema.

\begin{ejemplo}{la venta de la cosa ajena}
Quien tiene una cosa en su poder pero no es propietario, y la transmite, no hace adquirir la propiedad a quien la recibe: no puede transmitir un derecho mejor que el que tenía. El acto nace nulo. Pero si después el transmitente adquiere la propiedad (por ejemplo, porque el dueño se la vende o la hereda), el acto queda convalidado retroactivamente.
\end{ejemplo}

\section{Persecución y preferencia}

Son las dos facultades características del derecho real.

\begin{definicion}{\textit{ius persequendi} (persecución)}
El titular puede \textbf{perseguir la cosa} en poder de quien se encuentre, sin importar si cambió de titular: el derecho la sigue. Se la representa como una flecha que persigue la cosa sin importar en qué manos esté.
\end{definicion}

\begin{definicion}{\textit{ius preferendi} (preferencia)}
El titular es \textbf{preferido} a cualquier otro que pudiera tener sobre la cosa algún derecho real o personal de fecha posterior.
\end{definicion}

Es una diferencia clave con los derechos personales, donde no hay persecución. Ambas facultades derivan de la relación directa entre titular y cosa y de su carácter absoluto.

\cornell{legalidad y orden público}{
\cq{¿Qué consagra el art.~1884 y cuál es la sanción?}{El \textit{numerus clausus}: solo la ley crea y regula los derechos reales (elementos, contenido, adquisición, constitución, modificación, transmisión, duración y extinción). La sanción es la nulidad.}
\cq{¿Por qué las normas de derechos reales son de orden público y de ley federal?}{Porque expresan un esquema de sociedad único para toda la Nación (sistema de propiedad, defensa de las relaciones de poder y límites aceptados por la comunidad); solo la Nación legisla. Art.~12: las convenciones no pueden dejar sin efecto leyes de orden público.}
\cq{¿Qué articulan los arts.~15, 14 y 240?}{La propiedad privada (art.~15) y sus límites frente al ambiente y a los derechos de incidencia colectiva (arts.~14 y 240), con preeminencia del interés colectivo.}
\cq{Relacione \textit{nemo plus iuris} y convalidación.}{Regla general: nadie transmite un derecho mejor que el que tiene (acto nulo). Excepción (art.~1885): si el transmitente adquiere luego el derecho, el acto queda convalidado.}
\cq{¿Qué son la persecución y la preferencia?}{Persecución: seguir la cosa en poder de quien se encuentre. Preferencia: prevalecer sobre derechos reales o personales de fecha posterior. Son propias del derecho real.}
}{Solo la ley federal crea y regula los derechos reales (\textit{numerus clausus}), bajo normas de orden público y con sanción de nulidad. Ese carácter responde a un esquema de sociedad único y a la preeminencia del interés colectivo. El \textit{nemo plus iuris} limita la transmisión; la convalidación la sana si el transmitente adquiere después el derecho. La persecución y la preferencia son las facultades características del derecho real.}

% ------------------------------------------------------------
\chapter[El catálogo y la clasificación de los derechos reales]{El catálogo y la clasificación de los derechos reales}\label{cap:catalogo}
% ------------------------------------------------------------

Si solo la ley crea derechos reales, cabe conocer cuáles son. Este capítulo presenta la lista cerrada del artículo 1887 y las clasificaciones del artículo 1888, que serán la guía para ordenar los derechos que se estudian en las partes siguientes.

\section{Los catorce derechos reales (art.~1887)}

El artículo 1887 enumera taxativamente los derechos reales. Solo existen estos y no hay más.

\begin{table}[H]
\centering
\small
\begin{tabularx}{\textwidth}{>{\centering\arraybackslash}p{0.07\textwidth}>{\raggedright\arraybackslash}X>{\raggedright\arraybackslash}X}
\toprule
\textbf{N.º} & \textbf{Derecho real} & \textbf{Recae sobre cosa} \\
\midrule
1 & Dominio & Propia \\
2 & Condominio & Propia (varios titulares) \\
3 & Propiedad horizontal & Propia \\
4 & Conjuntos inmobiliarios & Propia \\
5 & Tiempo compartido & Propia \\
6 & Cementerio privado & Propia \\
7 & Superficie & Propia (con propiedad superficiaria) o ajena (sin propiedad superficiaria) \\
8 & Usufructo & Ajena \\
9 & Uso & Ajena \\
10 & Habitación & Ajena \\
11 & Servidumbre & Ajena \\
12 & Hipoteca & Ajena (garantía) \\
13 & Anticresis & Ajena (garantía) \\
14 & Prenda & Ajena (garantía) \\
\bottomrule
\end{tabularx}
\caption{Los derechos reales del art.~1887}
\end{table}
% TODO: verificar consistencia entre fuentes originales (una versión numera la superficie dos veces, según exista o no propiedad superficiaria, lo que daría 15 filas para 14 derechos; aquí se la cuenta una sola vez).

\begin{nota}
Regla mnemotécnica propuesta en los apuntes para recordar la lista: \textit{Do-Con-PH-CI-TC-CP-Su-Us-U-H-S-Hi-An-P}.
\end{nota}

La enumeración se ordena en este manual así: el dominio y sus límites (Parte~IV); el condominio y la medianería (Parte~V); la propiedad horizontal (Parte~VI); y los conjuntos inmobiliarios, el tiempo compartido y el cementerio privado (Parte~VII). El tratamiento de los restantes derechos reales no forma parte de las fuentes de este manual.

\section{Clasificación (art.~1888)}

\subsection{Sobre cosa propia y sobre cosa ajena}

Los derechos sobre \textbf{cosa propia} recaen sobre cosas del propio titular: dominio, condominio, propiedad horizontal, conjuntos inmobiliarios, tiempo compartido, cementerio privado y superficie (con propiedad superficiaria). Los derechos sobre \textbf{cosa ajena} recaen sobre cosas de otra persona y constituyen una \textbf{carga o gravamen} real para el propietario: usufructo, uso, habitación, servidumbre, superficie (sin propiedad superficiaria), hipoteca, anticresis y prenda.

\begin{ejemplo}{dominio e hipoteca}
Una persona es titular de dominio (derecho real sobre cosa propia) y constituye una hipoteca a favor de un banco para garantizar un crédito. Ha creado un nuevo derecho real: el banco tiene ahora un derecho real sobre esa cosa (hipoteca: derecho real sobre cosa ajena, de garantía y accesorio). Coexisten el dominio del titular y la hipoteca del banco.
\end{ejemplo}

\subsection{Principales y accesorios}

\begin{itemize}
    \item \textbf{Principales:} los derechos sobre cosa propia y los derechos de disfrute sobre cosa ajena (usufructo, servidumbre, uso, habitación, superficie). Son independientes: existen por sí solos.
    \item \textbf{Accesorios:} los derechos reales de garantía (hipoteca, prenda, anticresis). Garantizan una obligación principal, dependen de un crédito al que acceden y lo siguen.
\end{itemize}

\subsection{Registrables y no registrables}

Según su adquisición, transmisión o extinción requieran o no inscripción en un registro público. Son \textbf{registrables} los inmuebles (Registro de la Propiedad Inmueble) y los automotores (Registro de Automotores); son \textbf{no registrables} las cosas muebles en general.

\subsection{Ejercicio por posesión o sin ella}

Algunos derechos reales se ejercen mediante la \textbf{posesión} (poder de hecho sobre la cosa) y otros por otros modos; las fuentes mencionan para estos últimos la \textbf{cuasiposesión}. Este tema se desarrolla en la Parte~III.
% TODO: verificar consistencia entre fuentes originales (una versión contrapone «posesión» y «cuasiposesión»; otra, «con posesión» y «sin posesión»).

\section{Casos de aplicación}

\begin{ejercicio}{creación de un derecho real inexistente}
\textbf{Enunciado.} Pablo firma un contrato con Marta por el cual se crea un «derecho real de semi-usufructo» que no existe en ninguna ley. ¿Qué sucede?

\textbf{Resolución.} El acto es nulo (art.~1884 CCyC, \textit{numerus clausus}). Solo la ley puede crear derechos reales; los particulares no pueden crear uno que no esté en el artículo 1887.
\end{ejercicio}

\begin{ejercicio}{venta de un inmueble ajeno}
\textbf{Enunciado.} Juan vende a María un inmueble que pertenece a Carlos. Después, Juan hereda ese inmueble de Carlos. ¿Qué pasa con la venta?

\textbf{Resolución.} Originalmente el acto era nulo (\textit{nemo plus iuris}: Juan no podía transmitir lo que no tenía). Al heredar el inmueble, la transmisión queda convalidada (art.~1885 CCyC) y María es ahora legítima propietaria.
\end{ejercicio}

\begin{ejercicio}{hipoteca y venta del inmueble hipotecado}
\textbf{Enunciado.} El banco tiene una hipoteca sobre la casa de Roberto, quien la vende a Sofía. ¿Puede el banco ejecutar la hipoteca frente a Sofía?

\textbf{Resolución.} Sí. La hipoteca es un derecho real que otorga la persecución: el banco puede perseguir la cosa en manos de cualquier persona que la tenga, incluida Sofía, que adquirió la casa con la carga hipotecaria.
\end{ejercicio}

El apéndice~\ref{ap:ejercicios} reúne ejercicios adicionales de evaluación sobre toda la teoría general.

\cornell{catálogo y clasificación}{
\cq{¿Cuántos derechos reales hay y cuál es su fuente?}{Catorce (art.~1887), todos creados y regulados por la ley: dominio, condominio, propiedad horizontal, conjuntos inmobiliarios, tiempo compartido, cementerio privado, superficie, usufructo, uso, habitación, servidumbre, hipoteca, anticresis y prenda.}
\cq{Distinga cosa propia y cosa ajena (art.~1888).}{Propia: el derecho recae sobre cosa del propio titular. Ajena: recae sobre cosa de otro y es carga o gravamen para su propietario. La superficie es propia si hay propiedad superficiaria y ajena si no la hay.}
\cq{¿Qué diferencia a los principales de los accesorios?}{Los principales son independientes (cosa propia y disfrute de cosa ajena). Los accesorios, de garantía, dependen de un crédito y lo siguen.}
\cq{¿Qué son los derechos registrables?}{Los que recaen sobre cosas cuya adquisición, transmisión o extinción requiere inscripción en un registro público (inmuebles y automotores).}
\cq{Caso: el banco hipotecario y el comprador del inmueble.}{El banco puede perseguir la cosa en manos del comprador, porque la hipoteca es un derecho real con persecución.}
}{El art.~1887 fija una lista cerrada de catorce derechos reales y el art.~1888 los clasifica según recaigan sobre cosa propia o ajena, sean principales o accesorios, registrables o no registrables, y se ejerzan o no por posesión.}

% ============================================================
\part{Adquisición, publicidad, transmisión y extinción}
% ============================================================

% ------------------------------------------------------------
\chapter[Adquisición de los derechos reales: título y modo]{Adquisición de los derechos reales: título y modo}\label{cap:adquisicion}
% ------------------------------------------------------------

Conocidos los principios y el catálogo, corresponde preguntar cómo se llega a ser titular de un derecho real. El sistema reconoce distintas formas de adquisición; la más importante para los actos entre vivos exige la concurrencia de dos elementos: el título suficiente y el modo suficiente.

\section{Adquisición originaria y derivada}

\begin{table}[H]
\centering
\small
\begin{tabularx}{\textwidth}{>{\raggedright\arraybackslash}p{0.19\textwidth}>{\raggedright\arraybackslash}X>{\raggedright\arraybackslash}p{0.28\textwidth}}
\toprule
\textbf{Tipo} & \textbf{Concepto} & \textbf{Ejemplos} \\
\midrule
Originaria & El sistema reconoce la posibilidad de adquirir el derecho real en forma libre y absoluta, sin que derive de un vínculo jurídico anterior. & Cosas sin dueño (\textit{res nullius}); prescripción adquisitiva (usucapión); adquisición legal. \\
Derivada, entre vivos & El derecho real es transmitido por un enajenante a un adquirente: uno transmite y el otro adquiere. Puede ser a título oneroso o gratuito. & Compraventa, permuta, donación, dación en pago, cesión. \\
Derivada, por causa de muerte & La transmisión opera por el fallecimiento del titular: el heredero ocupa el lugar del causante y adquiere todo lo que este tenía. & Sucesión hereditaria. \\
\bottomrule
\end{tabularx}
\caption{Formas de adquisición de los derechos reales}
\end{table}

\begin{regla}{art.~1896 CCyC}
El juez no puede constituir derechos reales ni imponer su constitución.
\end{regla}

Es una expresión del principio de legalidad: el derecho real se adquiere por los modos previstos en la ley, no por decisión judicial. En toda adquisición derivada rige además el \textit{nemo plus iuris} (art.~399; véase el capítulo~\ref{cap:legalidad}), que impide adquirir un derecho mejor o más extenso que el que tenía el transmitente.

\section{La teoría del título y modo}

\begin{regla}{art.~1892 CCyC}
El derecho real se adquiere con la concurrencia del título suficiente más el modo suficiente.
\end{regla}

La adquisición derivada de derechos reales por actos entre vivos requiere necesariamente ambos elementos. Uno sin el otro no es suficiente.

\subsection{El título suficiente}

\begin{definicion}{título suficiente}
Es el acto jurídico revestido de las condiciones de fondo y de forma establecidas por la ley, que tiene por finalidad transmitir o constituir un derecho real. El acto jurídico causal de la transmisión constituye el título suficiente; su causa puede ser una compraventa, una cesión, una permuta, una dación en pago o una donación.
\end{definicion}

El título no debe identificarse en su sentido instrumental (el papel), sino en su \textbf{sentido abstracto o intangible}. Para ilustrarlo, el apunte recurre a una analogía: así como no es posible decir «te doy mi amor, dejámelo en el buzón», ante la falta de título suficiente no se resuelve afirmando que el título fue «dejado en el buzón». Entregar un instrumento que no plasma un acto jurídico causal equivale a no tener título suficiente.

Los artículos 259 y 260 aportan los conceptos de base:

\begin{itemize}
    \item \textbf{Art.~259, acto jurídico:} el acto voluntario lícito que tiene por fin inmediato la adquisición, modificación o extinción de relaciones o situaciones jurídicas.
    \item \textbf{Art.~260, acto voluntario:} el ejecutado con discernimiento, intención y libertad, que se manifiesta por un hecho exterior.
    \item \textbf{Art.~1017:} en materia de inmuebles, el acto jurídico causal debe plasmarse en \textbf{escritura pública}.
\end{itemize}

\subsection{El modo suficiente: la tradición}

El modo es el elemento material propio de los derechos reales que se ejercen por la posesión.

\begin{definicion}{tradición}
Es la entrega material de la cosa: una persona entrega la cosa a otra, que la recibe. El poder de hecho que tenía el sujeto se transmite, con la entrega, a quien la adquiere.
\end{definicion}

Quien entrega es el \textbf{tradente} y quien recibe es el \textbf{accipiens}. La tradición es \textbf{constitutiva} en todos los casos en que se adquieren derechos reales que se ejercen por la posesión: sin tradición no se adquirió el derecho real.

\begin{regla}{art.~1924 CCyC}
La tradición no puede suplirse por fórmulas escritas.
\end{regla}

Cuando una escritura pública dice que «en este acto las partes se entregan la tradición», se pretende suplantar un acto material por una fórmula escrita en el papel, lo cual no puede hacerse: la posesión no se transmite con palabras escritas. Las únicas excepciones son la \textit{traditio brevi manu} y el constituto posesorio (capítulo~\ref{cap:brevi}).

\section{Requisitos de la tradición traslativa de dominio}

No toda entrega material implica transferencia de dominio. Para que la tradición sea traslativa deben concurrir los siguientes requisitos.

\begin{table}[H]
\centering
\small
\begin{tabularx}{\textwidth}{>{\raggedright\arraybackslash}p{0.25\textwidth}>{\raggedright\arraybackslash}X>{\raggedright\arraybackslash}p{0.21\textwidth}}
\toprule
\textbf{Requisito} & \textbf{Descripción} & \textbf{Fundamento} \\
\midrule
Legitimación del transmitente & Debe ser realizada por el propietario: solo el dueño puede transmitir el dominio. & Art.~399, \textit{nemo plus iuris} \\
Acto causal (título suficiente) & Debe existir una causa válida: compraventa, donación, dación en pago, permuta, etc. & Art.~1892 \\
Capacidad de ambas partes & Tanto quien transmite como quien adquiere deben ser capaces. & Art.~1892 \\
Escritura pública (inmuebles) & El acto causal debe plasmarse en escritura pública. & Art.~1017 \\
Tradición material efectiva & La entrega física de la cosa, salvo los dos casos de excepción. & Art.~1924 \\
\bottomrule
\end{tabularx}
\caption{Requisitos de la tradición traslativa de dominio}
\end{table}

Si faltan los requisitos de legitimación, título y capacidad, lo que se transmite es meramente la relación de poder que el tradente tenía con la cosa, y no el derecho (véase el capítulo~\ref{cap:dinamica}).

\section{La inscripción registral}

El efecto de la inscripción registral depende de la clase de bien.

\begin{table}[H]
\centering
\small
\begin{tabularx}{\textwidth}{>{\raggedright\arraybackslash}p{0.17\textwidth}>{\raggedright\arraybackslash}p{0.19\textwidth}>{\raggedright\arraybackslash}X}
\toprule
\textbf{Bien} & \textbf{Efecto} & \textbf{Consecuencia práctica} \\
\midrule
Inmuebles & Declarativo & El derecho nace con el título más el modo (escritura más tradición). La inscripción solo lo hace oponible \textit{erga omnes}, es decir, conocido por toda la comunidad. \\
Automotores & Constitutivo & Sin inscripción (transferencia registral) no se es dueño del automotor: el derecho nace con la inscripción. \\
\bottomrule
\end{tabularx}
\caption{Efectos de la inscripción registral}
\end{table}

\cornell{adquisición de los derechos reales}{
\cq{Distinga adquisición originaria y derivada.}{Originaria: no deriva de un vínculo anterior (cosas sin dueño, prescripción adquisitiva, adquisición legal). Derivada: un enajenante transmite a un adquirente, entre vivos o por causa de muerte.}
\cq{¿Qué dispone el art.~1896 y qué principio expresa?}{El juez no puede constituir derechos reales ni imponer su constitución: principio de legalidad.}
\cq{¿Qué exige el art.~1892?}{Título suficiente más modo suficiente para la adquisición derivada entre vivos.}
\cq{¿Por qué el título se entiende en sentido abstracto?}{Porque es el acto jurídico causal que plasma la transmisión, no el instrumento: entregar un papel que no plasma un acto causal es no tener título.}
\cq{¿Por qué la tradición no puede suplirse por una cláusula de escritura (art.~1924)?}{Porque es un acto material; la posesión no se transmite con palabras escritas.}
\cq{Compare la inscripción registral en inmuebles y automotores.}{Inmuebles: declarativa (oponibilidad). Automotores: constitutiva (sin inscripción no hay dominio).}
}{El derecho real se adquiere por vía originaria o derivada. En la derivada entre vivos rige la teoría del título y modo (art.~1892): acto jurídico causal válido más tradición material, que no puede suplirse por fórmulas escritas. La inscripción es declarativa en inmuebles y constitutiva en automotores.}

% ------------------------------------------------------------
\chapter[Excepciones a la tradición material]{Excepciones a la tradición material}\label{cap:brevi}
% ------------------------------------------------------------

El ordenamiento admite solo dos casos en que la adquisición opera sin entrega material efectiva: la \textit{traditio brevi manu} y el constituto posesorio. Fuera de ellos, el principio de legalidad no permite apartarse del requisito de tradición. Ambos casos se explican por una misma idea: evitar un dispendio material inútil cuando la cosa ya está en poder de quien debe recibirla, o seguirá estándolo.

\section{\textit{Traditio brevi manu} (arts.~1892 y 1933)}

\begin{definicion}{\textit{traditio brevi manu}}
Es la excepción que permite que quien ya tenía la cosa como \textbf{tenedor} (por ejemplo, como locatario) pase a tenerla como \textbf{titular del derecho real} (propietario) por un acto jurídico causal, sin necesidad de tradición material, haciéndolo constar en la escritura. La relación de poder \textbf{asciende de categoría}: de tenencia a dominio.
\end{definicion}

\begin{ejemplo}{la inquilina que compra su departamento}
María alquiló durante muchos años un departamento. Se le presentó la posibilidad de comprarlo, coincidiendo con el interés del propietario en venderlo, y firmó la escritura pública de compraventa (título suficiente). Sin la \textit{traditio brevi manu}, debería contratar una mudadora, sacar sus muebles, salir en calidad de inquilina, esperar que el propietario ingresara y que luego se lo entregara como propietaria: un dispendio inútil. La ley admite que, haciéndolo constar en la escritura, María pase de tenedora (locataria) a propietaria sin ningún movimiento material.
\end{ejemplo}

Existe una variante: si el comprador es otro adquirente y María necesita seguir alquilando, puede continuar en el inmueble como inquilina del nuevo propietario. María, que tenía la cosa a nombre del propietario original, pasa a tenerla a nombre del adquirente sin movimiento material, con debida constancia notarial y fecha cierta.

\section{Constituto posesorio (arts.~1892 y 1923)}

\begin{definicion}{constituto posesorio}
Es el caso inverso: quien tenía la cosa como \textbf{propietario} la transmite a un tercero y pasa a tenerla como \textbf{tenedor} (locatario, comodatario, etc.), sin necesidad de tradición material y con la debida constancia notarial. La relación de poder \textbf{desciende de categoría}.
\end{definicion}

\begin{ejemplo}{el propietario endeudado de 2001}
En el año 2001, un señor con muchas deudas decide vender su departamento para pagarlas y propone al adquirente alquilárselo. Lo que cambia es la calidad en que se encuentra respecto de la cosa: antes era propietario, ahora es locatario. Sin el constituto posesorio, debería entregar materialmente la posesión al adquirente y luego reingresar como inquilino. La ley evita ese disparate: el propietario que vende permanece en el inmueble, pero como locatario del nuevo dueño, con debida constancia notarial y fecha cierta.
\end{ejemplo}

\section{Comparación}

\begin{table}[H]
\centering
\small
\begin{tabularx}{\textwidth}{>{\raggedright\arraybackslash}p{0.22\textwidth}>{\raggedright\arraybackslash}X>{\raggedright\arraybackslash}X}
\toprule
 & \textbf{\textit{Traditio brevi manu}} & \textbf{Constituto posesorio} \\
\midrule
Artículos & 1892 y 1933 & 1892 y 1923 \\
Situación previa & Tenedor (locatario, comodatario) & Propietario (titular del dominio) \\
Situación posterior & Propietario & Tenedor (locatario, comodatario) \\
Movimiento & Ascendente ($\uparrow$) & Descendente ($\downarrow$) \\
Ejemplo & La inquilina que compra su departamento & El propietario de 2001 que vende y queda como inquilino \\
Tradición material & No requerida & No requerida \\
Constancia & En la escritura & Notarial, con fecha cierta \\
\bottomrule
\end{tabularx}
\caption{Excepciones a la tradición material}
\end{table}
% TODO: verificar referencia presente en las fuentes originales (el art.~1933 se cita para la \textit{traditio brevi manu} en esta parte y para la obligación de restituir en el capítulo sobre efectos de las relaciones de poder; el art.~1923 se cita para el constituto posesorio y para la tradición como modo bilateral).

\cornell{excepciones a la tradición}{
\cq{¿Cuáles son los únicos casos de adquisición sin entrega material?}{La \textit{traditio brevi manu} y el constituto posesorio; fuera de ellos rige la tradición material.}
\cq{¿Qué es la \textit{traditio brevi manu}?}{Quien tenía la cosa como tenedor pasa a tenerla como propietario por un acto causal, sin entrega material, haciéndolo constar en la escritura. Es un ascenso de categoría.}
\cq{¿Qué es el constituto posesorio?}{Quien tenía la cosa como propietario la transmite y pasa a tenerla como tenedor, sin entrega material y con constancia notarial y fecha cierta. Es un descenso de categoría.}
\cq{¿Qué problema práctico resuelve cada una?}{Evitar una entrega material inútil: salir y reingresar al inmueble con otra calidad (María; el propietario de 2001).}
}{La tradición debe ser material y efectiva, salvo dos excepciones legales: la \textit{traditio brevi manu} (tenedor que pasa a propietario) y el constituto posesorio (propietario que pasa a tenedor), ambas con constancia en instrumento y sin movimiento material.}

% ------------------------------------------------------------
\chapter[Publicidad, adquisición legal, transmisión y extinción]{Publicidad, adquisición legal, transmisión y extinción}\label{cap:publicidad}
% ------------------------------------------------------------

\section{Publicidad (art.~1893)}

Los derechos reales deben ser respetados por todos y, si deben ser respetados por todos, \textbf{deben ser conocidos por todos}. De allí el principio de publicidad, regulado en el artículo 1893, que se vincula con la oponibilidad a terceros de las relaciones de poder.

\begin{table}[H]
\centering
\small
\begin{tabularx}{\textwidth}{>{\raggedright\arraybackslash}p{0.25\textwidth}>{\raggedright\arraybackslash}X}
\toprule
\textbf{Tipo de publicidad} & \textbf{Descripción} \\
\midrule
Posesoria & Desde la existencia misma del hombre en sociedad, el vínculo más claro con una cosa es tenerla: la posesión es en sí misma publicidad. Opera para toda clase de cosas. \\
Registral & La inscripción en un registro público permite que los terceros conozcan quién es el titular y hace oponible el derecho. Opera cuando se trata de cosas registrables. Su efecto es declarativo en inmuebles y constitutivo en automotores. \\
\bottomrule
\end{tabularx}
\caption{Tipos de publicidad de los derechos reales}
\end{table}

\section{Adquisición legal (art.~1894)}

Existen casos en que la ley misma dispone que el derecho se adquiere, sin que concurran título y modo en sentido estricto. Son casos de adquisición legal u originaria:

\begin{enumerate}
    \item el derecho real de \textbf{habitación del cónyuge supérstite}: al fallecer uno de los cónyuges, la ley reconoce al sobreviviente el derecho real de habitación sobre el inmueble asiento del hogar conyugal, gratuito y vitalicio, siempre que se cumplan los requisitos legales;
    \item el derecho real de \textbf{habitación del conviviente supérstite}: el Código incorporó al conviviente (unión convivencial) entre los beneficiarios de este derecho, manifestación del cambio de paradigma del nuevo cuerpo normativo, a veces llamada «el conviviente impensado»;
    \item los \textbf{condominios con indivisión forzosa} (capítulo~\ref{cap:medianeria});
    \item los derechos de los adquirentes de buena fe y a título oneroso, cuyo desarrollo se presenta en el capítulo~\ref{cap:muebles}.
\end{enumerate}
% TODO: verificar consistencia entre fuentes originales (una versión enumera tres casos de adquisición legal en el art.~1894 y otra agrega un cuarto; el régimen de cosas muebles se cita luego como «arts.~1894 a 1895»).

\section{Transmisión}

Todos los derechos, reales o personales, son en principio transmisibles. La transmisión de derechos reales se rige por la teoría del título y modo: quien transmite lo hace con título suficiente y modo, y quien adquiere los recibe de ese modo. El \textit{nemo plus iuris} es el límite infranqueable: nadie puede transmitir más derechos de los que tiene. Los casos particulares se analizan al tratar cada derecho real.

\section{Extinción (art.~1907)}

El derecho real se extingue en forma relativa o absoluta.

\begin{table}[H]
\centering
\small
\begin{tabularx}{\textwidth}{>{\raggedright\arraybackslash}p{0.17\textwidth}>{\raggedright\arraybackslash}X>{\raggedright\arraybackslash}p{0.16\textwidth}>{\raggedright\arraybackslash}p{0.22\textwidth}}
\toprule
\textbf{Tipo} & \textbf{Descripción} & \textbf{¿Continúa?} & \textbf{Ejemplo} \\
\midrule
Relativa & Se extingue para el titular actual. & Sí, en otro sujeto. & Venta, donación, dación en pago. \\
Absoluta & El derecho desaparece por completo; como el derecho real es un vínculo entre titular y cosa, destruida la cosa ese vínculo se extingue. & No. & Destrucción de la cosa. \\
Consolidación & Se reúnen en una misma persona el derecho real principal y el que lo desmembraba. & Se consolida en uno solo. & El usufructuario adquiere el dominio pleno; el nudo propietario adquiere el usufructo. \\
\bottomrule
\end{tabularx}
\caption{Extinción de los derechos reales}
\end{table}

La consolidación es análoga a lo que en los derechos personales se denomina confusión: el derecho menor queda absorbido.

\section{El derecho real como relación de poder}

El derecho real es una \textbf{relación de poder del titular respecto de la cosa}. Ese esquema, esa línea o flecha, es lo que permite perseguir la cosa, como ocurre en las hipotecas, aun cuando hubiera cambiado de titular. La relación es defendida y protegida por el sistema jurídico y debe ser respetada por toda la sociedad; de allí la necesidad de la publicidad, posesoria o registral, a pesar de los límites que la Nación decida imponer en función de los intereses de la comunidad.

Desde el nacimiento y a lo largo de toda su existencia, los seres humanos tienen contactos permanentes con las cosas. Ese contacto directo e inmediato tiene una protección jurídica que debe ser respetada, porque es también una manera de respetar la paz social.

Sin embargo, no siempre los vínculos de poder están en cabeza del titular del derecho real: no siempre las personas viven en las tierras de las que son propietarias. Existen diversas relaciones de poder y diversas calidades en que un sujeto puede vincularse con una cosa. A ellas se dedica la Parte~III.

\begin{conceptoclave}{síntesis de la adquisición y extinción}
La adquisición derivada por actos entre vivos exige título suficiente y modo suficiente (art.~1892), sin que la tradición pueda suplirse por fórmulas escritas, salvo la \textit{traditio brevi manu} y el constituto posesorio. Para su oponibilidad, los derechos reales se publicitan mediante la posesión o la inscripción registral. Se extinguen en forma relativa (continúan en otro titular) o absoluta (típicamente, por destrucción de la cosa), siempre en el marco de la preeminencia del interés colectivo que proclama el art.~240.
\end{conceptoclave}

\cornell{publicidad, adquisición legal y extinción}{
\cq{¿Por qué los derechos reales requieren publicidad?}{Porque deben ser respetados por todos y, por tanto, conocidos por todos (art.~1893). Puede ser posesoria o registral.}
\cq{¿Qué casos de adquisición legal enumera el art.~1894?}{Habitación del cónyuge supérstite, habitación del conviviente supérstite y condominios con indivisión forzosa; según una fuente, también los adquirentes de buena fe a título oneroso.}
\cq{Distinga extinción relativa y absoluta (art.~1907).}{Relativa: el derecho se extingue para un titular pero continúa en otro. Absoluta: desaparece, típicamente por destrucción de la cosa.}
\cq{¿Qué es la consolidación?}{La reunión en una persona del derecho real principal y del que lo desmembraba (usufructuario que adquiere el dominio); análoga a la confusión.}
\cq{¿Por qué el derecho real es una «relación de poder»?}{Porque es un vínculo directo del titular con la cosa, protegido por el sistema y exigible a todos; permite perseguir la cosa y exige publicidad.}
}{La publicidad (posesoria o registral) hace conocido y oponible el derecho real. La ley adquiere por sí ciertos derechos (habitación del cónyuge y del conviviente, condominios con indivisión forzosa). La extinción puede ser relativa, absoluta o por consolidación. El derecho real es una relación de poder protegida por el sistema jurídico.}

% ============================================================
\part{Relaciones de poder: posesión y tenencia}
% ============================================================

% ------------------------------------------------------------
\chapter[Las relaciones entre la persona y la cosa]{Las relaciones entre la persona y la cosa}\label{cap:poder}
% ------------------------------------------------------------

La tradición, la publicidad posesoria y la defensa de la posesión presuponen una pregunta previa: qué relaciones puede tener una persona con una cosa. Esas relaciones son \textbf{taxativas}: se denominan y regulan de manera específica y no pueden crearse otras distintas, pues serían nulas, en consonancia con el \textit{numerus clausus} del artículo 1884.

\section{Las relaciones de hecho reconocidas}

Las únicas relaciones de hecho que el ordenamiento reconoce son la \textbf{yuxtaposición local}, la \textbf{tenencia} y la \textbf{posesión}. Para la teoría posesoria interesan, según los fundamentos del anteproyecto, solo la posesión y la tenencia, con el agregado del \textbf{servidor de la posesión} al solo fin de la defensa extrajudicial. Las relaciones de poder se regulan en los artículos 1908 a 1940.

\subsection{La yuxtaposición local}

\begin{definicion}{yuxtaposición local}
Mero \textbf{contacto físico casual} entre una persona y una cosa.
\end{definicion}

Existe algo de voluntad (sin ella no habría contacto), pero no está dirigida hacia esa cosa en particular y no hacia otra.

\begin{ejemplo}{el colectivo}
Una persona sube a un colectivo, se apoya en una barra, se sienta y, cuando sube alguien que lo necesita más, cede el asiento. Hay contacto con las cosas, pero no está direccionado por el sujeto a esa cosa y no a otra.
\end{ejemplo}

\subsection{Los servidores de la posesión}

Quien utiliza una cosa en virtud de una relación de dependencia, hospedaje u hospitalidad se llama \textbf{servidor de la posesión} (art.~1911).

\begin{table}[H]
\centering
\small
\begin{tabularx}{\textwidth}{>{\raggedright\arraybackslash}p{0.17\textwidth}>{\raggedright\arraybackslash}X>{\raggedright\arraybackslash}X}
\toprule
\textbf{Vínculo} & \textbf{Ejemplo} & \textbf{Consecuencia jurídica} \\
\midrule
Dependencia & El obrero y su máquina de trabajo, el empleado y su computadora: la cuida y la usa todos los días, pero si lo despiden no se la lleva. & Quien sirve a la posesión no tiene \textit{animus} ni \textit{corpus} propio: posee en nombre del empleador. \\
Hospedaje & El huésped de un hotel al que asignan la habitación 22: si no hay agua, llama al conserje; con otra habitación, esa sería «su habitación». & El huésped no tiene posesión ni tenencia, solo un contacto derivado del servicio. \\
Hospitalidad & El amigo que visita una casa y se queda unos días: usa la cama, el baño, la heladera, y al irse no se lleva esas cosas. & El vínculo dura mientras dura la hospitalidad; no genera posesión ni tenencia. \\
\bottomrule
\end{tabularx}
\caption{Servidores de la posesión}
\end{table}

\section{Posesión, tenencia y servidor de la posesión}

\begin{definicion}{posesión (art.~1909 CCyC)}
Hay posesión cuando una persona, por sí o por medio de otra, ejerce un poder de hecho sobre una cosa, comportándose como titular de un derecho real, lo sea o no.
\end{definicion}

Tiene dos elementos: el \textbf{\textit{corpus}} (contacto material con la cosa) y el \textbf{\textit{animus}} o \textit{animus domini} (ánimo de dueño): el poseedor se comporta como si la cosa fuera propia, sin reconocer en otro un señorío superior, lo sea o no el dueño.

\begin{definicion}{tenencia (art.~1910 CCyC)}
Hay tenencia cuando una persona, por sí o por medio de otra, ejerce un poder de hecho sobre una cosa y se comporta como representante del poseedor.
\end{definicion}

Tiene un solo elemento, el \textit{corpus}: no hay \textit{animus}, porque el tenedor reconoce en otro la propiedad. El caso típico es el inquilino, que tiene relación directa con la cosa y numerosos derechos sobre ella, pero reconoce en otro la propiedad: paga el alquiler y, si le reclaman por desperfectos, responde que la propiedad no es suya.

\begin{table}[H]
\centering
\small
\begin{tabularx}{\textwidth}{>{\raggedright\arraybackslash}p{0.2\textwidth}>{\raggedright\arraybackslash}X>{\raggedright\arraybackslash}X>{\raggedright\arraybackslash}X}
\toprule
 & \textbf{Poseedor} & \textbf{Tenedor} & \textbf{Servidor} \\
\midrule
Elementos & \textit{Corpus} y \textit{animus} & Solo \textit{corpus} & Ni \textit{corpus} ni \textit{animus} propios \\
Actitud ante la cosa & Se comporta como titular, lo sea o no & Reconoce en otro la propiedad & Usa la cosa por dependencia, hospedaje u hospitalidad \\
Ejemplo & Fermín, que se lleva el teléfono & Juan, que lo tiene hasta que vuelva su dueño & El empleado y su computadora \\
\bottomrule
\end{tabularx}
\caption{Poseedor, tenedor y servidor de la posesión}
\end{table}
% TODO: verificar referencia presente en las fuentes originales (el art.~1911 se cita en una parte para el servidor de la posesión y, en otra, para la presunción de posesión).

\begin{ejemplo}{el teléfono: Juan y Fermín}
Una persona (A) espera un llamado y debe ausentarse un momento. Le pide a Juan que le tenga el teléfono y que, si llaman, avise que ya vuelve. Llega Carolina y le pide que se lo preste porque perdió el suyo; Juan responde que no puede, porque el teléfono no es suyo sino de A, y cuando A regresa, Juan espontáneamente le extiende el teléfono. \textbf{Juan es tenedor:} reconoce en otro un señorío superior.

En otra versión, A deja el teléfono sobre el escritorio y Fermín lo toma, lo guarda en su mochila, sale del aula, habla en el bar y lo lleva a su casa. Se comporta como si fuera dueño: no reconoció la propiedad de A, pues de haberla reconocido no se lo habría llevado. \textbf{Fermín es poseedor:} tiene \textit{corpus} y \textit{animus} y se comporta como propietario, lo sea o no.
\end{ejemplo}

\cornell{relaciones de poder}{
\cq{¿Cuáles son las relaciones de hecho reconocidas y por qué son taxativas?}{Yuxtaposición local, tenencia y posesión. Son taxativas por el \textit{numerus clausus}: no pueden crearse otras bajo pena de nulidad.}
\cq{¿Qué es la yuxtaposición local?}{Mero contacto físico casual con una cosa, sin voluntad dirigida a esa cosa en particular (el pasajero del colectivo).}
\cq{¿Qué es el servidor de la posesión?}{Quien usa la cosa por dependencia, hospedaje u hospitalidad. No tiene posesión ni tenencia; solo importa para la defensa extrajudicial.}
\cq{Defina posesión y tenencia por sus elementos.}{Posesión (art.~1909): \textit{corpus} más \textit{animus domini}. Tenencia (art.~1910): solo \textit{corpus}; el tenedor reconoce en otro la propiedad y se comporta como representante del poseedor.}
\cq{Aplique al teléfono: ¿quién es tenedor y quién poseedor?}{Juan es tenedor, porque reconoce el señorío de A. Fermín es poseedor, porque se lleva el teléfono comportándose como dueño.}
}{Las relaciones de poder reconocidas son taxativas. La posesión exige \textit{corpus} y \textit{animus domini}; la tenencia, solo \textit{corpus}, reconociendo en otro la propiedad; el servidor de la posesión usa la cosa por dependencia, hospedaje u hospitalidad y solo interesa para la defensa extrajudicial.}

% ------------------------------------------------------------
\chapter[Presunciones, clasificación y defensa de la posesión]{Presunciones, clasificación y defensa de la posesión}\label{cap:clasificacion}
% ------------------------------------------------------------

\section{Presunciones y reglas posesorias}

La ley resuelve la imposibilidad práctica de conocer el verdadero vínculo que cada persona tiene con las cosas que usa mediante un régimen de presunciones. Quien tiene una relación de hecho con una cosa y se comporta como titular del derecho real se presume poseedor, lo sea o no. Esto hace posible la vida en sociedad: sería imposible verificar cada relación de poder que se observa cotidianamente, por lo que, si no existe un interés legítimo en conocer el verdadero contenido jurídico de una relación, se respetan las múltiples relaciones de hecho observadas.

\begin{regla}{presunciones de posesión, legitimidad y buena fe}
\begin{itemize}
    \item \textbf{Art.~1911 (posesión):} toda relación de poder se presume posesión, excepto prueba en contrario.
    \item \textbf{Art.~1916 (legitimidad):} las relaciones de poder se presumen legítimas, a menos que exista prueba en contrario.
    \item \textbf{Art.~1919 (buena fe):} la posesión se presume de buena fe, a menos que exista prueba en contrario o que el título sea de mala fe.
\end{itemize}
\end{regla}

Otras reglas en materia posesoria:

\begin{itemize}
    \item \textbf{Innecesariedad del título (art.~1917):} el poseedor no tiene obligación de producir su título, salvo que deba exhibirlo como obligación inherente a su posesión.
    \item \textbf{Fecha y extensión (art.~1914):} si media título, se presume que la relación de poder comienza desde la fecha del título y tiene la extensión que en él se indica.
    \item \textbf{Relatividad de los vicios (art.~1921):} los vicios de la posesión son relativos respecto de aquel contra quien se ejercen; aun si la posesión es viciosa, solo puede alegar el vicio quien lo sufrió.
    \item \textbf{Inmutabilidad de la causa:} nadie puede cambiar por sí mismo la causa de su posesión. Quien comenzó como tenedor continúa siéndolo; quien era poseedor de mala fe continúa siéndolo.
\end{itemize}

\section{Clasificación de la posesión}

\begin{table}[H]
\centering
\small
\begin{tabularx}{\textwidth}{>{\raggedright\arraybackslash}p{0.2\textwidth}>{\raggedright\arraybackslash}X>{\raggedright\arraybackslash}X}
\toprule
\textbf{Tipo} & \textbf{Definición} & \textbf{Ejemplo} \\
\midrule
Legítima & Ejercicio de un derecho real adquirido conforme a la ley: título suficiente, modo suficiente y capacidad de ambas partes. & El propietario que compró con escritura pública e hizo la tradición. \\
Ilegítima de buena fe & Sin título suficiente, pero el poseedor cree que tiene derecho a poseer (error excusable). & Quien compra por boleto de compraventa (falta la escritura); quien adquiere un libro usado sin saber que era robado. \\
Ilegítima de mala fe simple & El poseedor duda de la legitimidad de su adquisición, pero igualmente adquiere. & Quien compra un reloj de alta gama a precio ínfimo en un lugar inusual, sin preguntarse por qué se lo venden tan barato. \\
Ilegítima de mala fe viciosa, muebles & Posesión adquirida mediante vicios específicos de las cosas muebles. & Hurto, estafa, abuso de confianza. \\
Ilegítima de mala fe viciosa, inmuebles & Posesión adquirida mediante vicios específicos de los inmuebles. & Clandestinidad, violencia, abuso de confianza. \\
\bottomrule
\end{tabularx}
\caption{Clasificación de la posesión}
\end{table}

El vicio común a muebles e inmuebles es el \textbf{abuso de confianza}.

\begin{ejemplo}{abuso de confianza en un inmueble}
Un propietario presta a un amigo un inmueble desocupado que va a vender pronto. Cuando le informa que lo vendió y que debe entregar la posesión, el amigo responde que él es poseedor, que tiene \textit{animus} y que allí se queda. Hay una intervención del título basada en un abuso de confianza (véase el capítulo~\ref{cap:dinamica}).
\end{ejemplo}

\begin{importante}
Lo que se califica es el \textbf{vínculo}, no la persona. El derecho califica de legítima o ilegítima la relación que la persona tiene con la cosa, no su rostro, su nombre ni su historia: una persona honesta puede tener una posesión ilegítima respecto de una cosa, y un delincuente puede ser poseedor legítimo del inmueble que adquirió con título suficiente y tradición.
\end{importante}

\section{La defensa de la posesión}

La posesión se protege aunque sea ilegítima o viciosa, por razones de orden público:

\begin{enumerate}
    \item para proteger a la persona del poseedor, el anterior y el nuevo;
    \item para proteger la paz social, evitando la violencia privada: quien tiene derecho a la posesión debe recuperarla por vía legal, no por la fuerza;
    \item incluso el verdadero propietario que toma la posesión por la fuerza puede ser demandado por el poseedor despojado.
\end{enumerate}

La carga de la prueba recae siempre en quien controvierte la relación de poder y no en el poseedor: su posesión se presume legítima y de buena fe.

\cornell{presunciones, clasificación y defensa}{
\cq{¿Qué presunciones rigen y para qué sirven?}{De posesión (art.~1911), de legitimidad (1916) y de buena fe (1919). Permiten la vida social sin verificar cada relación de poder; la prueba contraria corresponde a quien controvierte.}
\cq{¿Qué significa que los vicios sean relativos (art.~1921)?}{Solo puede alegar el vicio quien lo sufrió.}
\cq{¿Qué quiere decir que nadie puede cambiar por sí la causa de su posesión?}{Que el tenedor sigue siendo tenedor y el poseedor de mala fe sigue siéndolo, aunque pase el tiempo.}
\cq{Clasifique la posesión.}{Legítima; ilegítima de buena fe; de mala fe simple (duda); de mala fe viciosa (hurto, estafa, abuso de confianza en muebles; clandestinidad, violencia, abuso de confianza en inmuebles).}
\cq{¿Qué se califica: la persona o el vínculo?}{El vínculo con la cosa.}
\cq{¿Por qué se defiende incluso la posesión viciosa?}{Para proteger a la persona y la paz social, evitando la violencia privada.}
}{Las relaciones de poder se presumen posesión, legítima y de buena fe. La posesión se clasifica según su legitimidad y vicios; se califica el vínculo y no a la persona. Se protege incluso la posesión ilegítima para preservar la paz social.}

% ------------------------------------------------------------
\chapter[Dinámica de las relaciones de poder]{Dinámica de las relaciones de poder}\label{cap:dinamica}
% ------------------------------------------------------------

Este capítulo sigue la vida de una relación de poder: cómo se adquiere, qué efectos produce, cómo puede alterarse su causa, cómo se acumula en el tiempo y cómo se extingue.

\section{Adquisición}

\begin{regla}{art.~1922 CCyC, adquisición de poder}
Para adquirir una relación de poder sobre una cosa, esta debe ser establecida voluntariamente: a) por sujeto capaz, excepto las personas menores de edad, para quienes es suficiente que tengan diez años; b) por medio de contacto con la cosa, de la posibilidad física de establecerlo, o cuando ella ingresa en el ámbito de custodia del adquirente.
\end{regla}

Además, la cosa debe estar libre de toda otra relación excluyente (\textbf{posesión vacua}, art.~1926). Los modos de adquisición son:

\begin{itemize}
    \item \textbf{unilateral:} el apoderamiento;
    \item \textbf{bilateral:} la tradición, por excelencia (art.~1923): una persona (el tradente) entrega la cosa a otra, que ocupa la posición que tenía el tradente;
    \item \textbf{originaria:} la prescripción adquisitiva, entre otros modos.
\end{itemize}

\begin{table}[H]
\centering
\small
\begin{tabularx}{\textwidth}{>{\raggedright\arraybackslash}X>{\raggedright\arraybackslash}X}
\toprule
\textbf{Tradición traslativa de dominio} & \textbf{Tradición posesoria (mera)} \\
\midrule
Requiere: (1) que la realice el legitimado, es decir, el titular del derecho; (2) título suficiente (acto jurídico causal); (3) capacidad de ambas partes. & Si no se dan esos requisitos, se transmite meramente la relación de poder que el tradente tiene con la cosa. Nadie puede transmitir un derecho mejor o más extenso que el que tenía. \\
\bottomrule
\end{tabularx}
\caption{Tradición traslativa y tradición posesoria}
\end{table}

Si un poseedor no titular transmite la posesión, solo transmite el hecho de la posesión; pero si el propietario (legitimado, capaz y con título suficiente) transmite, lo que transmite es el derecho de poseer.

\section{Efectos}

\subsection{Derechos inherentes a la posesión}

Son los que tiene todo aquel que se encuentra en poder de una cosa, sea poseedor o tenedor; por ejemplo, los límites al dominio.

\begin{ejemplo}{ventana en pared medianera}
Un vecino, como propietario, tiene amplias facultades, pero también límites: no puede abrir una ventana de vista en la pared medianera, porque quien está del otro lado tiene derecho a que no se construya. Obligación y derecho son las dos caras de una misma moneda.
\end{ejemplo}

\subsection{Obligaciones del poseedor y del tenedor}

\begin{regla}{art.~1933 CCyC}
El poseedor o tenedor debe restituir la cosa a quien tenga derecho de reclamarla, aunque no se la haya dado a él.
\end{regla}

Quien tiene una relación de poder con la cosa y no tiene derecho real tiene la \textbf{obligación de restituirla} y el \textbf{derecho a que le sean restituidas las mejoras} realizadas en su conservación.

\subsection{La \textit{laudatio auctoris}}

El tenedor tiene la \textbf{obligación legal de identificar al poseedor} (\textit{laudatio auctoris}) cuando recibe una demanda dirigida al propietario, bajo apercibimiento de daños y perjuicios.

\begin{ejemplo}{el inquilino que no avisa}
El inmueble está ocupado por un inquilino, que es tenedor. Recibe una demanda de reivindicación y sostiene que, como no es el titular, la demanda no lo afecta: no la contesta ni avisa al propietario. La conducta es gravísima, porque priva al propietario de la posibilidad de defenderse en juicio.
\end{ejemplo}

\section{Intervención del título}

La regla general es que nadie puede cambiar por sí mismo la especie de su relación de poder: ni el tiempo, ni la voluntad interna, ni el pensamiento modifican el título.

\begin{itemize}
    \item \textbf{Intervención unilateral:} el tenedor, por actos exteriores, manifiesta su voluntad de privar al poseedor de la cosa.
    \item \textbf{Intervención bilateral:} existe acuerdo entre poseedor y tenedor (art.~1915).
\end{itemize}

\begin{ejemplo}{intervención unilateral}
El propietario presta su propiedad a un amigo para que viva allí hasta la venta. Cuando se produce la venta, le pide que desocupe; el amigo responde que supuso otra cosa y que ahora posee, porque posee y su posesión se presume legítima. Ha intervertido unilateralmente el título. Del mismo modo, quien recibió el teléfono en tenencia afirma ante quien se lo entregó que nadie se lo dejó y que es suyo: hay un acto exterior. No hay conversión natural por el paso del tiempo ni por haber pensado de otro modo.
\end{ejemplo}

\section{Unión de posesiones}

\begin{regla}{art.~1901 CCyC}
El heredero continúa la posesión de su causante. El sucesor particular puede unir su posesión a la de sus antecesores, siempre que derive inmediatamente de las otras. En la prescripción breve las posesiones unidas deben ser de buena fe y estar unidas por un vínculo jurídico.
\end{regla}

La unión de posesiones es una \textbf{suma temporal} del tiempo de dos posesiones de distintos poseedores, útil para alcanzar el número de años requerido para la \textbf{prescripción adquisitiva}.

\begin{ejemplo}{unión de posesiones}
Un poseedor que posee desde hace diez años cede su posesión por escritura pública; el nuevo poseedor acumula ese tiempo al propio para alcanzar los veinte años de la prescripción larga.
\end{ejemplo}

\section{Actos posesorios}

\begin{regla}{art.~1928 CCyC}
Constituyen actos posesorios sobre la cosa inmueble: su cultura, percepción de frutos, amojonamiento o impresión de signos materiales, mejora, exclusión de terceros y, en general, su apoderamiento por cualquier modo que se obtenga.
\end{regla}

Son actos propios de quien ejerce el \textit{animus domini} y que un tenedor difícilmente realizaría, como construir, hacer mejoras o plantaciones en un inmueble ajeno que va a dejar al cabo del período locativo, salvo que se trate de un acto conservatorio. Quien invoca la posesión debe acreditar que tuvo la disposición física del inmueble y que esa relación era la propia de un poseedor, mediante actos de intervención material sobre la cosa.

\section{Extinción}

\begin{regla}{art.~1931 CCyC}
La posesión y la tenencia se extinguen cuando el sujeto que la ejerce: a) se extingue la cosa; b) otro priva al titular de la cosa; c) se ha imposibilitado el ejercicio de la posesión en forma perdurable; d) por el abandono.
\end{regla}

\begin{ejemplo}{abandono y pérdida}
Un colchón en la pila de residuos se presume abandonado, pues nadie guarda un colchón así: se extingue la relación de poder por abandono. Si un teléfono cae al mar mientras su poseedor saca una foto, es imposible recuperarlo y la relación de poder se extingue.
\end{ejemplo}

\cornell{dinámica de las relaciones de poder}{
\cq{¿Qué se exige para adquirir una relación de poder (arts.~1922 y 1926)?}{Voluntad de un sujeto capaz (diez años para menores), contacto con la cosa, posibilidad física de establecerlo o ingreso en el ámbito de custodia, y que la cosa esté libre de otra relación excluyente (posesión vacua).}
\cq{Distinga tradición traslativa de dominio y tradición posesoria.}{La traslativa requiere legitimado, título suficiente y capacidad. Sin ellos se transmite solo la relación de poder del tradente.}
\cq{¿Qué obligaciones tienen el poseedor y el tenedor sin derecho real (art.~1933)?}{Restituir la cosa a quien pueda reclamarla y derecho al reembolso de mejoras de conservación. El tenedor debe además identificar al poseedor (\textit{laudatio auctoris}).}
\cq{Distinga intervención unilateral y bilateral del título.}{Unilateral: el tenedor, con actos exteriores, pretende privar al poseedor. Bilateral: acuerdo entre ambos (art.~1915). El tiempo o el pensamiento no cambian el título.}
\cq{¿Para qué sirve la unión de posesiones (art.~1901)?}{Sumar el tiempo de posesiones sucesivas para la prescripción adquisitiva; en la breve, deben ser de buena fe y unidas por vínculo jurídico.}
\cq{¿Cómo se extinguen la posesión y la tenencia (art.~1931)?}{Por extinción de la cosa, privación por otro, imposibilidad perdurable de ejercicio y abandono.}
}{La relación de poder se adquiere voluntariamente por un sujeto capaz mediante contacto o custodia, sobre una cosa libre de otra relación excluyente. Genera derechos inherentes y obligaciones de restitución, no cambia de causa por voluntad propia salvo intervención del título, puede sumarse a la de los antecesores y se extingue por pérdida de la cosa, privación, imposibilidad o abandono.}

% ------------------------------------------------------------
\chapter[Adquisición legal sobre cosas muebles: art.~1895]{Adquisición legal de derechos reales sobre cosas muebles}\label{cap:muebles}
% ------------------------------------------------------------

Este régimen, de gran interés en materia de orden público, resuelve un conflicto entre dos intereses legítimos: el del propietario que quiere conservar su cosa y el del tercero que la adquirió confiando en la apariencia.

\section{Marco: clasificación de las cosas}

Las cosas son muebles o inmuebles. Conforme al art.~227, son muebles las que pueden desplazarse, por sí mismas o por una fuerza externa. Dentro de las muebles se distinguen las registrables y las no registrables. El régimen estudiado aquí corresponde a las muebles \textbf{no registrables}.

\section{El art.~1895}

\begin{regla}{art.~1895 CCyC}
La posesión de buena fe del subadquirente de cosas muebles no registrables que no sean hurtadas o perdidas es suficiente para adquirir los derechos reales principales, excepto que el verdadero propietario pruebe que la adquisición fue gratuita.
\end{regla}

Requisitos: \textbf{posesión de buena fe} del subadquirente y que la cosa \textbf{no sea hurtada ni perdida}. Con solo ellos, por decisión legal, el poseedor adquiere la propiedad y la defiende frente a todos.

\begin{importante}
El \textbf{título oneroso no es requisito} de la adquisición: el poseedor adquiere la propiedad la haya recibido a título oneroso o gratuito. El título oneroso solo importa para repeler la acción reivindicatoria del anterior propietario, y la carga de probar la gratuidad recae sobre quien reivindica.
\end{importante}

La norma se aplica cuando el propietario se desprende voluntariamente de la cosa, aunque su voluntad hubiera estado viciada. Quien recibió la cosa fue tenedor y, defraudando la confianza, la transmitió a un tercero, el \textbf{subadquirente}.

\begin{ejemplo}{el libro de Derechos Reales}
Un propietario presta un libro de Derechos Reales a una persona, quien, defraudando su confianza, lo vende o regala a Juan. Hay tradición y Juan tiene la posesión, de buena fe; la cosa no fue robada ni perdida, porque el propietario la entregó voluntariamente. Cuando el propietario se encuentra con Juan e identifica el libro por una dedicatoria, pretende reivindicarlo, y Juan se lo niega.

Cada parte tiene su verdad, pero no es posible multiplicar los libros; hay que decidir. El artículo 1895, en consonancia con la mayoría de los ordenamientos modernos, \textbf{sacrifica al propietario y protege al subadquirente}. La única posibilidad del propietario es probar que la adquisición de Juan fue gratuita, lo que presenta dificultades fácticas y jurídicas.
\end{ejemplo}

\section{Seguridad estática y seguridad dinámica}

\begin{table}[H]
\centering
\small
\begin{tabularx}{\textwidth}{>{\raggedright\arraybackslash}X>{\raggedright\arraybackslash}X}
\toprule
\textbf{Seguridad estática} & \textbf{Seguridad dinámica} \\
\midrule
Protege al propietario actual, a quien ya tiene el derecho real y desea conservarlo (en el ejemplo, el dueño del libro). & Protege la circulación de los bienes y la seguridad de las transacciones, a quien adquiere confiando en la apariencia (Juan, el subadquirente de buena fe). \\
Favorece la estabilidad del derecho ya constituido. & Favorece el tráfico jurídico y la circulación rápida de las cosas muebles, que no dejan rastro instrumental. \\
El art.~1895 la sacrifica. & El art.~1895 la privilegia, elección casi uniforme en los sistemas occidentales modernos. \\
\bottomrule
\end{tabularx}
\caption{Seguridad estática y dinámica en las cosas muebles}
\end{table}

El Código adopta así la máxima «\textbf{posesión vale título}»: adquisición legal de derechos reales sobre muebles por subadquirentes de buena fe cuando la cosa no ha sido hurtada ni perdida. La misma tensión se presenta en los inmuebles, con otra solución, en el capítulo~\ref{cap:adq-dominio}.

\cornell{adquisición legal sobre muebles}{
\cq{¿Qué requisitos exige el art.~1895?}{Posesión de buena fe del subadquirente y que la cosa no sea hurtada ni perdida; se aplica a muebles no registrables cuando el propietario se desprendió voluntariamente.}
\cq{¿Es requisito el título oneroso?}{No para adquirir. Solo importa para repeler la reivindicación: el propietario puede recuperar si prueba que la adquisición fue gratuita.}
\cq{¿Qué hace el art.~1895 con los dos intereses en conflicto?}{Sacrifica al propietario (seguridad estática) y protege al subadquirente (seguridad dinámica).}
\cq{¿Qué significa «posesión vale título»?}{El subadquirente de buena fe de una cosa no hurtada ni perdida adquiere el derecho real por la sola posesión.}
}{El art.~1895 privilegia la seguridad de las transacciones: el subadquirente de buena fe de una cosa mueble no registrable, no hurtada ni perdida, adquiere el derecho real principal, salvo que el propietario pruebe que adquirió gratuitamente.}

% ============================================================
\part{El dominio}
% ============================================================

% ------------------------------------------------------------
\chapter[Concepto, caracteres y clasificación del dominio]{Concepto, caracteres y clasificación del dominio}\label{cap:dominio}
% ------------------------------------------------------------

Quien es dueño de una casa puede, como propietario, hacer casi todo con ella: venderla, alquilarla, reformarla. Esa propiedad puede ser plena o estar afectada por una hipoteca o un usufructo. Cabe preguntarse además cómo se llegó a ser propietario, por compra con escritura o por herencia. Y, aun siendo dueño, no es posible hacer ruido hasta las tres de la mañana, tapar la ventana del vecino ni bloquear el río. Ese es el recorrido de esta parte: concepto y caracteres del dominio, su clasificación, su adquisición y, finalmente, sus límites.

\section{Las disposiciones comunes aplicadas al dominio}

El CCyC utiliza, además de los principios rectores aplicables a todo el ordenamiento, partes generales y disposiciones comunes que no se repiten en la regulación de cada derecho real. Por eso, cada vez que se aborda un derecho real particular, es metodológicamente fundamental «exponer en acto» lo aprendido sobre las disposiciones comunes, lo que implica aplicar la nueva metodología del Código en la actividad cotidiana. Aplicado al dominio, el ejercicio consiste en afirmar que:

\begin{itemize}
    \item es un \textbf{derecho real} y, por tanto, un poder jurídico, de \textbf{estructura legal};
    \item se ejerce directamente sobre un objeto y, muy excepcionalmente, sobre un derecho cuando la ley específicamente lo indique;
    \item sus normas son de \textbf{orden público} y no pueden ser modificadas por los particulares, bajo pena de nulidad;
    \item es un derecho real sobre \textbf{cosa propia}, que se ejerce por la \textbf{posesión};
    \item se adquiere con \textbf{título suficiente y modo}, y puede adquirirse en forma \textbf{originaria o derivada}.
\end{itemize}

\section{Dominio y propiedad}

Para determinar si «dominio» y «propiedad» son sinónimos, y qué protege la CN cuando garantiza el derecho de propiedad, hay que distinguir dos acepciones.

\begin{table}[H]
\centering
\small
\begin{tabularx}{\textwidth}{>{\raggedright\arraybackslash}p{0.23\textwidth}>{\raggedright\arraybackslash}X>{\raggedright\arraybackslash}p{0.17\textwidth}}
\toprule
\textbf{Concepto} & \textbf{Alcance} & \textbf{Referencia} \\
\midrule
Propiedad en sentido amplio & Abarca derechos personales y reales: los bienes materiales e inmateriales que integran el patrimonio. & CN, arts.~14 y 17 \\
Propiedad en sentido restringido (técnico) & Es el derecho real de dominio: el más extenso y completo que puede recaer sobre una cosa determinada. & CCyC, art.~1941 \\
\bottomrule
\end{tabularx}
\caption{Propiedad en sentido amplio y restringido}
\end{table}

El art.~14 de la CN garantiza a todos los habitantes usar y disponer de la propiedad, y el art.~17 declara que la propiedad es inviolable. La CN se refiere a la propiedad en sentido amplio: protege tanto derechos personales como reales.

\section{Regulación normativa}

El dominio se regula en tres niveles: el Título Tercero del Libro Cuarto (arts.~1941 a 1982), las disposiciones generales en materia de derechos reales (arts.~1882 a 1907) y los principios generales del ordenamiento (arts.~10, 11, 12 y 14).

\section{Concepto}

\begin{definicion}{dominio perfecto (art.~1941 CCyC)}
El dominio es el derecho real en virtud del cual una cosa se encuentra sometida a la voluntad y a la acción de una persona. El dominio pleno o perfecto es el que reúne las facultades de usar, gozar y disponer material y jurídicamente de una cosa, dentro de los límites previstos por la ley. Se presume perfecto hasta que se pruebe lo contrario.
\end{definicion}

Es el derecho más extenso que puede tener una persona sobre una cosa. Los demás derechos reales no son sino fragmentos del dominio: \textbf{desmembraciones} que el titular, que reúne todas las facultades, va constituyendo al separar alguna y conferirla a otro, mediante derechos de disfrute o de garantía. Cuando se habla de «dominio» sin más, se hace referencia al dominio perfecto, pues el art.~1941 lo presume perfecto.

\section{Caracteres del dominio perfecto}

El dominio perfecto es absoluto, exclusivo y perpetuo.

\subsection{Absoluto}

Otorga todas las facultades (usar, gozar y disponer, material y jurídicamente) reunidas en una misma persona. No significa que el titular pueda hacer «lo que quiere» sin restricción: el ejercicio no debe ser abusivo y está limitado por normas de orden público, en interés de la comunidad o de la vecindad (capítulo~\ref{cap:limites}).

\subsection{Exclusivo}

\textbf{Art.~1943:} el dominio es exclusivo y no puede tener más de un titular. Sin exclusividad no hay dominio: si dos personas tuvieran un derecho real sobre la misma cosa, sería \textbf{condominio}, un derecho rotundamente diferente (Parte~V). \textbf{Art.~1944 (facultad excluyente):} el propietario puede excluir a cualquiera del uso y goce de la cosa, remover los objetos puestos en ella y encerrar los inmuebles con muros, cercos o fosos.

\subsection{Perpetuo}

No tiene límites en el tiempo y subsiste con independencia de que el titular lo ejerza o no; se perpetúa incluso más allá de la vida humana, razón por la cual existe el sistema de sucesiones. Quien poseyera la cosa por el tiempo y con las calidades legales adquirirá el dominio por prescripción adquisitiva.

\section{Facultades del titular}

\begin{itemize}
    \item \textbf{Disponer jurídicamente:} incluye \textbf{enajenar} (transmitir) a título oneroso (venta, permuta) o gratuito (donación, dación en pago); «enajenar» no es sinónimo de vender.
    \item \textbf{Constituir derechos reales} sobre la cosa: usufructo, hipoteca, servidumbre.
    \item \textbf{Administrar:} por ejemplo, dar la cosa en comodato, gratuito u oneroso.
    \item \textbf{Abandonar:} si recae sobre inmuebles, debe hacerse por escritura pública.
    \item \textbf{Disponer materialmente:} construir, destruir (con los límites legales) y transformar.
    \item \textbf{Usar y gozar:} obtener los frutos, civiles (locaciones) o naturales (cosechas).
\end{itemize}

\begin{ejemplo}{el usufructo}
El titular de dominio tiene todas las facultades. Al constituir un usufructo, «saca» de lo que tiene el uso y el goce y los pone en el usufructuario, que será titular de un derecho real sin ser propietario, con la fortaleza de los derechos reales: orden público, publicidad, contenido patrimonial y protección legal. Tiene todo lo que tiene el titular de dominio, menos la propiedad.
\end{ejemplo}

\section{Extensión (art.~1945)}

El dominio de una cosa comprende los objetos que forman un todo con ella o son sus accesorios. El dominio de un inmueble se extiende al subsuelo y al espacio aéreo, y se presume que pertenecen al dueño todas las construcciones, plantaciones y forestaciones existentes (salvo propiedad horizontal o derecho de superficie). Como decían los romanos, el dueño lo es «desde el cielo hasta el infierno».

\begin{ejemplo}{accesorios y construcciones}
Quien compra un inmueble adquiere las cañerías de agua y gas incrustadas en sus paredes; quien compra un vehículo adquiere las ruedas, aunque separadas sean cosas muebles. Todo lo construido en el terreno pertenece al dueño del terreno.
\end{ejemplo}

\section{Clasificación (art.~1946)}

\begin{table}[H]
\centering
\small
\begin{tabularx}{\textwidth}{>{\raggedright\arraybackslash}p{0.17\textwidth}>{\raggedright\arraybackslash}X>{\raggedright\arraybackslash}p{0.15\textwidth}>{\raggedright\arraybackslash}p{0.18\textwidth}}
\toprule
\textbf{Tipo} & \textbf{Característica} & \textbf{Afecta el carácter} & \textbf{Normativa} \\
\midrule
Perfecto & Sin gravamen, sin plazo, sin condición & Ninguno & Art.~1941 \\
Revocable & Sometido a plazo o condición resolutoria (máximo diez años) & Perpetuo & Arts.~1946 y 1965 \\
Fiduciario & Ejercido durante un plazo para luego transmitir & Perpetuo & Arts.~1701 y ss. \\
Desmembrado & Gravado con un derecho real a favor de un tercero & Absoluto & Normas propias de cada derecho \\
\bottomrule
\end{tabularx}
\caption{Tipos de dominio}
\end{table}

El art.~1946 establece que los tipos de dominio imperfecto son solo tres y que son de orden público.

\begin{itemize}
    \item \textbf{Revocable (art.~1965):} sometido a condición o plazo resolutorios, con máximo de diez años. Cumplida la condición o vencido el plazo, el dueño debe transmitir la cosa a quien se la transmitió originariamente. Si a los diez años no se revoca, queda convertido en dominio perfecto.
    \item \textbf{Fiduciario (arts.~1701 y ss.):} regulado en el Libro Tercero junto con el contrato de fideicomiso; lo ejerce el adquirente durante un tiempo, al cabo del cual debe transmitir la propiedad. El \textbf{fideicomiso} es el contrato; el \textbf{dominio fiduciario} es el derecho real.
    \item \textbf{Desmembrado:} gravado con un derecho real a favor de un tercero (usufructo, hipoteca); se rige por las normas de cada derecho constituido y afecta el carácter absoluto.
\end{itemize}

\cornell{concepto y caracteres del dominio}{
\cq{¿Dominio y propiedad son sinónimos?}{No. La propiedad en sentido amplio (CN, arts.~14 y 17) abarca derechos personales y reales; en sentido restringido es el dominio (art.~1941).}
\cq{Defina el dominio perfecto (art.~1941).}{Derecho real por el que una cosa está sometida a la voluntad y acción de una persona, con facultades de usar, gozar y disponer material y jurídicamente, dentro de los límites legales; se presume perfecto.}
\cq{¿Cuáles son sus caracteres y qué significa cada uno?}{Absoluto (reúne todas las facultades, sin ser ilimitado), exclusivo (un solo titular; si no, es condominio; art.~1943) y perpetuo (sin límite temporal).}
\cq{¿Qué comprende la extensión del dominio (art.~1945)?}{Accesorios, subsuelo, espacio aéreo y construcciones, plantaciones y forestaciones, salvo propiedad horizontal o superficie.}
\cq{Distinga dominio revocable, fiduciario y desmembrado.}{Revocable: plazo o condición resolutoria de hasta diez años. Fiduciario: se ejerce por un plazo para luego transmitir. Desmembrado: gravado con derecho real de un tercero.}
}{El dominio es el derecho real más extenso, que reúne las facultades de usar, gozar y disponer; es absoluto, exclusivo y perpetuo, se extiende a accesorios, subsuelo, espacio aéreo y lo construido, y se clasifica en perfecto y tres tipos de imperfecto.}

% ------------------------------------------------------------
\chapter[Adquisición del dominio]{Adquisición del dominio}\label{cap:adq-dominio}
% ------------------------------------------------------------

\section{Modos de adquisición}

Los modos de adquirir el dominio son los hechos y actos jurídicos que la ley prevé para adjudicar en propiedad una cosa mueble o inmueble a una persona. Los hechos ilícitos (hurto, robo) no son modos reconocidos. Los artículos 1947 a 1963 regulan los distintos modos, con casuística específica: animales de caza, enjambres, pesca, acrecimiento de inmuebles por hechos de la naturaleza (aluvión, avulsión), plantaciones, etc.

\begin{table}[H]
\centering
\small
\begin{tabularx}{\textwidth}{>{\raggedright\arraybackslash}p{0.15\textwidth}>{\raggedright\arraybackslash}p{0.16\textwidth}>{\raggedright\arraybackslash}X>{\raggedright\arraybackslash}p{0.2\textwidth}}
\toprule
\textbf{Criterio} & \textbf{Subtipo} & \textbf{Descripción} & \textbf{Ejemplo} \\
\midrule
Por el acto & Entre vivos & Transmisión en vida de ambas partes & Compraventa, donación \\
Por el acto & Mortis causa & Transmisión por fallecimiento & Herencia \\
Por el objeto & Título universal & Se adquieren todos los derechos del causante & Herencia universal \\
Por el objeto & Título singular & Se adquiere un bien determinado & Legado de cosa cierta \\
Por la causa & Originario & La ley permite tomar la cosa por sí mismo & Apropiación, prescripción \\
Por la causa & Derivado & Se transmite de un titular anterior & Tradición (título y modo) \\
\bottomrule
\end{tabularx}
\caption{Clasificación de los modos de adquisición del dominio}
\end{table}

\textbf{Apropiación (art.~1947).} Son susceptibles de apropiación las cosas muebles sin dueño o abandonadas por su dueño. Los animales domésticos y domesticados no son apropiables.

\textbf{Modos derivados.} Remiten a la teoría del título y modo (capítulo~\ref{cap:adquisicion}); el modo derivado por excelencia es la tradición, y uno sin el otro elemento no basta. Si un poseedor no titular transmite la posesión, solo transmite el hecho de la posesión; si lo hace el propietario con capacidad y título suficiente, transmite el derecho de poseer.

\section{Protección del subadquirente de buena fe y a título oneroso}

Este tema refleja la fortaleza del derecho real: ante varias personas con derechos sobre una misma cosa, determina a priori hacia quién se inclina la protección de la ley.

\subsection{Regla y excepciones}

\textbf{Art.~399:} nadie puede transmitir a otro un derecho mejor ni más extenso que el que tiene, salvo las excepciones expresamente establecidas. La expresión «salvo» abre la puerta a excepciones; una es la \textbf{convalidación} (art.~1885), que en este desarrollo se describe así: si alguien transmite una cosa de la que no era propietario y luego la adquiere antes de que se declare la nulidad del acto, la transmisión se convalida. % TODO: verificar consistencia entre fuentes originales (esta descripción del art.~1885, referida a muebles no registrables y a la nulidad declarada, difiere de la redacción transcripta en el capítulo sobre legalidad).

\subsection{El subadquirente de inmuebles (art.~392)}

\begin{regla}{art.~392 CCyC}
Todos los derechos reales o personales transmitidos a terceros sobre un inmueble o muebles registrables por una persona que ha resultado adquirente en virtud de un acto nulo quedan sin ningún valor, pudiendo serle opuestos por quienes tengan interés legítimo, excepto respecto de los subadquirentes de buena fe y a título oneroso. Los subadquirentes no pueden ampararse en su buena fe si el acto se ha realizado sin intervención del titular del derecho.
\end{regla}

\begin{ejemplo}{Diego, Miriam y Juan}
Diego (enajenante) vende a Miriam (adquirente). Miriam vende luego a Juan (subadquirente). Después se declara nulo el acto entre Diego y Miriam (por ejemplo, por incapacidad de Diego). ¿Qué sucede con los derechos de Juan?
\end{ejemplo}

\begin{table}[H]
\centering
\small
\begin{tabularx}{\textwidth}{>{\raggedright\arraybackslash}p{0.2\textwidth}>{\raggedright\arraybackslash}X>{\raggedright\arraybackslash}X}
\toprule
\textbf{Seguridad} & \textbf{A quién protege} & \textbf{En el ejemplo} \\
\midrule
Estática & Al propietario original, quien tiene el derecho real & Protege a Diego: el inmueble volvería a él \\
Dinámica & Al adquirente en el tráfico, quien confió en la apariencia & Protege a Juan: consolida su adquisición \\
Opción del art.~392 & Seguridad dinámica, para el subadquirente de buena fe y título oneroso & Juan hace suyo el inmueble \\
\bottomrule
\end{tabularx}
\caption{Seguridad estática y dinámica en los inmuebles}
\end{table}

Para que Juan quede protegido se requiere: haber adquirido de \textbf{buena fe}, a \textbf{título oneroso}, y que el enajenante (Diego) haya \textbf{participado} de alguna manera en el acto después declarado nulo. En ningún caso hay protección para Miriam, el adquirente intermedio, que siempre sufre los efectos de la nulidad.

\begin{importante}
\textbf{Excepción a la excepción:} si hubo sustitución de personalidad (alguien se hizo pasar por Diego), la protección \textbf{no} alcanza a Juan, aunque haya actuado de buena fe y a título oneroso, porque el enajenante debe haber participado en el acto nulo.
\end{importante}

\begin{nota}
En el ejercicio profesional se actúa sabiendo que, en infinidad de situaciones, no es posible atender la justicia de cada uno. Ante la tensión entre derechos todos atendibles, la norma se ve forzada a decidir a priori sobre quién recae la protección, con la mirada puesta en el orden público y el interés de la comunidad. La misma elección en materia de cosas muebles se estudió en el capítulo~\ref{cap:muebles}.
\end{nota}

\cornell{adquisición del dominio}{
\cq{¿Qué son los modos de adquisición y cuáles no se reconocen?}{Hechos y actos jurídicos previstos por la ley para adjudicar una cosa en propiedad; los hechos ilícitos (hurto, robo) no lo son.}
\cq{Clasifique los modos de adquisición.}{Por el acto (entre vivos o mortis causa), por el objeto (título universal o singular) y por la causa (originario o derivado).}
\cq{¿Qué es la apropiación (art.~1947)?}{Modo originario sobre cosas muebles sin dueño o abandonadas; los animales domésticos y domesticados no son apropiables.}
\cq{En el caso Diego–Miriam–Juan, ¿qué requisitos protegen a Juan (art.~392)?}{Buena fe, título oneroso y participación del enajenante en el acto luego anulado. No se protege a Juan si hubo sustitución de personalidad; Miriam nunca está protegida.}
\cq{¿Qué tipo de seguridad privilegia el art.~392?}{La dinámica: la del subadquirente que confió en la apariencia, en detrimento del propietario original.}
}{El dominio se adquiere por modos originarios (apropiación, prescripción) o derivados (título y modo). El subadquirente de inmuebles de buena fe y a título oneroso queda protegido frente a la nulidad del acto anterior si el titular participó en él, con prioridad de la seguridad dinámica.}

% ------------------------------------------------------------
\chapter[Límites al dominio]{Límites al dominio}\label{cap:limites}
% ------------------------------------------------------------

\section{¿Es ilimitado el dominio?}

Si el dominio es el derecho más extenso, cabe preguntarse si es ilimitado. El ejercicio del dominio, en Argentina y en la gran mayoría de los sistemas occidentales, reconoce infinidad de límites, tanto en interés de los vecinos como de la comunidad. Esos límites se han ido ampliando y continúan modelándose en armonía con los derechos humanos, nuevo eje rector de coherencia del ordenamiento nacional.

\section{Evolución histórica}

El art.~2513 del Código de Vélez Sarsfield otorgaba al titular la facultad de desnaturalizar la cosa y degradarla; sin embargo, en una nota al art.~2506, el propio Vélez citaba una ley de Partidas según la cual el dominio debe ejercerse «según Dios y según fuero». La facultad de destruir la cosa fue suprimida por la Ley~17.711, que incorporó la noción de \textbf{abuso del derecho}.

\begin{table}[H]
\centering
\small
\begin{tabularx}{\textwidth}{>{\raggedright\arraybackslash}p{0.22\textwidth}>{\raggedright\arraybackslash}X>{\raggedright\arraybackslash}X}
\toprule
\textbf{Momento} & \textbf{Regulación} & \textbf{Criterio de abuso} \\
\midrule
Ley~17.711 & Introdujo la noción de abuso del derecho & Cuando no se tenían en cuenta las razones del legislador al dictar la norma \\
CCyC (art.~10) & Modela todo el ordenamiento & Cuando el ejercicio contraría los principios en que se asienta el sistema normativo \\
\bottomrule
\end{tabularx}
\caption{Evolución del abuso del derecho}
\end{table}

\section{La función social de la propiedad}

El art.~21 del Pacto de San José de Costa Rica reconoce el derecho de toda persona al uso y goce de sus bienes y la facultad de la ley de subordinarlos al interés social. El Pacto fue incorporado a la Constitución mediante el art.~75, inc.~22, y el CCyC lo refleja en su articulado.

La expresión «función social de la propiedad» fue suprimida del art.~15 tal como lo votó el Congreso, aunque figuraba en el proyecto original; se sostiene que la supresión no modifica ni aminora la función social, que rige en el ordenamiento nacional y forma parte expresamente de catorce constituciones provinciales. Al tiempo que limita el ejercicio absoluto del derecho del titular, también lo protege del ejercicio ilimitado o abusivo de los demás.

\begin{importante}
El dominio en Argentina es un derecho conformado e integrado por la función social de la propiedad. No hay contradicción entre ser propietario y tener límites.
\end{importante}

\section{Clasificación de los límites (art.~1970)}

Las limitaciones al dominio privado en interés público se rigen por el derecho administrativo; el uso del dominio sobre inmuebles debe ejercerse conforme a las normas administrativas de cada jurisdicción, y los límites del Capítulo sobre relaciones de vecindad rigen en subsidio.

\begin{table}[H]
\centering
\small
\begin{tabularx}{\textwidth}{>{\raggedright\arraybackslash}p{0.15\textwidth}>{\raggedright\arraybackslash}p{0.19\textwidth}>{\raggedright\arraybackslash}X>{\raggedright\arraybackslash}p{0.19\textwidth}}
\toprule
\textbf{Tipo} & \textbf{Fundamento} & \textbf{Regulación} & \textbf{Indemnización} \\
\midrule
Interés público & Beneficio de la comunidad & Derecho administrativo (ilimitadas en número y clase) & No, en principio \\
Interés de los vecinos & Convivencia y tutela de la vecindad & CCyC, arts.~1970 a 1982 & No, en principio, salvo daño efectivo \\
\bottomrule
\end{tabularx}
\caption{Tipos de restricciones al dominio}
\end{table}

Son ejemplos de interés público las normas de higiene, tranquilidad, horarios de actividades, velocidades máximas y mínimas y prohibición de estacionar.

\section{Restricciones en interés de los vecinos}

No son taxativas. El CCyC regula expresamente: cláusulas de inenajenabilidad (art.~1972), inmisiones inmateriales (art.~1973), camino de sirga (art.~1974), aguas (arts.~1975 a 1977), luces y vistas (arts.~1978 y 1979), árboles, arbustos y plantas (art.~1980) y la obligación de soportar instalaciones provisorias y el paso de personas en obras linderas. Siempre primarán las normas administrativas del lugar: el derecho privado cede ante el orden público.

\subsection{Cláusulas de inenajenabilidad (art.~1972)}

Son nulas si obligan al propietario a no enajenar a título oneroso o a no constituir derechos reales; son válidas solo cuando se obliga a no enajenar a persona determinada. Los donantes pueden prohibir enajenar por un máximo de diez años y los condóminos pueden obligarse a no dividir por un máximo de diez años. Si el dueño no desea enajenar, no enajena, pues el dominio es perpetuo; lo que la ley anula es el compromiso futuro, reflejo de la fortaleza del derecho real y de la seguridad jurídica.

\subsection{Inmisiones inmateriales (art.~1973)}

Las molestias que ocasionan el humo, calor, olores, luminosidad, ruidos, vibraciones o inmisiones similares por actividades en inmuebles vecinos no deben exceder la \textbf{normal tolerancia}, teniendo en cuenta las condiciones del lugar y aunque medie autorización administrativa. Los jueces pueden disponer la cesación o disminución de las molestias hasta lo aceptable, y/o la indemnización de los daños, ponderando especialmente el respeto debido al uso regular de la propiedad, la prioridad en el uso, el interés general y las exigencias de la producción.

\begin{table}[H]
\centering
\small
\begin{tabularx}{\textwidth}{>{\raggedright\arraybackslash}p{0.2\textwidth}>{\raggedright\arraybackslash}X>{\raggedright\arraybackslash}X}
\toprule
\textbf{Criterio} & \textbf{Descripción} & \textbf{Ejemplo} \\
\midrule
Uso regular de la propiedad & Cómo se usa la propiedad que genera la inmisión & Una fundición realiza un uso regular de su propiedad industrial \\
Prioridad en el uso & Quién estaba primero en el lugar & Si la escuela estaba antes, quien compra al lado sabe a qué se expone \\
Interés general & Si la actividad sirve a la comunidad & Una escuela o un hospital tienen un interés público que pesa en la ponderación \\
Exigencias de la producción & Si la actividad no puede realizarse sin la inmisión & Una fundición no puede funcionar sin generar calor \\
\bottomrule
\end{tabularx}
\caption{Criterios de la normal tolerancia}
\end{table}

\begin{ejemplo}{una fundición y una escuela}
Una fundición de hierro generaba calor que se percibía desde el inmueble contiguo. El juez ponderó la prioridad en el uso (la fundición estaba primero), el interés de la producción, las exigencias de la actividad y el uso regular de la propiedad; impuso medidas para amortiguar el calor (aislantes y materiales), pero no ordenó el cese de la actividad. Del mismo modo, quien adquiere un inmueble al lado de una escuela primaria no puede reclamar por el ruido del recreo: pesan la prioridad en el uso y el interés general, y los valores inmobiliarios reflejan la zona.
\end{ejemplo}

\subsection{Camino de sirga (art.~1974)}

El dueño de un inmueble colindante con las orillas de los cauces o situado en las márgenes de los canales debe dejar libre una franja de quince metros, a la que puede acceder para posibilitar su uso. Originalmente era mayor porque se usaba para jalar embarcaciones por la orilla. Es una restricción legal de interés público (facilitar la navegación) que no genera indemnización; el propietario no puede construir, plantar ni realizar obras que afecten la navegación o el tránsito de personas o animales, y quien se sienta perjudicado está legitimado para pedir la remoción. El Estado puede reglamentar el uso pero no construir allí: tendría que expropiar, fundado en una ley especial. Tampoco puede el titular obstaculizar el curso de las aguas ni modificar su velocidad o dirección (salvo obras defensivas), bajo pena de remoción e indemnización.

\subsection{Aguas (art.~1976)}

El titular del fundo inferior debe recibir las aguas, arena y piedras que descienden naturalmente del fundo superior y no puede realizar obras que alteren ese escurrimiento. Si lo hace, el titular del fundo superior puede pedir la remoción e incluso la indemnización.

\subsection{Luces y vistas (arts.~1978 y 1979)}

\begin{table}[H]
\centering
\small
\begin{tabularx}{\textwidth}{>{\raggedright\arraybackslash}p{0.22\textwidth}>{\raggedright\arraybackslash}p{0.2\textwidth}>{\raggedright\arraybackslash}X>{\raggedright\arraybackslash}p{0.1\textwidth}}
\toprule
\textbf{Tipo} & \textbf{Función} & \textbf{Distancias exigidas} & \textbf{Art.} \\
\midrule
Vistas frontales & Permiten ver el inmueble vecino & Mínimo tres metros del lindero & 1978 \\
Vistas laterales u oblicuas & Permiten ver en ángulo & Mínimo sesenta centímetros del lindero & 1978 \\
Luces & Solo permiten entrada de luz y aire, no la vista & Superiores a 1,80 m desde el piso, salvo fijas o no transparentes & 1979 \\
\bottomrule
\end{tabularx}
\caption{Luces y vistas}
\end{table}

\begin{ejemplo}{vista y elevación del vecino}
En principio no puede impedirse que el lindero ejerza su derecho a elevar su construcción aunque prive de vista y luz. Se menciona un caso en la Patagonia en que la elevación de un vecino privó directamente la vista del Nahuel Huapi.
\end{ejemplo}

\subsection{Árboles, arbustos y plantas (art.~1980)}

El Código introduce un criterio subjetivo: ya no se fijan distancias mínimas de plantación, sino que nadie puede tener árboles, arbustos u otras plantas que causen molestias que excedan la normal tolerancia. El vecino perjudicado puede exigir que sean retirados, cortados o podados, y puede cortar él mismo las \textbf{raíces} que penetren en su inmueble, pero \textbf{no} las ramas, que solo puede cortar el propietario o a su pedido. Por ejemplo, las raíces del gomero pueden levantar los pisos y veredas de la casa vecina. Los apuntes mencionan casos de extrema violencia, incluso homicidios, por conflictos de árboles limítrofes.

\section{Regla general final}

\begin{importante}
En todos los casos de límites al dominio primarán las normas administrativas del lugar (alturas de ventanas, luces, distancias de árboles, etc.), por preeminencia del orden público sobre el derecho privado. Los límites no son taxativos: el orden público es dinámico porque son dinámicos los intereses de la comunidad, y los límites pueden modificarse en la medida en que respondan a un interés de orden público.
\end{importante}

\cornell{límites al dominio}{
\cq{¿Por qué el dominio absoluto no es ilimitado?}{Porque su ejercicio no debe ser abusivo y reconoce límites de interés público y vecinal, que se amplían en armonía con los derechos humanos.}
\cq{¿Cómo evolucionó el abuso del derecho?}{Vélez admitía desnaturalizar la cosa; la Ley~17.711 incorporó el abuso del derecho (razones del legislador); el CCyC (art.~10) lo vincula con los principios del sistema.}
\cq{¿Qué lugar ocupa la función social de la propiedad?}{Integra el dominio: limita su ejercicio absoluto y lo protege del abuso de otros; no hay contradicción entre ser propietario y tener límites.}
\cq{Distinga límites de interés público y de vecindad (art.~1970).}{Los primeros se rigen por el derecho administrativo (ilimitados en número y clase); los segundos, por los arts.~1970 a 1982, no taxativos. Ambos, en principio, sin indemnización.}
\cq{¿Cómo se determina la normal tolerancia (art.~1973)?}{Ponderando uso regular de la propiedad, prioridad en el uso, interés general y exigencias de la producción.}
\cq{Compare luces y vistas, y el régimen de raíces y ramas.}{Vistas: 3~m (frontales) y 60~cm (laterales); luces: más de 1,80~m. El vecino puede cortar raíces, no ramas.}
}{El dominio absoluto está sometido a límites de interés público (derecho administrativo) y vecinal (arts.~1970 a 1982), fundados en la función social de la propiedad y en el abuso del derecho. Las normas administrativas primarán siempre y los límites no son taxativos.}

% ============================================================
\part{Condominio y medianería}
% ============================================================

% ------------------------------------------------------------
\chapter[Concepto y caracteres del condominio]{Concepto y caracteres del condominio}\label{cap:condominio}
% ------------------------------------------------------------

El condominio es el primer derecho real en que aparece «el otro»: la pluralidad de titulares sobre una misma cosa. De esa circunstancia nacen sus caracteres, sus facultades restringidas y la abundancia de normas pensadas para el conflicto.

\section{El condominio en el sistema de derechos reales}

Es uno de los catorce derechos reales del art.~1887 y comparte los rasgos comunes: es un derecho real sobre cosa propia, absoluto y de contenido patrimonial, regulado por normas de orden público, oponible \textit{erga omnes} a partir de la publicidad (posesoria o registral), se ejerce por la posesión y se adquiere por título suficiente y modo.

\section{Relevancia práctica}

Es uno de los derechos reales de mayor aplicación cotidiana: aparece en herencias donde varios herederos quedan como titulares de un bien, en parejas que compran un inmueble o un vehículo, en familias que comparten campos y en empresas que cotitularizan activos. Permite asesorar a los condóminos antes del conflicto, redactar convenios de uso y goce, iniciar o defender acciones de partición, resolver conflictos sobre uso exclusivo y valor locativo y proteger a los acreedores que quieran ejecutar una parte indivisa. La vida humana genera condominio por herencia, por afecto o por negocio.

\section{Un derecho autónomo y esencialmente inestable}

El condominio \textbf{no es un dominio «de dos»}: es un derecho autónomo con normas propias. La diferencia fundamental es la \textbf{exclusividad}: en el dominio el titular es único dueño; en el condominio aparece el otro y, con él, la posibilidad del conflicto, por lo que su regulación está plagada de normas para cuando el conflicto surge.

\begin{regla}{Vélez Sársfield}
«El condominio es esencialmente inestable.»
\end{regla}

Las relaciones humanas son inestables y la finitud es propia de los seres humanos. Quienes adquieren en condominio se miran a los ojos al hacerlo, pero con el tiempo pueden no tener el mismo rostro, la misma historia, los mismos afectos ni los mismos intereses, y por la muerte de alguno sus herederos ocuparán su lugar. Para fines de inversión existen figuras más adecuadas, como las sociedades o los fideicomisos.

\begin{ejemplo}{Juan y Carolina}
Juan y Carolina, pareja con un proyecto de vida en común, compraron un departamento en la Ciudad Autónoma de Buenos Aires, donde residen, y luego un automóvil que usan de manera compartida. Este caso se retoma en los capítulos siguientes para analizar los conflictos posibles.
\end{ejemplo}

\section{Concepto (art.~1983)}

\begin{definicion}{condominio (art.~1983 CCyC)}
Condominio es el derecho real de propiedad sobre una cosa que pertenece en común a varias personas y que corresponde a cada una por una parte indivisa.
\end{definicion}

Las partes de los condóminos se \textbf{presumen iguales}, salvo que la ley o el título dispongan una proporción diferente: si uno aportó el 70\,\% y el otro el 30\,\%, ese será el porcentaje de cada uno; si el título nada dice, 50\,\% y 50\,\%.

\section{Caracteres}

\begin{table}[H]
\centering
\small
\begin{tabularx}{\textwidth}{>{\raggedright\arraybackslash}p{0.2\textwidth}>{\raggedright\arraybackslash}X>{\raggedright\arraybackslash}X}
\toprule
\textbf{Carácter} & \textbf{Descripción} & \textbf{Ejemplo (Juan y Carolina)} \\
\midrule
Unidad de objeto & Recae sobre una sola cosa, mueble o inmueble & El auto o el departamento es una sola cosa \\
Pluralidad de sujetos & Dos o más titulares & Juan y Carolina (o sus herederos) \\
Partes ideales (indivisas) & Ninguno tiene una parte material determinada & Ninguno es dueño de las ruedas delanteras ni de las traseras \\
\bottomrule
\end{tabularx}
\caption{Caracteres del condominio}
\end{table}

Si los condóminos no quisieran continuar juntos, no podría resolverse que cada uno se lleve dos ruedas: el derecho no recae sobre las ruedas, ni sobre el baño o la cocina del departamento, sino sobre una \textbf{parte indivisa}. Cada uno puede decir que es su auto, pero lo es con otro.

\begin{ejemplo}{fallecimiento de un condómino}
Si Carolina fallece y deja tres hijos, el condominio pasa a ser de cuatro personas: Juan, que conserva su mitad, y los tres herederos de Carolina, que comparten la otra mitad en partes iguales por sucesión \textit{mortis causa}.
\end{ejemplo}

\cornell{concepto y caracteres del condominio}{
\cq{¿Por qué el condominio no es un «dominio de dos»?}{Porque es un derecho autónomo: falta la exclusividad. El otro condómino introduce el conflicto, por lo que la regulación se orienta a resolverlo.}
\cq{¿Qué significa que sea «esencialmente inestable»?}{Que las relaciones humanas y la vida de los titulares cambian (conflictos, muerte, herederos), lo que pone en riesgo la comunidad.}
\cq{Defina el condominio (art.~1983) y la presunción de partes.}{Derecho real de propiedad sobre una cosa común a varias personas, por parte indivisa. Las partes se presumen iguales salvo proporción distinta en la ley o el título.}
\cq{¿Cuáles son sus caracteres?}{Unidad de objeto, pluralidad de sujetos y partes ideales o indivisas, no materiales.}
\cq{¿Qué ocurre si fallece un condómino?}{Sus herederos ocupan su lugar y comparten su parte indivisa.}
}{El condominio es un derecho real autónomo, de unidad de objeto, pluralidad de sujetos y partes indivisas presumidas iguales, esencialmente inestable por la presencia del otro titular y por la finitud humana.}

% ------------------------------------------------------------
\chapter[Facultades, gastos y administración del condominio]{Facultades, gastos y administración del condominio}\label{cap:facultades-cond}
% ------------------------------------------------------------

Las facultades de los condóminos dependen de aquello sobre lo que recaen: son \textbf{amplísimas sobre la parte indivisa} y \textbf{restringidas sobre la cosa en común}, precisamente porque está el otro.

\section{Facultades sobre la parte indivisa (art.~1989)}

\begin{regla}{art.~1989 CCyC}
El condómino puede enajenar y gravar su parte indivisa sin el consentimiento de los restantes condóminos. Los acreedores pueden embargarla y ejecutarla sin esperar el resultado de la partición, que les es inoponible.
\end{regla}

El condómino puede \textbf{enajenar} (vender, donar o transferir), \textbf{gravar} (por ejemplo, hipotecar) su parte, y sus acreedores pueden embargarla y llevarla a subasta.

\section{Facultades sobre la cosa en común}

El condómino puede \textbf{usar y gozar} la cosa, pero de forma \textbf{conjunta} con los demás y \textbf{sin alterar su destino}; puede usarla en forma exclusiva siempre que no excluya al otro ni cambie el destino.

\subsection{Convenio de uso y goce exclusivo (art.~1987)}

\begin{regla}{art.~1987 CCyC}
Los condóminos pueden convenir el uso y goce alternado de la cosa común o que se ejercite de manera exclusiva y excluyente sobre determinadas partes materiales.
\end{regla}

Es una de las novedades más importantes del Código de 2014: permite pacificar una cuestión interna que a veces provocaba conflictos serios, sobre todo tras el fallecimiento de alguno (personas que durante veinte o treinta años convivieron sin dificultad y, al morir uno, el otro se dice tan dueño de la casa de atrás como de la de adelante).

\begin{ejemplo}{dos casas sobre un terreno}
Sobre un terreno en condominio hay una casa adelante y otra atrás. Registralmente es un solo inmueble, pero los condóminos pueden convenir que uno use exclusivamente la de adelante y el otro la de atrás. Sin convenio, ambos son titulares de una parte indivisa de todo: un poco del terreno, un poco de cada casa.
\end{ejemplo}

\subsection{Uso exclusivo no pactado y valor locativo (art.~1988)}

\begin{regla}{art.~1988 CCyC}
El condómino que usa la cosa común en forma exclusiva está obligado, excepto pacto en contrario, a satisfacer una indemnización, desde que el otro condómino lo solicita.
\end{regla}

Si Juan y Carolina se separan, ella se va y Juan queda viviendo solo en el departamento, ¿puede ella reclamar alquileres por todos esos años? El \textbf{valor locativo} es una compensación como si se pagara un alquiler, pero \textbf{no es un alquiler}: el condómino usa la cosa ejerciendo su facultad de usar y gozar, y si nada se le pide se presume que lo hace con el consentimiento del otro. Solo procede \textbf{desde que el otro condómino lo notificó fehacientemente}: si Carolina notificó a los tres días, procede desde entonces; si notificó dos o tres años después, solo desde esa fecha. No puede reclamar compensaciones anteriores a la notificación.

\section{Disposición, gastos y deudas}

Sobre la cosa en su totalidad, ningún condómino puede disponer materialmente de una parte determinada sin autorización de los demás, ni realizar actos de disposición jurídica, como alquilarla, sin el consentimiento de la mayoría. Si son tres condóminos y dos dicen que sí y uno que no, este no puede oponerse.

Los \textbf{gastos comunes} se afrontan en proporción al porcentaje de participación (si Juan tiene el 70\,\% pagará el 70\,\%; si el título nada dice, por mitades). Si un condómino contrata con un tercero (por ejemplo, un servicio de pintura), él es el obligado ante el tercero, pero puede recuperar internamente de los demás su parte; esto genera debates sobre si los gastos fueron de conservación o de mejora, necesarios o no.

\section{Administración y asamblea (arts.~1993 y 1994)}

\begin{regla}{art.~1994 CCyC}
Cualquier condómino puede convocar a una asamblea de condóminos, citando a los demás en forma fehaciente, indicando la finalidad y con antelación suficiente. La resolución de la mayoría absoluta de los condóminos computada según el valor de las partes indivisas es obligatoria para todos.
\end{regla}

Presenta dificultades prácticas en condominios de solo dos titulares, donde no resulta claro cómo convocar ni dónde alcanzar acuerdo; si hay empate, el art.~1994 establece que \textbf{decide la suerte}. El régimen es útil en condominios con varios titulares, pero no es de aplicación amplia.

\cornell{facultades y administración del condominio}{
\cq{¿Qué puede hacer el condómino con su parte indivisa (art.~1989)?}{Enajenarla y gravarla sin consentimiento de los demás; sus acreedores pueden embargarla y ejecutarla sin esperar la partición.}
\cq{¿Y sobre la cosa en común?}{Usar y gozar conjuntamente y sin alterar el destino; no disponer materialmente de una parte ni alquilar sin consentimiento de la mayoría.}
\cq{¿Qué permite el art.~1987?}{Convenir uso alternado o exclusivo y excluyente de partes materiales, pacificando conflictos internos.}
\cq{¿Cuándo procede el valor locativo (art.~1988) y qué naturaleza tiene?}{Desde la notificación fehaciente del condómino que lo reclama; no es un alquiler, sino una compensación por uso exclusivo. Antes se presume consentimiento.}
\cq{¿Cómo se reparten los gastos y cómo decide la asamblea (art.~1994)?}{Gastos en proporción a la parte; la asamblea decide por mayoría absoluta según el valor de las partes, y en caso de empate, la suerte.}
}{Las facultades del condómino son amplísimas sobre su parte indivisa y restringidas sobre la cosa común. El uso y goce puede pactarse; el uso exclusivo no pactado da lugar a valor locativo desde la notificación; los gastos se reparten según las partes y la asamblea decide por mayoría computada por valor.}

% ------------------------------------------------------------
\chapter[Partición e indivisión en el condominio]{Partición e indivisión en el condominio}\label{cap:particion}
% ------------------------------------------------------------

\section{Clasificación}

\begin{table}[H]
\centering
\small
\begin{tabularx}{\textwidth}{>{\raggedright\arraybackslash}p{0.18\textwidth}>{\raggedright\arraybackslash}X>{\raggedright\arraybackslash}X>{\raggedright\arraybackslash}X}
\toprule
\textbf{Tipo} & \textbf{Definición} & \textbf{Consecuencia} & \textbf{Plazo máximo} \\
\midrule
Sin indivisión forzosa & No hay causa legal que impida dividir & Cualquier condómino puede pedir la partición en cualquier momento & Sin plazo: acción imprescriptible \\
Con indivisión forzosa & Hay una causa legal que restringe la partición & Partición temporalmente impedida según la fuente & Diez años (convención y testamento); cinco años (sentencia) \\
\bottomrule
\end{tabularx}
\caption{Tipos de condominio según la partición}
\end{table}

La indivisión forzosa surge únicamente de la ley, aunque con distintas fuentes. Se distingue la \textbf{indivisión temporaria}, de plazo determinado, y la \textbf{perdurable}, vinculada a la naturaleza del objeto, como en la medianería (capítulo~\ref{cap:medianeria}).

\section{Condominio sin indivisión forzosa}

\begin{regla}{arts.~1996 y 1999 CCyC}
Todo condómino puede, en cualquier tiempo, pedir la partición de la cosa. La acción es imprescriptible (art.~1996). El derecho a pedir la partición es irrenunciable por tiempo indeterminado (art.~1999).
\end{regla}

Es imprescriptible, al igual que el dominio es perpetuo, e irrenunciable por tiempo indeterminado: no puede pactarse definitivamente que nunca se pedirá.

\begin{ejemplo}{la cláusula prohibitiva de la partición}
Una pareja que compra «el departamento de sus sueños» pide incluir una condición: que nunca puedan dividirlo. Quien asesora no puede garantizarlo, porque el derecho a pedir la partición es irrenunciable por tiempo indeterminado y, si Carolina muere, es dudoso para qué querrían sus herederos tener el departamento con Juan.
\end{ejemplo}

\section{Condominio con indivisión forzosa temporaria: fuentes}

\begin{table}[H]
\centering
\small
\begin{tabularx}{\textwidth}{>{\raggedright\arraybackslash}p{0.18\textwidth}>{\raggedright\arraybackslash}p{0.12\textwidth}>{\raggedright\arraybackslash}X>{\raggedright\arraybackslash}p{0.2\textwidth}}
\toprule
\textbf{Fuente} & \textbf{Art.} & \textbf{Regla} & \textbf{Plazo} \\
\midrule
Convención & 2000 & Los condóminos pueden convenir suspender la partición; si no fijan plazo, o es incierto o superior, se considera celebrada por diez años & Hasta diez años, renovable por igual término \\
Sentencia judicial & 2001 & A petición de parte, cuando la partición es nociva para el interés de los condóminos o de alguno, el juez puede disponerla con carácter diferido & Hasta cinco años, renovable judicialmente por otros cinco \\
Acto de última voluntad & 2330 & El testador puede imponer a los herederos, aun mayores de edad, la indivisión de la herencia o de algún bien & Hasta diez años \\
\bottomrule
\end{tabularx}
\caption{Fuentes de la indivisión temporaria}
\end{table}

\begin{itemize}
    \item \textbf{Convención.} Si Carolina muere, sus herederos deben cumplir el pacto, que también los obliga.
    \item \textbf{Sentencia.} Puede \textbf{demorar} una partición que no tenía restricción, si se demuestra perjudicial para todos, o \textbf{anticipar} una partición restringida, a pedido de quienes acrediten circunstancias graves.
    \item \textbf{Última voluntad.} Es muy frecuente en la práctica («les dejo todo a mis hijos, pero que nunca lo separen»), pero el plazo máximo es de diez años y no puede ser indefinido.
\end{itemize}

\section{Partición nociva y derechos humanos}

La partición nociva existía ya en el Código derogado. El ejemplo clásico es el campo sembrado: si uno pide la partición en plena cosecha, el otro puede oponerse judicialmente porque sería antieconómica y perjudicial para todos.

\begin{importante}
Ya existen fallos que demoran la partición con fundamento en los derechos humanos, los derechos del niño, de los ancianos y de los enfermos: algunos jueces la demoran cinco años y, tras nueva evaluación, otros cinco, ya no por perjuicio económico común sino por el interés del condómino anciano o enfermo que vive en el inmueble.
\end{importante}

Esto no significa que la partición no pueda pedirse: es un derecho rotundo. Ante un caso concreto y cierta prueba, el juez puede demorarla o anticiparla.

\cornell{partición e indivisión}{
\cq{¿Qué caracteres tiene el derecho a pedir la partición (arts.~1996 y 1999)?}{Es imprescriptible e irrenunciable por tiempo indeterminado.}
\cq{¿Puede pactarse que nunca se divida?}{No: solo se puede suspender la partición por hasta diez años (art.~2000), renovables.}
\cq{Enumere las fuentes de indivisión temporaria con sus plazos.}{Convención (hasta diez años, art.~2000), sentencia (hasta cinco años renovables por otros cinco, art.~2001) y testamento (hasta diez años, art.~2330).}
\cq{¿Qué es una partición nociva?}{La que perjudica el interés de los condóminos o de alguno; ejemplo: pedirla en plena cosecha.}
\cq{¿Qué rol cumplen hoy los derechos humanos?}{Fundan fallos que demoran la partición en interés de condóminos ancianos, enfermos o niños, sin negar el derecho a pedirla.}
}{La partición puede pedirse siempre salvo indivisión forzosa, que solo la ley establece, con fuentes convencional, judicial o testamentaria de plazo máximo diez o cinco años. La partición nociva y los derechos humanos permiten demorarla o anticiparla según el caso.}

% ------------------------------------------------------------
\chapter[Medianería: estructura, fundamentos y clasificación de los muros]{Medianería: estructura, fundamentos y clasificación de los muros}\label{cap:medianeria}
% ------------------------------------------------------------

La medianería es la regulación legal de muros, cercos y fosos dentro del condominio con indivisión forzosa de origen legal. Para quien ejerce en la práctica inmobiliaria, notarial o registral, comprender este régimen permite asesorar a propietarios sobre derechos y obligaciones en construcciones colindantes, resolver conflictos sobre muros y cercos, redactar escrituras que reflejen la verdadera naturaleza del inmueble, orientar trámites registrales y evitar litigios costosos. Todo acto de compraventa, construcción o modificación de un inmueble toca este régimen. Cada muro es una especie de acuerdo silencioso entre vecinos: ambos necesitan encerrarse (derecho) y ambos deben contribuir (obligación). El estudio avanza de lo más simple (qué tipos de muros hay) a lo más complejo (quién paga qué y cuándo).

\section{Estructura del condominio con indivisión forzosa}

Es un condominio de \textbf{origen legal} con dos facetas, reguladas en el Título Cuarto del Libro Cuarto.

\begin{definicion}{condominio sobre accesorios indispensables (arts.~2004 y 2005)}
Condominio sobre accesorios indispensables para el uso común de dos inmuebles colindantes, como pasillos de acceso y pasadizos comunes. Cada condómino puede usarlos conforme a su destino, sin afectar al otro, y no puede pedir la partición.
\end{definicion}

\begin{definicion}{condominio sobre muros, cercos y fosos (arts.~2006 a 2030)}
Condominio sobre muros de cerramiento entre inmuebles colindantes. Es la faceta compleja, por su vinculación con la medianería.
\end{definicion}

\section{Un límite al derecho de propiedad}

La medianería otorga a cada propietario de inmuebles colindantes urbanos el \textbf{derecho y la obligación recíprocos} de construir un muro de cerramiento. En este régimen se advierte la potencia del derecho real: el muro \textbf{no requiere la conformidad del vecino}; quien construye puede hacerlo aunque el vecino se oponga o no contribuya, y la ley autoriza a edificar \textbf{la mitad del muro sobre el terreno del vecino}. Quien construye adelanta el cumplimiento de una obligación legal que también pesaba sobre el vecino y por eso tiene derecho a reclamar. Con las normativas actuales, los muros de cerramiento miden habitualmente \textbf{tres metros de alto por treinta centímetros de espesor} («pared de treinta»).

\section{Fundamentos}

\begin{itemize}
    \item \textbf{Constitucionales.} CN, art.~18 (inviolabilidad del domicilio) y art.~19 («Las acciones privadas de los hombres que de ningún modo ofendan al orden y a la moral pública, ni perjudiquen a un tercero, están solo reservadas a Dios, y exentas de la autoridad de los magistrados»): protegen la privacidad e intimidad personal y familiar, que dependen de la posibilidad de encerrarse.
    \item \textbf{Facultad de exclusión y encerramiento (art.~1944).} Si cada propietario tiene esa facultad, su vecino también; la faceta obligacional surge de la normativa del condominio con indivisión forzosa.
    \item \textbf{Económicos.} Permite economizar el suelo, escaso y no siempre cultivable; el ahorro se multiplica por la enorme cantidad de inmuebles urbanos.
    \item \textbf{Técnicos.} Solidez: los inmuebles se «abrazan», se apoyan unos en otros y se fortalecen mutuamente.
\end{itemize}

\section{Clasificación de los muros}

Siguiendo la estructura didáctica de Marina Mariani de Vidal, se formulan dos preguntas: \textbf{dónde está} el muro (criterio físico) y \textbf{de quién es} (criterio jurídico).

\subsection{Clasificación física}

\begin{table}[H]
\centering
\small
\begin{tabularx}{\textwidth}{>{\raggedright\arraybackslash}p{0.2\textwidth}>{\raggedright\arraybackslash}p{0.12\textwidth}>{\raggedright\arraybackslash}X}
\toprule
\textbf{Muro} & \textbf{Art.~2006} & \textbf{Definición} \\
\midrule
Lindero, separativo o divisorio & inc.~a & Sinónimos: demarca el inmueble y lo delimita del colindante. \\
Incaballado & inc.~b & Lindero que se asienta parcialmente en cada inmueble; lo típico es quince centímetros en cada terreno («pared de 30»: 15 y 15). \\
Contiguo & inc.~c & Lindero que se asienta totalmente en uno de los inmuebles, de modo que el filo coincide con el límite separativo. \\
De elevación & inc.~d & Lindero que excede la altura del muro de cerramiento ordinario (tres metros), sea incaballado o contiguo, en la altura que supera esa medida. \\
Enterrado & inc.~h & Ubicado debajo del nivel del suelo, sin servir de cimiento a una construcción en superficie. \\
\bottomrule
\end{tabularx}
\caption{Clasificación física de los muros}
\end{table}

Por ejemplo, un muro de cinco metros es medianero hasta los tres y de elevación en los dos metros que superan esa altura.

\subsection{Clasificación jurídica: el muro medianero}

\begin{definicion}{muro medianero}
Es el lindero común, perteneciente en condominio a ambos colindantes. Surge desde la construcción.
\end{definicion}

Desde que uno o ambos vecinos construyen el muro, sea incaballado o contiguo, el muro \textbf{pertenece a los dos}, aunque no haya sido construido en el terreno de ambos; por eso la ley permite que Juan lo construya avanzando quince centímetros sobre el terreno de su vecina María. El condominio es de \textbf{fuente legal}: nace medianero desde la construcción, hasta los tres metros de altura.

\begin{table}[H]
\centering
\small
\begin{tabularx}{\textwidth}{>{\raggedright\arraybackslash}p{0.28\textwidth}>{\raggedright\arraybackslash}X}
\toprule
\textbf{Característica} & \textbf{Descripción} \\
\midrule
Origen legal & El condominio surge por ley desde la construcción. \\
Titularidad & Pertenece a ambos condóminos por igual (50\,\% y 50\,\%). \\
Derechos del no constructor & Usar el muro, anclar construcciones, reclamar contribución. \\
\bottomrule
\end{tabularx}
\caption{Características del muro medianero}
\end{table}

El muro puede ser construido por ambos vecinos, por uno solo, con comunidad de gastos o pagado por uno solo. % TODO: contenido incompleto o ambiguo en las notas originales (la transcripción del pasaje sobre la construcción compartida o unilateral es confusa).
Si lo pagó uno solo, el vecino que no lo construyó es igualmente titular, porque el muro es medianero desde la construcción. Como el muro construido por Juan pertenece a María en un 50\,\%, ella, aunque no haya contribuido, puede \textbf{anclar} construcciones al muro, \textbf{elevar} construcciones por sobre los tres metros, \textbf{prolongar} cimientos hacia abajo y \textbf{destruir y reconstruir} el muro si es necesario para su obra.

\cornell{medianería: estructura, fundamentos y muros}{
\cq{¿Cuáles son las dos facetas del condominio con indivisión forzosa?}{Accesorios indispensables (arts.~2004 y 2005, como pasillos) y muros, cercos y fosos (arts.~2006 a 2030).}
\cq{¿Por qué la medianería muestra la potencia del derecho real?}{Porque el muro puede construirse sin conformidad del vecino, incluso sobre la mitad de su terreno, al adelantar una obligación legal recíproca.}
\cq{¿Cuáles son sus fundamentos?}{Constitucionales (CN arts.~18 y 19: privacidad e intimidad), facultad de encerramiento (art.~1944), económicos (ahorro de suelo) y técnicos (solidez).}
\cq{Distinga muro incaballado, contiguo, de elevación y enterrado.}{Incaballado: asentado parcialmente en cada inmueble. Contiguo: totalmente en uno. De elevación: la altura que excede tres metros. Enterrado: bajo el nivel del suelo sin servir de cimiento.}
\cq{¿De quién es el muro y desde cuándo?}{Es medianero, condominio legal por mitades, desde que se construye (hasta tres metros); el vecino que no pagó es igualmente titular con derechos plenos.}
}{La medianería impone a los propietarios colindantes el derecho y la obligación recíprocos de cerramiento, fundada en la privacidad, la economía de suelo y la solidez. Los muros se clasifican por su ubicación y por su titularidad; el muro nace medianero desde su construcción.}

% ------------------------------------------------------------
\chapter[Régimen del muro medianero, presunciones y medianería rural]{Régimen del muro medianero, presunciones y medianería rural}\label{cap:regimen-medianero}
% ------------------------------------------------------------

\section{Acción de cobro de medianería}

\begin{definicion}{acción de medianería (art.~2019)}
Desde la construcción del muro nace el derecho a reclamar el cobro de la medianería, haya utilizado o no el vecino el muro.
\end{definicion}

La mora es \textbf{automática} (art.~886): no requiere interpelación ni demanda previa. La acción prescribe a los \textbf{cinco años} desde la construcción (art.~2560); pasado ese plazo se extingue el derecho a reclamar, incluso si el vecino comienza a usar el muro después. La defensa es la prescripción liberatoria.

\begin{ejemplo}{terreno baldío y uso posterior}
Se construye y paga un muro; no vive nadie al lado, pasa el tiempo y nadie lo utiliza. A los seis años, el propietario comienza a construir y usa el muro: la acción de cobro ya prescribió. Puede usar los cimientos necesarios para tres metros y la pared de treinta, porque el muro es suyo en condominio.
\end{ejemplo}

\section{Muro de elevación y muro enterrado}

En un mismo muro conviven dos regímenes: hasta tres metros, \textbf{medianero}; desde tres metros hacia arriba, \textbf{privativo} del constructor. Para la parte que excede tres metros, la obligación nace cuando el vecino efectivamente la usa, apoyando en el segmento elevado. El muro enterrado sigue un régimen similar: la obligación de pagar nace desde la utilización específica (art.~2019).

\section{Imposibilidad de liberarse}

No es posible liberarse de la obligación de pagar, ni de los gastos de conservación o reconstrucción, mediante abdicación o declinación de la parte. Como el muro ya nació medianero, para liberarse habría que transmitir el 50\,\% al vecino, lo que requiere nuevos planos y escritura pública: un procedimiento costoso.

\section{Presunciones legales (arts.~2010 a 2013)}

Cuando hay construcciones antiguas y es difícil determinar la naturaleza del muro, la ley establece presunciones que admiten prueba en contrario, pero son fuertes.

\begin{definicion}{presunción de medianería (art.~2010)}
A menos que se pruebe lo contrario, un muro lindero entre dos edificios que supera los tres metros se presume medianero hasta esa altura, y la parte que los supera se presume privativa del dueño del edificio más alto.
\end{definicion}

No se aplican cuando el muro separa patios, huertos o jardines (y no un edificio), ni cuando hay documentos o signos que identifican la verdadera naturaleza o permiten destruir la presunción.

\section{Medianería rural}

En predios rurales la división no se hace con muros sino con \textbf{cercas vivas} (árboles alineados), \textbf{zanjas}, cercos de alambre u otras estructuras (arts.~2032 y 2033); el titular tiene la facultad de levantar o excavar el cerramiento, como se ve al costado de una ruta. Los \textbf{árboles y arbustos} plantados en la línea divisoria pueden ser medianeros, contiguos o incaballados, según su ubicación respecto de los muros, cercos y fosos linderos, tanto en predios rurales como urbanos (arts.~2034 a 2036). Si causan perjuicio a un condómino, puede exigir su remoción en cualquier tiempo, e incluso su reemplazo. % TODO: contenido incompleto o ambiguo en las notas originales (el apunte no precisa el alcance del reemplazo).

\section{Síntesis}

\begin{conceptoclave}{medianería}
La medianería es una institución del derecho real que limita el derecho de propiedad en favor de la vida urbana moderna. Nace por ley desde la construcción de un muro entre inmuebles colindantes. Sus fundamentos son constitucionales (arts.~18 y 19), económicos (ahorro de terreno) y técnicos (solidez). Ambos propietarios son condóminos con derechos y obligaciones recíprocos, con independencia de quién financió la obra. La ley presume medianería hasta tres metros y prevé regímenes especiales para los muros de elevación y enterrados; las zonas rurales tienen regímenes propios. La acción de cobro prescribe a los cinco años.
\end{conceptoclave}

Para resolver un caso conviene distinguir: (1) dónde está el muro (ubicación física); (2) de quién es (naturaleza jurídica); y (3) qué régimen le corresponde según su tipo.

\cornell{régimen del muro medianero}{
\cq{¿Cuándo nace la acción de cobro y cuándo prescribe?}{Nace desde la construcción, aunque el vecino no use el muro (art.~2019); la mora es automática y prescribe a los cinco años, sin renacer por uso posterior.}
\cq{¿Qué régimen tiene el muro de elevación?}{Medianero hasta tres metros y privativo del constructor por encima; la obligación sobre el excedente nace cuando el vecino lo usa.}
\cq{¿Puede un vecino liberarse declinando su parte?}{No: debería transmitir su 50\,\% por escritura pública y con nuevos planos.}
\cq{¿Qué presume el art.~2010 y cuándo no se aplica?}{Medianería hasta tres metros y privatividad del dueño del edificio más alto por encima. No se aplica entre patios, huertos o jardines ni con documentos o signos en contrario.}
\cq{¿Cómo se divide un predio rural?}{Con cercas vivas, zanjas o alambrados (arts.~2032 y 2033); los árboles divisorios pueden ser medianeros (arts.~2034 a 2036) y su remoción puede exigirse si perjudican.}
}{El cobro de la medianería nace con la construcción, con mora automática y prescripción a los cinco años. La ley presume medianería hasta tres metros, con regímenes diferenciados para elevación y muros enterrados. No es posible liberarse de las obligaciones sin transmitir la parte, y la medianería rural se ordena con cercas vivas, zanjas y alambrados.}

% ============================================================
\part{Propiedad horizontal}
% ============================================================

% ------------------------------------------------------------
\chapter[Concepto y naturaleza de la propiedad horizontal]{Concepto y naturaleza de la propiedad horizontal}\label{cap:ph}
% ------------------------------------------------------------

La propiedad horizontal, regulada en los artículos 2037 a 2069, es uno de los regímenes más presentes en la práctica profesional: la mayoría de la población urbana vive o trabaja en inmuebles sometidos a él. Conocerlo permite asesorar en la compraventa de unidades funcionales (verificando deudas de expensas y reglamentos), redactar o revisar reglamentos y contratos de administración, iniciar juicios ejecutivos por expensas, defender o representar consorcistas en conflictos de asamblea, asesorar sobre prehorizontalidad y articularlo con el usufructo, el embargo, la subasta y la protección de la vivienda.

\section{Del dominio y el condominio a la propiedad horizontal}

El \textbf{dominio} tiene un titular exclusivo y le otorga un poder directo e inmediato sobre la cosa, con las facultades de usar, gozar y disponer (\textit{ius utendi, ius fruendi, ius abutendi}); sobre un inmueble, el dominio se extiende a todo lo plantado y construido.

\begin{ejemplo}{Juan y sus hijos}
Juan adquiere un terreno, lo construye y luego edifica sobre esa casa una planta alta con dos departamentos, con la ilusión de que cada uno sea usado por un hijo. Con el tiempo, decide enajenarles los departamentos. Materialmente se observan tres departamentos, pero jurídicamente existe \textbf{un solo inmueble}, el de Juan; si sus acreedores lo ejecutan, ejecutarán el inmueble con todo lo plantado y construido. Juan puede disponer de toda la cosa o de una parte indivisa: si dona a cada hijo el 25\,\%, nace un \textbf{condominio}, con unidad de objeto, pluralidad de sujetos y partes indivisas, sin que se haya transmitido un departamento determinado a cada hijo.
\end{ejemplo}

En el ámbito de los derechos reales, \textbf{enajenar} es transmitir el derecho de dominio, sea por compraventa, donación o dación en pago.

\begin{table}[H]
\centering
\small
\begin{tabularx}{\textwidth}{>{\raggedright\arraybackslash}X>{\raggedright\arraybackslash}X}
\toprule
\textbf{Dominio} & \textbf{Condominio} \\
\midrule
Titular único & Pluralidad de sujetos \\
Poder directo e inmediato & Partes indivisas o ideales \\
Se extiende a lo construido y plantado & No se corresponde con una parte material determinada \\
Un solo inmueble jurídico & Un solo inmueble jurídico (misma unidad de objeto) \\
\bottomrule
\end{tabularx}
\caption{Dominio y condominio}
\end{table}

\section{El derecho real de propiedad horizontal}

La cuestión es cómo cada departamento puede pertenecer en propiedad a cada titular. La propiedad horizontal permite constituir \textbf{unidades independientes dentro de un mismo inmueble}: un \textbf{dominio exclusivo sobre partes privativas} inescindiblemente unido a un \textbf{condominio sobre las partes comunes}. No es simplemente dominio más condominio: es un derecho único de estructura dual.

\begin{definicion}{propiedad horizontal (art.~2037 CCyC)}
La propiedad horizontal es el derecho real que se ejerce sobre un inmueble propio que otorga a su titular facultades de uso, goce y disposición material y jurídica que se ejercen sobre partes privativas y sobre partes comunes de un edificio, de conformidad con lo que establece este Título y el respectivo reglamento de propiedad horizontal.
\end{definicion}

Un mismo inmueble puede ser objeto de dominio, de condominio o estar afectado a propiedad horizontal.

\subsection{Origen}

Fue creada por la \textbf{Ley~13.512} en 1948 y cumplió una gran función social. Antes, el Código de Vélez Sársfield prohibía específicamente la división horizontal de los inmuebles. Su contenido fue incorporado y actualizado en los arts.~2037 a 2069 del CCyC (Ley~26.994, vigente desde agosto de 2015).

\subsection{Inmuebles afectables y trascendencia}

No solo pueden afectarse propiedades en altura, sino también las de una sola planta: con frente y fondo, «en tira» y dúplex. El sistema es, según los apuntes, una de las mayores creaciones en materia de derechos reales del siglo XX y superó largamente sus objetivos originales; a 70 años de su creación % TODO: la cifra de años figura así en las notas originales.
se demuestra su gran trascendencia social. Prueba de ello es que los \textbf{conjuntos inmobiliarios} son una regulación especial de la propiedad horizontal, por lo que se los llama \textbf{propiedades horizontales especiales} (capítulo~\ref{cap:conjuntos}).

\cornell{concepto de la propiedad horizontal}{
\cq{¿Qué es la propiedad horizontal?}{Derecho real autónomo (arts.~2037 a 2069): dominio exclusivo sobre partes privativas inescindiblemente unido a condominio sobre partes comunes.}
\cq{¿En qué se diferencia del dominio y del condominio?}{El dominio tiene un titular y recae sobre todo el inmueble; el condominio, pluralidad de sujetos con partes indivisas; la propiedad horizontal permite dominio exclusivo sobre cada unidad.}
\cq{¿Por qué el condominio no permite que cada hijo tenga su departamento?}{Porque las partes son indivisas e ideales: no corresponden a una parte material determinada y sigue habiendo un solo inmueble.}
\cq{¿De dónde proviene y qué inmuebles puede abarcar?}{De la Ley~13.512 (1948), incorporada al CCyC; vale para edificios en altura y también para propiedades de una planta, en tira y dúplex.}
\cq{¿Qué son las propiedades horizontales especiales?}{Los conjuntos inmobiliarios, regulación especial de la propiedad horizontal.}
}{La propiedad horizontal articula un dominio exclusivo sobre la unidad con un condominio inescindible sobre lo común. Creada por la Ley~13.512 y recogida por el CCyC, permite unidades independientes dentro de un mismo inmueble.}

% ------------------------------------------------------------
\chapter[Afectación, unidad funcional y partes comunes]{Afectación, unidad funcional y partes comunes}\label{cap:afectacion}
% ------------------------------------------------------------

\section{Requisitos para afectar un inmueble}

El inmueble debe ser material y jurídicamente apto. Cada unidad funcional debe tener salida independiente a la vía pública, acceso común de todos los propietarios a los servicios de agua y tanques, y cumplir las normativas locales de construcción. En inmuebles ya construidos pueden requerirse obras, con costos adicionales y, a veces, imposibilidad material por falta de espacio.

\begin{table}[H]
\centering
\small
\begin{tabularx}{\textwidth}{>{\raggedright\arraybackslash}X>{\raggedright\arraybackslash}X}
\toprule
\textbf{Requisitos administrativos} & \textbf{Requisitos jurídicos} \\
\midrule
Plano de construcción habitual & Confección del reglamento de copropiedad \\
Plano de mensura (esencial) & Escritura pública del reglamento \\
Aprobación por la municipalidad local & Inscripción en el Registro de la Propiedad Inmueble (constitutiva) \\
\bottomrule
\end{tabularx}
\caption{Requisitos para someter un inmueble a propiedad horizontal}
\end{table}

El \textbf{plano de mensura} es esencial: sin él no existiría la propiedad horizontal. Divide e individualiza los sectores privativos y comunes y fija el porcentual que cada parte privativa tiene sobre la cosa común. La \textbf{inscripción} del reglamento es \textbf{constitutiva}: hace nacer el estado de propiedad horizontal.

\section{La unidad funcional}

El objeto del derecho es la \textbf{unidad funcional}, sobre la que recae el dominio exclusivo (parte privativa), inescindiblemente unida a la parte común. Pueden ser departamentos, locales o cocheras, siempre con salida independiente a la vía pública, directa o por un pasaje común.

\section{Partes comunes (arts.~2040 y 2041)}

Son las destinadas a permitir el acceso a las partes privativas (como el ascensor) y a garantizar el funcionamiento del consorcio. Ningún propietario tiene derecho exclusivo sobre ellas y siempre se determinan en el reglamento. Algunas pueden estar destinadas al uso exclusivo de un propietario, como una azotea que es común pero está destinada al uso exclusivo de una unidad.

\begin{regla}{art.~2041 CCyC, cosas y partes necesariamente comunes}
No pueden recibir otro destino en el reglamento, bajo pena de nulidad: el terreno (siempre común); los muros exteriores y los que separan cada unidad funcional; el piso de cada unidad (techo del vecino inferior); el techo de cada unidad (piso del vecino superior).
\end{regla}

\section{La teoría del cubo de aire}

\begin{ejemplo}{el octavo piso}
Quien vive en el octavo piso tiene un piso que es común (techo del vecino del séptimo) y un techo que es común (piso del vecino del noveno). ¿Cuál es entonces el objeto de su derecho?
\end{ejemplo}

El objeto es el \textbf{cubo de aire}: el volumen interior que el plano de mensura identificó como su unidad funcional (en el ejemplo, la unidad «8.º C»), que la arquitectura y la ingeniería hicieron materialmente posible. El Código lo denomina «volumen». El titular es dueño, además, de un \textbf{porcentual indiviso} del terreno inescindiblemente unido a la unidad; de otro modo no tendría ni terreno, ni paredes laterales, ni piso, porque todo ello es común. % TODO: contenido incompleto o ambiguo en las notas originales (la referencia a la terraza del piso 12 y al piso ocho pisos más abajo no es clara).

\cornell{afectación y estructura jurídica}{
\cq{¿Qué requisitos administrativos y jurídicos exige la afectación?}{Administrativos: plano de construcción y plano de mensura aprobado. Jurídicos: reglamento en escritura pública e inscripción constitutiva en el Registro de la Propiedad.}
\cq{¿Por qué es esencial el plano de mensura?}{Porque sin él no existe la propiedad horizontal: individualiza partes privativas y comunes y fija el porcentual.}
\cq{¿Cuál es el objeto del derecho?}{La unidad funcional (departamento, local, cochera) con dominio exclusivo, unida inescindiblemente a la parte común.}
\cq{¿Qué es el cubo de aire?}{El volumen interior identificado como unidad funcional; el titular no es dueño de paredes, pisos ni techos, que son comunes.}
\cq{¿Qué partes son necesariamente comunes (art.~2041)?}{Terreno, muros exteriores y separativos de unidades, y piso y techo de cada unidad; no pueden recibir otro destino bajo pena de nulidad.}
}{La afectación requiere plano de mensura aprobado y reglamento en escritura inscripto con efecto constitutivo. El objeto es la unidad funcional como cubo de aire con porcentual indiviso del terreno; terreno, muros, pisos y techos son necesariamente comunes.}

% ------------------------------------------------------------
\chapter[El consorcio y las expensas]{El consorcio de propietarios y las expensas}\label{cap:consorcio}
% ------------------------------------------------------------

\section{El consorcio como persona jurídica (art.~2044)}

Cuando se transmite al menos una unidad, es decir, cuando hay al menos dos propietarios diferentes, nace el \textbf{consorcio}.

\begin{definicion}{consorcio (art.~2044 CCyC)}
El conjunto de los propietarios de las unidades funcionales constituye la persona jurídica consorcio. Tiene personería jurídica y se rige por las normas del Capítulo y las del reglamento de propiedad horizontal.
\end{definicion}

Tiene nombre y domicilio, capacidad de adquirir derechos y contraer obligaciones, patrimonio formado por las expensas, una voluntad común que se forma en las asambleas y un órgano ejecutivo, el \textbf{administrador}, representante legal.

\subsection{Responsabilidad de los propietarios (art.~143)}

\begin{ejemplo}{falla del ascensor}
Falla el ascensor y un tercero sufre un daño: demanda al consorcio. Responde la persona jurídica, con su patrimonio de expensas.
\end{ejemplo}

Se cita la postura de Marina Mariani de Vidal, con la que el docente manifiesta coincidir: ser propietario no puede ser un velo de indemnidad para no responder por los daños que puede ocasionar la cosa, en virtud de la responsabilidad objetiva. La responsabilidad de los propietarios es siempre \textbf{subsidiaria e ilimitada}.

\section{Facultades y obligaciones de los propietarios}

\begin{table}[H]
\centering
\small
\begin{tabularx}{\textwidth}{>{\raggedright\arraybackslash}X>{\raggedright\arraybackslash}X}
\toprule
\textbf{Facultades} & \textbf{Obligaciones y prohibiciones} \\
\midrule
Usar y gozar de la unidad (facultades materiales) & Cumplir el reglamento de copropiedad \\
Enajenar la unidad funcional & Sostener económicamente al consorcio (pagar expensas) \\
Constituir derechos reales & Permitir el acceso a la unidad para reparaciones \\
Realizar actos de administración (alquilar) & No destinar la unidad a un fin diferente al admitido por el reglamento \\
Facultades siempre unidas a la parte proporcional de la cosa común & No realizar actividades riesgosas o peligrosas que afecten al edificio \\
\bottomrule
\end{tabularx}
\caption{Facultades y obligaciones de los propietarios}
\end{table}

El art.~2069 faculta al consorcio a accionar para que cesen las infracciones, por la vía procesal más rápida. Por ejemplo, ante una pérdida en un departamento cuyo titular no permite el acceso, con riesgo de inundar pisos inferiores, el consorcio tiene acción, incluso para pedir el desalojo si está ocupado por terceros.

\section{Las expensas}

\subsection{Concepto y fundamento (art.~2048)}

Para Mariani de Vidal, las expensas son la savia que alimenta al consorcio: sin el sostenimiento de los gastos por todos los propietarios, el sistema sucumbe. Por eso la normativa de orden público es rigurosa y el crédito por expensas goza de grandes privilegios.

\begin{ejemplo}{falta de pago de expensas}
Si la mayoría no paga, el consorcio no puede pagar la luz de los servicios comunes: dejan de funcionar el ascensor, la luz del edificio y el agua. Los apuntes proponen los casos de Juan (piso 14, con secuelas de un accidente) y de María (embarazo a término, piso 17, cuyo médico debería subir por la escalera con un respirador eléctrico, sin luz), que muestran por qué el pago es esencial.
\end{ejemplo}

\begin{regla}{art.~2048 CCyC}
Cada propietario debe atender los gastos de conservación y reparación de su propia unidad funcional. Asimismo, debe pagar las expensas comunes ordinarias de administración y las extraordinarias de reparación o mejora que resuelva la asamblea, así como las que, estando determinadas en el reglamento, se devenguen en forma proporcional a su parte indivisa.
\end{regla}

\subsection{Tipos}

\begin{table}[H]
\centering
\small
\begin{tabularx}{\textwidth}{>{\raggedright\arraybackslash}X>{\raggedright\arraybackslash}X}
\toprule
\textbf{Ordinarias} & \textbf{Extraordinarias} \\
\midrule
Regulares y periódicas & No siempre existen \\
Sueldo del encargado, luz, seguridad, seguros, mantenimiento de ascensores & Cambio del sistema de ascensores, portero eléctrico, pintura del edificio \\
No requieren aprobación previa de la asamblea & Deben ser decididas por la asamblea en todos los casos \\
Las liquida el administrador mensualmente & Exceden las facultades ordinarias del administrador \\
\bottomrule
\end{tabularx}
\caption{Expensas ordinarias y extraordinarias}
\end{table}

\subsection{Obligados al pago (arts.~2049 y 2059)}

Están obligados los \textbf{propietarios} (obligación esencial e irrenunciable), los \textbf{poseedores} por cualquier título, los \textbf{usufructuarios} (el nudo propietario también es responsable), el \textbf{adquirente} (adquiere también la deuda existente) y el \textbf{enajenante} (continúa respondiendo por las deudas al tiempo de la enajenación).

\begin{importante}
Los \textbf{locatarios} se obligan ante el consorcio solo si se obligaron a pagar las expensas, en una relación de derecho personal con sus locadores; si no pagan, es motivo de desalojo. El consorcio \textbf{no tiene acción directa} contra ellos.
\end{importante}

\subsection{La expensa como obligación \textit{propter rem}}

La doctrina no acuerda si es \textit{propter rem} en sentido propio; al menos sería una \textit{propter rem sui generis}. La característica de la obligación \textit{propter rem} es que \textbf{sigue a la cosa}, pero en las expensas el enajenante no queda desobligado.

\begin{table}[H]
\centering
\small
\begin{tabularx}{\textwidth}{>{\raggedright\arraybackslash}X>{\raggedright\arraybackslash}X}
\toprule
\textbf{Adquirente} & \textbf{Enajenante} \\
\midrule
Por deudas devengadas \textbf{antes} de la adquisición: responde solo con la cosa & Continúa respondiendo con todo su patrimonio por las deudas anteriores \\
Por deudas devengadas \textbf{durante} su titularidad: responde con la cosa y con todo su patrimonio & Ya no tiene la cosa: responde exclusivamente con su patrimonio \\
\bottomrule
\end{tabularx}
\caption{Responsabilidad del adquirente y del enajenante}
\end{table}

\begin{ejemplo}{deuda al enajenar}
Una unidad (octavo C) tiene una deuda de expensas de \$10.000. Su titular la enajena y la adquiere Juan. Juan responde desde ese momento por los \$10.000 y el enajenante también. El consorcio tiene acción contra ambos; quien pague repetirá contra quien corresponda.
\end{ejemplo}

\begin{importante}
\textbf{Inoponibilidad.} Todo convenio entre enajenante y adquirente sobre las deudas de expensas es inoponible al consorcio (art.~2049), por ejemplo, una cláusula de la escritura por la que el adquirente se hace cargo de la deuda a cambio de un menor precio. Es un instrumento de derecho personal que produce efectos solo entre las partes.
\end{importante}

\subsection{El plenario «Servicios Eficientes c/ Sabrá»}

Plenario de la Cámara Civil (1999), obligatorio para todas las salas por diez años. A partir de él, que siguió aplicándose aun después de cesar su obligatoriedad, el adquirente en subasta, que se libera de impuestos y tasas, no se libera de las deudas de expensas cuando el monto obtenido en la subasta no alcanza a solventarlas. % TODO: contenido incompleto o ambiguo en las notas originales (el pasaje original contiene una frase confusa sobre el precio obtenido).
Al cumplirse diez años, en 2009, el Dr. Jorge Alterini, entonces director de la revista La Ley, solicitó al docente un estudio publicado en esa revista, que constató que la doctrina era aplicada por todas las salas de la Cámara Civil y había llegado a la justicia comercial.

\subsection{Privilegios y protecciones}

\begin{itemize}
    \item \textbf{Prescripción bienal:} el art.~2562 fija dos años para ejecutar el cobro de expensas.
    \item \textbf{Privilegio especial} en ejecuciones individuales: frente a otros acreedores que subastan el inmueble.
    \item No se les opone la afectación del inmueble al \textbf{régimen de protección de la vivienda} (anterior bien de familia).
    \item Se autorizan \textbf{intereses punitorios} además de moratorios.
    \item \textbf{Procedimiento ejecutivo:} el certificado de deuda del administrador (suscripto también por el consejo de propietarios, si existe) es título ejecutivo que habilita embargo y subasta.
\end{itemize}

\cornell{consorcio y expensas}{
\cq{¿Qué es el consorcio y cuándo nace?}{La persona jurídica formada por los propietarios (art.~2044), con nombre, domicilio, patrimonio de expensas, voluntad por asamblea y representante (administrador). Nace con al menos dos propietarios.}
\cq{¿Responden los propietarios por las deudas del consorcio?}{Responde primero el consorcio; la responsabilidad de los propietarios es subsidiaria e ilimitada (postura de Mariani de Vidal).}
\cq{Distinga expensas ordinarias y extraordinarias.}{Ordinarias: regulares, sin aprobación previa, liquidadas mensualmente. Extraordinarias: no siempre existen y las decide la asamblea.}
\cq{¿Quiénes deben pagar expensas y qué pasa con adquirente y enajenante?}{Propietarios, poseedores, usufructuarios, adquirente y enajenante; locatarios solo frente a su locador. Los convenios entre enajenante y adquirente son inoponibles al consorcio.}
\cq{¿Qué dijo el plenario «Servicios Eficientes c/ Sabrá»?}{El adquirente en subasta responde por las expensas en la medida en que el producido no las cubra; la doctrina se mantuvo y llegó a la justicia comercial.}
\cq{¿Qué privilegios tiene el crédito por expensas?}{Prescripción bienal, privilegio especial, inoponibilidad del régimen de vivienda, intereses punitorios y vía ejecutiva con certificado de deuda.}
}{El consorcio es una persona jurídica sostenida por las expensas, cuya obligación recae sobre propietarios, poseedores y aun sobre quien ya enajenó. La expensa es una obligación \textit{propter rem sui generis} muy tutelada, y la doctrina del plenario protege el cobro frente a la subasta.}

% ------------------------------------------------------------
\chapter[Gobierno del consorcio, reglamento, extinción y prehorizontalidad]{Gobierno del consorcio, reglamento, extinción y prehorizontalidad}\label{cap:gobierno-ph}
% ------------------------------------------------------------

\section{El administrador (art.~2067)}

Es el \textbf{representante legal} del consorcio y un \textbf{mandatario}: administra su patrimonio, paga impuestos y sueldos (el encargado, el ayudante), convoca a la asamblea, redacta el orden del día (la identificación clara, concreta y fehaciente de los puntos a tratar), ejecuta sus decisiones y recauda y administra las expensas ordinarias y las extraordinarias decididas. Debe rendir cuentas y tiene la responsabilidad propia de todo mandatario.

\section{Las asambleas (arts.~2058 a 2063)}

La asamblea es el órgano deliberativo y la única vía para resolver asuntos de interés común: forma la voluntad de la persona jurídica. Su decisión obliga a presentes, ausentes y a quienes votaron en contra, siempre que se cumplan las mayorías y se trate el orden del día.

\begin{table}[H]
\centering
\small
\begin{tabularx}{\textwidth}{>{\raggedright\arraybackslash}X>{\raggedright\arraybackslash}X}
\toprule
\textbf{Ordinaria} & \textbf{Extraordinaria} \\
\midrule
Se reúne al menos una vez al año & Se convoca cuando lo exige la situación \\
Trata los temas regulares y la rendición de cuentas & Puede convocarla el administrador, propietarios con más del 5\,\% o autoconvocarse \\
Evalúa presupuestos de gastos o reparaciones & Para autoconvocarse se necesitan los 2/3 de la totalidad de propietarios \\
 & Trata temas urgentes o que no admiten demora \\
\bottomrule
\end{tabularx}
\caption{Asambleas ordinarias y extraordinarias}
\end{table}

\begin{table}[H]
\centering
\small
\begin{tabularx}{\textwidth}{>{\raggedright\arraybackslash}p{0.19\textwidth}>{\raggedright\arraybackslash}X>{\raggedright\arraybackslash}X}
\toprule
\textbf{Mayoría} & \textbf{Descripción} & \textbf{Aplicación} \\
\midrule
Absoluta & Doble cómputo: más del 50\,\% del porcentual y más de la mitad de los votos por unidad funcional & Obras nuevas o mejoras importantes y costosas \\
Unanimidad & 100\,\% de los propietarios & Obras que afecten la estructura en beneficio de uno o pocos; que modifiquen los porcentuales de contribución y propiedad; hipotecar el terreno; cambiar el destino de partes comunes \\
Dos tercios & 2/3 de la totalidad de propietarios & Modificar el reglamento \\
De valor & Más de la mitad del valor del inmueble & Destrucción parcial o total: demolición, venta o reconstrucción \\
\bottomrule
\end{tabularx}
\caption{Tipos de mayorías}
\end{table}

\section{El consejo de copropietarios (art.~2064)}

Órgano \textbf{facultativo} de fiscalización que controla y acompaña la actividad del administrador; cuando existe, suscribe el certificado de deuda.

\section{El reglamento (art.~2056)}

Es un \textbf{contrato de adhesión}: cada nuevo propietario adhiere al contrato que dio vida al sistema. Se redacta en escritura pública, rige derechos y obligaciones e integra el título suficiente de la unidad funcional.

\begin{regla}{art.~2056 CCyC, cláusulas obligatorias}
Son obligatorias para su validez, entre otras: la determinación del terreno; las unidades funcionales y las complementarias; la ubicación de las partes privativas; las partes comunes; el porcentual para el pago de las expensas; las facultades de las asambleas; las mayorías requeridas para las distintas decisiones.
\end{regla}

\section{Extinción del sistema}

La propiedad horizontal recae sobre cosas y carece de suelo propio: si el edificio se destruye, subsiste un \textbf{condominio sobre el terreno}. Ante destrucción o grave deterioro, la \textbf{mayoría de valor} resuelve demolición, reconstrucción o venta. Esa reconstrucción permite al propietario que votó en contra enajenar su parte; si no hay oferentes, el consorcio puede llegar a adquirirla.

\section{Prehorizontalidad (arts.~2070 y 2071)}

Se aplica a los contratos de compraventa de las constructoras sobre unidades a construir y luego someter a propiedad horizontal. Protege a los adquirentes por boleto, titulares de derechos personales, que no pueden adquirir el derecho real porque el inmueble aún no está afectado. Se impone al propietario (generalmente la constructora) la obligación de contratar un \textbf{seguro} en favor de los adquirentes. \textbf{No se aplica} cuando la propiedad horizontal resulta de una partición de condominio, en inmuebles del dominio privado del Estado ni en construcciones con financiamiento público o fideicomisos de organismos oficiales.

\cornell{gobierno del consorcio y prehorizontalidad}{
\cq{¿Qué papel cumple el administrador?}{Representante legal y mandatario: recauda y liquida expensas, paga sueldos e impuestos, convoca asambleas, redacta el orden del día, ejecuta decisiones y rinde cuentas (art.~2067).}
\cq{Distinga asamblea ordinaria y extraordinaria.}{Ordinaria: al menos anual, temas regulares y cuentas. Extraordinaria: cuando la situación lo exige; convocable por el administrador, por propietarios con más del 5\,\% o autoconvocada con 2/3.}
\cq{Mencione las mayorías y un caso de cada una.}{Absoluta (doble cómputo): obras importantes. Unanimidad: cambiar porcentuales o destino de partes comunes. Dos tercios: modificar el reglamento. De valor: demolición, venta o reconstrucción.}
\cq{¿Qué naturaleza tiene el reglamento?}{Contrato de adhesión en escritura pública, parte del título suficiente, con cláusulas obligatorias (art.~2056).}
\cq{¿Cómo se extingue el sistema?}{Por destrucción: queda un condominio sobre el terreno; la mayoría de valor decide demolición, venta o reconstrucción.}
\cq{¿Qué protege la prehorizontalidad y cuándo no se aplica?}{A los adquirentes por boleto antes de la afectación, mediante seguro; no se aplica a particiones de condominio, inmuebles del dominio privado del Estado ni financiamiento público.}
}{La voluntad del consorcio se forma en asambleas con mayorías según la trascendencia del acto y la ejecuta el administrador, bajo fiscalización facultativa del consejo. El reglamento es un contrato de adhesión; ante la destrucción subsiste un condominio sobre el terreno, y la prehorizontalidad protege a quienes compran antes de que nazca el derecho real.}

% ============================================================
\part{Nuevas formas de propiedad: conjuntos inmobiliarios, tiempo compartido y cementerios privados}
% ============================================================

% ------------------------------------------------------------
\chapter[Conjuntos inmobiliarios]{Conjuntos inmobiliarios}\label{cap:conjuntos}
% ------------------------------------------------------------

Quien ejerza como abogado, escribano o asesor inmobiliario se enfrentará con operaciones sobre inmuebles en countries, barrios cerrados, parques industriales o clubes náuticos. Saber bajo qué régimen está organizado cada emprendimiento, derechos reales o personales, determina qué derechos tiene el cliente, qué riesgos enfrenta, si puede hipotecar el bien o afectarlo a la protección de la vivienda y qué ocurre si el desarrollista quiebra.

\section{Marco histórico}

Bajo la denominación de \textbf{conjuntos inmobiliarios}, el CCyC regula formas de organización habituales durante décadas bajo los nombres de barrios cerrados, countries y clubes de campo: nuevas manifestaciones de propiedad inmobiliaria con amplio desarrollo incluso antes de que existiera una ley nacional.

\subsection{Organización bajo derechos personales}

La mayoría se organizaba en el ámbito de los \textbf{derechos personales}: la \textbf{tierra madre} pertenecía en dominio a una persona jurídica (por ejemplo, una sociedad anónima) con dominio absoluto, exclusivo y perpetuo. Lo que adquiría el comprador \textbf{no era un derecho real, sino una acción}, que, por la autonomía de la voluntad, le conferiría derechos sobre un lote.

\subsection{Los riesgos}

Muchos emprendimientos funcionaron sin dificultad, pero otros revelaron la fragilidad del sistema: la tierra madre fue subastada en ejecuciones individuales o la persona jurídica quebró. Cuando se subasta la tierra, se subasta con lo plantado y lo construido, y los compradores perdían todo.

\section{Regulación provincial y debate doctrinal}

Ante la ausencia de una ley nacional, la Provincia de Buenos Aires dictó una norma, inconstitucional o susceptible de serlo, que exigía afectar los nuevos emprendimientos a propiedad horizontal. Sus fundamentos de inconstitucionalidad eran dos: las provincias carecen de competencia para legislar en propiedad, reservada al Congreso, y en la propiedad horizontal clásica el \textbf{terreno es bien común}, incompatible con la naturaleza de un country, pues quien adquiere en un barrio cerrado difícilmente quiera compartir con todos el terreno donde construirá. La Provincia sabía esto, pero no encontraba otro modo de dar seguridad; la sanción de una ley nacional era imprescindible. Marina Mariani de Vidal caracterizó a los inmuebles de countries organizados bajo propiedad horizontal como «grandes castillos construidos sobre arena».

Gran parte de la doctrina sostuvo que debía legislarse en el ámbito de los \textbf{derechos reales}, mediante un derecho real autónomo o una propiedad horizontal especial que comparta lo común y respete la propiedad individual del terreno. El proyecto del Código incluía en el art.~1887 a los conjuntos inmobiliarios como derechos reales, pero su art.~2075 admitía organizarlos también como derechos personales o mixtos, a elección del desarrollista. La ponencia presentada en el Quinto Congreso de Derecho Privado rechazó esa posibilidad: al urbanizador le conviene el ámbito personal, más flexible y menos oneroso porque no exige plano aprobado, escritura pública de cada propietario ni inscripción registral para la oponibilidad.

El adquirente atiende a la ubicación, el precio, el tipo de construcción, los servicios, la distancia a las lagunas, el tiempo de traslado al trabajo y la cercanía de un colegio; en esas cuestiones «el derecho debe guardar un respetuoso silencio». Pero difícilmente comprenda las implicancias de estar bajo derechos personales o reales, y esa función protectora corresponde al derecho: cautelar el orden público y el interés social es función del derecho real, que considera la finitud humana (de allí las sucesiones), la oponibilidad \textit{erga omnes} y la extensión del derecho más allá de la vida del adquirente.

\begin{table}[H]
\centering
\small
\begin{tabularx}{\textwidth}{>{\raggedright\arraybackslash}X}
\toprule
\textbf{Ventajas exclusivas del titular de un derecho real} \\
\midrule
Afectar el inmueble al régimen de protección de la vivienda familiar \\
Derecho de habitación vitalicio y gratuito del cónyuge o conviviente supérstite \\
Persecución y preferencia \\
Constituir nuevos derechos reales, como usufructo e hipoteca \\
Oponibilidad \textit{erga omnes} mediante inscripción registral \\
\bottomrule
\end{tabularx}
\caption{Ventajas del derecho real frente al derecho personal}
\end{table}

\section{Regulación en el CCyC}

Se los incorpora en el art.~1887 como derechos reales y, en el tránsito a la sanción definitiva, se modificó el art.~2075 eliminando la posibilidad de constituirlos como derechos personales. A partir de la sanción (31 de agosto de 2015), todo conjunto inmobiliario debe constituirse exclusivamente como derecho real. % TODO: verificar referencia presente en las fuentes originales (los apuntes transcriben el art.~1887 con un texto que parece una paráfrasis).

\begin{regla}{art.~2075, tercer párrafo}
Todos los conjuntos inmobiliarios preexistentes a la sanción del nuevo ordenamiento que se hubieran establecido como derechos personales, o en los que coexistan derechos personales y derechos reales, se deben adecuar a las previsiones normativas que regulan este derecho real.
\end{regla}

\begin{regla}{art.~2073 CCyC, tipos}
Son conjuntos inmobiliarios los clubes de campo, barrios cerrados o privados, parques industriales, empresariales o náuticos, y todo otro emprendimiento urbanístico independientemente del destino de vivienda permanente o temporaria, laboral, comercial o empresarial que tenga, comprendidos asimismo aquellos que contemplan usos mixtos.
\end{regla}

No es el destino lo que convierte al emprendimiento en conjunto inmobiliario, sino los elementos del art.~2074.

\begin{regla}{art.~2074 CCyC, elementos característicos}
Sin que su enumeración sea taxativa: cerramiento, partes comunes y privativas, estado de indivisión forzosa y perpetua de las partes, lugares y bienes comunes, reglamento que establece órganos de funcionamiento, limitaciones y restricciones a los derechos particulares y régimen disciplinario, obligación de contribuir con los gastos y cargas comunes y entidad con personería jurídica que agrupe a los propietarios de las unidades privativas.
\end{regla}

\subsection{La adecuación de los preexistentes}

El art.~2075 ordena adecuarlos, pero \textbf{no fija plazo} para los organizados como derechos personales o mixtos \textbf{ni sanción} para quienes no lo hagan; coexistirán durante años conjuntos anteriores a 2015 como derechos personales y posteriores como derechos reales. La adecuación requiere planos de mensura aprobados, reglamento, escritura pública de afectación e inscripción registral de cada unidad; es más conveniente y segura, pero de \textbf{elevadísimos costos}, lo que explica la ausencia de plazo y sanción.

\section{Estructura y funcionamiento}

Los conjuntos son una \textbf{propiedad horizontal especial}. En la propiedad horizontal lo esencialmente común es el terreno; en el conjunto, conforme al art.~2077, cada propietario tiene la \textbf{propiedad exclusiva del lote}: así como en aquella el derecho recae sobre un «cubo de aire», aquí recae sobre la tierra.

\begin{regla}{art.~2077 CCyC, unidad funcional}
El derecho de propiedad del propietario en un conjunto inmobiliario se ejerce sobre unidades funcionales privativas, cuya superficie se destina a uso privado. La unidad funcional puede encontrarse construida o en proceso de construcción; debe contar con independencia funcional vinculada a su destino y salida a la vía pública por vías o elementos comunes o en forma directa.
\end{regla}

\begin{table}[H]
\centering
\small
\begin{tabularx}{\textwidth}{>{\raggedright\arraybackslash}X>{\raggedright\arraybackslash}X}
\toprule
\textbf{Propiedad horizontal clásica} & \textbf{Conjunto inmobiliario} \\
\midrule
El terreno es siempre común a todos los copropietarios & Cada propietario tiene propiedad exclusiva de su lote \\
Se requiere edificio ya construido para afectar & Puede estar construido o en proceso de construcción \\
Típicamente, edificios en altura o plantas divididas & Típicamente, lotes individuales con cerramiento perimetral \\
La unidad es un «cubo de aire» & La unidad incluye la tierra (el lote) \\
Destino predominantemente residencial o comercial en altura & Cualquier destino: residencial, industrial, náutico, deportivo \\
Arts.~2037 a 2072 & Arts.~2073 a 2086 (propiedad horizontal especial) \\
\bottomrule
\end{tabularx}
\caption{Propiedad horizontal clásica y conjunto inmobiliario}
\end{table}

\subsection{Partes comunes y privativas}

Son necesariamente comunes las partes del terreno destinadas a vías de circulación, acceso y comunicación; las destinadas a actividades deportivas, recreativas y sociales; y lo que el reglamento determine (art.~2079, que se transcribe en lo esencial). Son privativas las unidades funcionales, construidas o en construcción, con independencia funcional y salida a la vía pública, que por la naturaleza del emprendimiento suele ser indirecta, a través del cerramiento perimetral.

\subsection{Persona jurídica, reglamento y normas urbanísticas}

El conjunto conforma una \textbf{persona jurídica}, como el consorcio, y se constituye con un \textbf{reglamento} que define partes propias y comunes, funcionamiento y asambleas, régimen disciplinario, gastos y su cálculo y el porcentual de contribución. Se aplican además las \textbf{normas urbanísticas} de cada jurisdicción: son normas administrativas, y no puede constituirse un conjunto en cualquier lugar; localización, límites perimetrales y aspectos edilicios son regulados por las normas provinciales y municipales.

\subsection{Facultades, obligaciones y límites del reglamento}

El propietario debe ejercer su derecho conforme a la normativa, el reglamento y las normas de convivencia, respetar los valores constructivos, paisajísticos, arquitectónicos y ecológicos y pagar las expensas en la proporción del reglamento, que puede fijar otras contribuciones por servicios (por ejemplo, polo). Existen restricciones que no se observan en la propiedad horizontal clásica: régimen de invitados, reglas sobre uso de la pileta y canchas, canon adicional por ciertas instalaciones y restricciones arquitectónicas, paisajísticas y ecológicas.

\begin{importante}
El reglamento puede \textbf{restringir} el uso de servicios y espacios comunes, pero \textbf{nunca prohibir} que un propietario venda o transmita su unidad, porque limitaría la transmisión inmobiliaria. Sí puede establecer un \textbf{derecho de preferencia} a favor del consorcio o de los demás propietarios (igualar la oferta antes que un tercero).
\end{importante}

\section{Conclusiones}

Todo conjunto constituido desde el 31 de agosto de 2015 solo puede organizarse como derecho real; los preexistentes personales o mixtos deben adecuarse aunque no haya plazo ni sanción; la adecuación es costosa y compleja, por lo que coexistirán ambas organizaciones durante años; y el profesional debe identificar el régimen de cada emprendimiento para asesorar correctamente.

\cornell{conjuntos inmobiliarios}{
\cq{¿Qué problema tenían los emprendimientos organizados como derechos personales?}{El comprador solo tenía una acción de una sociedad titular de la tierra madre; si la tierra se subastaba o la sociedad quebraba, perdía todo, incluso lo construido.}
\cq{¿Por qué la propiedad horizontal clásica no servía y era inconstitucional la norma bonaerense?}{Porque obligaba a compartir el terreno, contradictorio con un country, y las provincias carecen de competencia para legislar en propiedad.}
\cq{¿Qué dispone el CCyC sobre su constitución y los preexistentes (arts.~1887 y 2075)?}{Desde el 31/8/2015 solo como derecho real; los anteriores deben adecuarse, sin plazo ni sanción, lo que se explica por los elevados costos.}
\cq{¿Qué elementos los caracterizan (arts.~2073 y 2074)?}{No el destino sino el cerramiento, partes comunes y privativas interdependientes, indivisión forzosa, reglamento, contribución a gastos y personería jurídica.}
\cq{Compare con la propiedad horizontal clásica.}{En el conjunto el lote (tierra) es privativo; en la clásica el terreno es común y la unidad es un cubo de aire. El conjunto puede estar en construcción y tener cualquier destino.}
\cq{¿Qué límite tiene el reglamento sobre la transmisión?}{No puede prohibir la venta, pero sí establecer derecho de preferencia.}
}{Los conjuntos inmobiliarios son una propiedad horizontal especial con lote exclusivo, partes comunes, reglamento y persona jurídica. Desde 2015 solo pueden constituirse como derecho real y los anteriores deben adecuarse, aunque sin plazo ni sanción.}

% ------------------------------------------------------------
\chapter[Tiempo compartido]{Tiempo compartido}\label{cap:tc}
% ------------------------------------------------------------

Quien adquiere una semana en un complejo turístico, o una parcela en un cementerio privado, confía en que su adquisición perdurará. El Código intentó dar seguridad jurídica a esas decisiones convirtiéndolas en derechos reales, regulación que todavía genera dudas. El tiempo compartido (arts.~2087 a 2102) se enumera en el art.~1887, inc.~e). Quien ejerce la abogacía enfrenta clientes que lo adquieren sin comprender sus derechos: conocer la figura permite asesorar sobre riesgos y garantías, redactar o revisar contratos de afectación, detectar irregularidades en la comercialización y actuar ante ejecuciones hipotecarias que afecten a los usuarios.

\section{Concepto y objeto}

\begin{definicion}{tiempo compartido (art.~2087 CCyC)}
Hay tiempo compartido si uno o más bienes están afectados al uso periódico y por turno para alojamiento, hospedaje, comercio, turismo, industria u otros fines, que ejercen sobre ellos un derecho que los habilita a su uso y goce. Puede recaer sobre inmuebles o muebles.
\end{definicion}

Se lo asocia con el turismo por ser la modalidad conocida, pero el Código pretende abrir paso a nuevos negocios; por ejemplo, un tomógrafo u otro equipo costoso de medicina nuclear usado por turnos, por un lugar o por otro, durante una semana o período determinado.

\subsection{La modalidad turística en el ámbito de los derechos personales}

La persona adquiere, por un contrato de derecho personal, una \textbf{unidad de tiempo}, generalmente de siete días (o catorce) al año en un complejo vacacional, sin pagar servicios o con un costo muy inferior. Existen sistemas de intercambio dentro del país o en el exterior, como \textbf{Interval}. Quienes lo usan habitualmente no han tenido dificultades, pero como lo lleva adelante una persona jurídica, los inmuebles pueden ser subastados o la persona quebrar, comprometiendo el derecho adquirido por muchos años o a perpetuidad.

\begin{importante}
La regulación como derecho real busca proteger a los adquirentes de los avatares de las personas jurídicas: \textbf{quiebra, inhibición y subasta} de los inmuebles afectados.
\end{importante}

\section{Afectación e inscripción}

Cuando recae sobre inmuebles, el emprendedor debe afectar el inmueble mediante \textbf{escritura pública} con \textbf{inscripción registral} para su oponibilidad \textit{erga omnes} (art.~2089): quien pretenda comprar el inmueble para su propio proyecto, al pedir el informe registral, verá que está afectado a tiempo compartido. Con carácter previo a los anuncios y la comercialización, debe realizarse la afectación en un \textbf{registro de prestadores} y en el Registro de la Propiedad conforme a la ley especial (art.~2092). Según las notas, no existe aún ley especial más que la que regula los tiempos compartidos turísticos, por lo que serían los únicos comercializables; esta ausencia es uno de los mayores obstáculos para la plena operatividad del derecho real.

\section{Sujetos intervinientes}

\begin{table}[H]
\centering
\small
\begin{tabularx}{\textwidth}{>{\raggedright\arraybackslash}p{0.22\textwidth}>{\raggedright\arraybackslash}X}
\toprule
\textbf{Figura} & \textbf{Rol} \\
\midrule
Propietario & Titular del inmueble; puede afectarlo (o el condómino junto a los demás). \\
Emprendedor & Organiza y gestiona el emprendimiento; puede afectar con autorización del propietario, como mandatario. \\
Comercializador & Ofrece y vende las unidades de tiempo al público. \\
Administrador & Gestiona el funcionamiento; puede ser el propio emprendedor o un tercero. \\
Usuarios & Adquirentes del uso periódico: el sujeto que la normativa busca proteger. \\
\bottomrule
\end{tabularx}
\caption{Figuras del tiempo compartido (art.~2100, según las notas)}
\end{table}

Solo puede afectar el inmueble el \textit{dominus}: el propietario o el condómino que, con las partes ideales de todos, conforma el derecho al todo; también el emprendedor con autorización, como mandatario. En definitiva, el titular del dominio, por medio de un administrador.

\section{Interrogantes sobre la naturaleza del derecho}

Según las notas, el art.~2011 establece que al derecho del adquirente «se le aplican las normas de los derechos reales». % TODO: verificar referencia presente en las fuentes originales (el art.~2011 figura en otro capítulo de los apuntes como regla de medianería; la numeración citada aquí puede ser errónea).
De allí surgen interrogantes: ¿se le aplican las normas de los derechos reales, o es un derecho real? ¿Cómo se aplica la norma de un derecho real a un derecho personal, o viceversa? ¿Por qué ha de ser necesariamente uno u otro? Y sobre la titularidad: el propietario tiene un derecho real indiscutido, pero los usuarios, ¿tienen un derecho real sobre cosa propia?

\begin{importante}
\textbf{Interrogante central sin resolver:} el Código dice que al derecho del adquirente se le aplican las normas de los derechos reales, pero no queda claro si es un derecho real en sí mismo. Generará un extenso trabajo de la doctrina y la jurisprudencia.
\end{importante}

\section{Protección de los usuarios}

Al afectar, el inmueble debe estar \textbf{libre de gravámenes} y el propietario, emprendedor, administrador y comercializador \textbf{no deben estar inhibidos} para disponer. El Código admite hipotecas posteriores a la afectación, que \textbf{no son oponibles} a los usuarios (art.~2093): los acreedores del propietario no afectan el derecho de los usuarios.

Las notas señalan dificultades prácticas: si la hipoteca no se paga, el acreedor hipotecario tiene un derecho real con preferencia y puede subastar el inmueble. Es habitual que el titular hipoteque el inmueble precisamente para realizar el emprendimiento, y un cliente puede adquirir una unidad en un complejo importante de la zona norte cuyo inmueble estaba hipotecado.

Los usuarios tienen derechos y deberes: ejercer el uso en la fecha y del modo que correspondan y pagar las expensas necesarias; estas tienen \textbf{cobro ejecutivo} (arts.~2096 y 2097).

\section{Relación de consumo}

El Código incorpora el tiempo compartido como una \textbf{relación de consumo}, algo que ya se había hecho antes por la gran problemática que generaba (art.~2100, según las notas): las relaciones entre propietario, emprendedor, comercializador, administrador y usuarios se rigen por las normas de consumo. % TODO: verificar referencia presente en las fuentes originales (el art.~2100 se cita tanto para la enumeración de sujetos como para la relación de consumo).
Todos deben cumplir las normas nacionales, provinciales o municipales: la regulación del derecho real es nacional, pero el funcionamiento es competencia de las provincias.

\cornell{tiempo compartido}{
\cq{¿Qué es el tiempo compartido (art.~2087)?}{Uso periódico y por turno de uno o más bienes, inmuebles o muebles, para alojamiento, turismo, comercio, industria u otros fines.}
\cq{¿Qué riesgos tenía en derechos personales y qué busca el derecho real?}{Subasta o quiebra de la persona jurídica comprometía el derecho; el derecho real protege contra quiebra, inhibición y subasta.}
\cq{¿Cómo se afecta el inmueble?}{Por escritura pública con inscripción registral (art.~2089), previa a la comercialización y con registro de prestadores (art.~2092).}
\cq{¿Quiénes intervienen?}{Propietario, emprendedor, comercializador, administrador y usuarios.}
\cq{¿Qué interrogante central plantea el Código?}{Si el derecho del adquirente es un derecho real o solo se le aplican las normas de los derechos reales.}
\cq{¿Cómo se protege al usuario frente a hipotecas posteriores?}{Son inoponibles al usuario (art.~2093), aunque subsisten dificultades prácticas si el acreedor ejecuta.}
}{El tiempo compartido es un uso periódico y por turno de bienes afectados mediante escritura inscripta, protegido frente a acreedores posteriores y sometido a normas de consumo. Su naturaleza real sigue siendo discutida y falta la ley especial de registro.}

% ------------------------------------------------------------
\chapter[Cementerios privados]{Cementerios privados}\label{cap:cp}
% ------------------------------------------------------------

Los cementerios privados (arts.~2103 a 2113) son el segundo derecho real de esta parte que plantea, como el tiempo compartido, interrogantes sobre su naturaleza.

\section{Antecedentes}

La organización de los cementerios estuvo ligada a lo religioso: se ubicaban en torno a las iglesias. A partir de \textbf{1821} se produce la \textbf{secularización} y surgen los \textbf{cementerios públicos}, del dominio público municipal, destinados a dar sepultura; el derecho de los particulares sobre las sepulturas es una \textbf{concesión de uso} administrativa. Con el aumento de la población, se saturaron y hace unos cincuenta años surgieron los \textbf{cementerios privados}. Un marco provincial de ejemplo: la Ordenanza General~221 de 1978 y la Ley~9.000/94 autorizan cementerios privados con fines especulativos (de lucro), destinados a la inhumación de restos humanos; el poder de policía mortuoria y el servicio de inhumaciones y cremación es ejercido por las municipalidades.

\section{Concepto y estructura}

\begin{definicion}{cementerio privado (art.~2113 CCyC)}
Se consideran cementerios privados los inmuebles de propiedad privada afectados a la inhumación de restos humanos. Requieren habilitación municipal, cerramiento, y tendrán partes comunes y partes privativas sujetas a indivisión forzosa.
\end{definicion}

El Código pretende regularlos como derecho real, el que mayor seguridad confiere. Aunque la ley no lo aclara, pueden afectar el inmueble el titular de dominio y todos los condóminos que representen el cien por ciento de la propiedad.

\begin{table}[H]
\centering
\small
\begin{tabularx}{\textwidth}{>{\raggedright\arraybackslash}X>{\raggedright\arraybackslash}X}
\toprule
\textbf{Partes comunes (indivisión forzosa)} & \textbf{Partes exclusivas (dominio privativo)} \\
\midrule
Calles y circulaciones internas & Parcelas individuales de sepultura \\
Servicios centrales & Sepulcros y construcciones individuales \\
Capillas y lugares de culto & Derecho a inhumar, exhumar y trasladar restos \\
Instalaciones generales del predio & \\
\bottomrule
\end{tabularx}
\caption{Partes comunes y exclusivas del cementerio privado}
\end{table}

\section{Reglamento y derechos del titular}

\begin{regla}{art.~2105 CCyC, reglamento de afectación y uso}
Debe contener: descripción del inmueble; identificación de las partes y lugares comunes y de los servicios comunes; identificación de las parcelas de dominio exclusivo; modo de pago del canon locativo y los gastos; normas de acceso y circulación de titulares y visitantes; normas administrativas sobre inhumaciones, exhumaciones y cremaciones. Debe existir un registro de inhumados y de sepulturas.
\end{regla}

El \textbf{poder de policía mortuoria} está siempre en cabeza del Estado. El titular de la parcela puede realizar inhumaciones, exhumaciones, reducciones y traslados y construir sepulcros, conforme a las normas del cementerio y el poder de policía mortuoria. El administrador debe garantizar el funcionamiento de las instalaciones y servicios comunes para que quien adquirió pueda usar los servicios contratados.

\section{Inembargabilidad de las parcelas (art.~2110)}

\begin{regla}{art.~2110 CCyC}
Las parcelas exclusivas son inembargables e inejecutables, salvo por saldo del precio de compra, precio de construcción de sepulcros, expensas, tasas, impuestos y contribuciones; en esos casos son embargables y ejecutables.
\end{regla}

\begin{importante}
La inembargabilidad es la diferencia fundamental entre el cementerio privado como derecho real y como derecho personal: protege el destino del bien frente a las deudas del propietario del cementerio, de gran trascendencia para quien adquiere con miras al descanso eterno.
\end{importante}

\section{Relación de consumo y aplicación de normas de derechos reales}

\begin{regla}{arts.~2111 y 2112 CCyC}
La relación entre el propietario y el administrador del cementerio y los usuarios es una relación de consumo sujeta a las normas del derecho del consumidor (art.~2111). Al derecho de sepultura se le aplican las normas de los derechos reales (art.~2112).
\end{regla}

Se plantea el mismo interrogante que en el tiempo compartido: sobre el titular de dominio no hay duda de que tiene un derecho real sobre cosa propia, pero si el derecho de sepultura es un derecho real, ¿por qué habría que aplicarle las normas de los derechos reales?

\begin{importante}
En ambos casos el Código usa la misma fórmula, «se le aplican las normas de los derechos reales», sin afirmar que sea un derecho real en sí mismo. Esta ambigüedad es el principal interrogante que deberá resolver la doctrina y la jurisprudencia. Se trata de regulaciones muy novedosas que requerirán tiempo.
\end{importante}

\cornell{cementerios privados}{
\cq{¿Cómo evolucionó la organización de los cementerios?}{Del vínculo religioso a la secularización de 1821 (cementerios públicos, concesión de uso) y, por saturación, a los cementerios privados con fines de lucro.}
\cq{¿Qué es un cementerio privado (art.~2113)?}{Inmueble privado afectado a inhumación, con habilitación municipal, cerramiento y partes comunes y privativas en indivisión forzosa.}
\cq{¿Qué partes son comunes y cuáles exclusivas?}{Comunes: calles, servicios centrales, capillas e instalaciones. Exclusivas: parcelas y sepulcros, con el derecho de inhumar, exhumar y trasladar.}
\cq{¿Qué contiene el reglamento (art.~2105)?}{Descripción del inmueble, partes y servicios comunes, parcelas exclusivas, canon y gastos, normas de acceso y administrativas, y registro de inhumados y sepulturas.}
\cq{¿Por qué es relevante la inembargabilidad (art.~2110)?}{Protege las parcelas de las deudas del propietario del cementerio, salvo saldo de precio, construcción de sepulcros, expensas, tasas e impuestos.}
\cq{¿Qué ambigüedad comparten tiempo compartido y cementerios?}{Ambos usan la fórmula «se aplican las normas de los derechos reales» sin afirmar que sean derechos reales.}
}{El cementerio privado es un inmueble afectado a la inhumación, con partes comunes en indivisión forzosa y parcelas exclusivas inembargables, bajo poder de policía mortuoria estatal y normas de consumo. Como el tiempo compartido, plantea el interrogante de si es un derecho real o solo se le aplican sus normas.}

\begin{conceptoclave}{consideración conjunta}
Tanto en el tiempo compartido como en los cementerios privados hay muchos interrogantes que no pueden resolverse en esta etapa y que serán despejados con un largo trabajo de la doctrina y la jurisprudencia: son regulaciones novedosas que buscan dar seguridad jurídica convirtiendo en derechos reales adquisiciones antes organizadas en el ámbito personal.
\end{conceptoclave}

% ============================================================
% BACK MATTER
% ============================================================
\backmatter

\chapter*{Referencias mencionadas en las fuentes}
\addcontentsline{toc}{chapter}{Referencias mencionadas en las fuentes}
\markboth{Referencias}{Referencias}

Los apuntes originales no incluyen una bibliografía formal. Se conservan, sin completarlas con información externa, las referencias doctrinales, normativas y jurisprudenciales mencionadas en el texto.

\begin{itemize}
    \item Código Civil y Comercial de la Nación (Ley~26.994), en vigencia desde agosto de 2015; Código Civil de Vélez Sársfield; Ley~13.512 (propiedad horizontal, 1948); Ley~17.711; Constitución Nacional, arts.~14, 17, 18, 19 y 75, inc.~22.
    \item Pacto de San José de Costa Rica, art.~21.
    \item Ordenanza General~221 (1978) y Ley~9.000/94 de la Provincia (cementerios privados).
    \item Mariani de Vidal, Marina: estructura didáctica de la clasificación de los muros; concepción de las expensas como «savia» del consorcio; responsabilidad subsidiaria e ilimitada de los propietarios; caracterización de los countries bajo propiedad horizontal como «grandes castillos construidos sobre arena».
    \item Cámara Nacional de Apelaciones en lo Civil, plenario «Servicios Eficientes S.A. c/ Sabrá, Victor» (1999).
    \item Alterini, Jorge: referido como director de la revista La Ley, que solicitó un estudio sobre la doctrina del plenario diez años después (2009).
    \item Quinto Congreso de Derecho Privado: ponencia sobre conjuntos inmobiliarios.
\end{itemize}
% TODO: verificar referencia presente en las fuentes originales (las fuentes no ofrecen datos bibliográficos completos de estas referencias; no se completaron).

\appendix
\chapter{Ejercicios de evaluación sobre la teoría general}\label{ap:ejercicios}

Las siguientes consignas abarcan la Parte~I y los principios de la Parte~II, organizadas por niveles: básico, intermedio y avanzado.

\section{Nivel básico}

\begin{ejercicio}{1}
\textbf{Consigna.} Defina la obligación según el art.~724 e identifique sus tres elementos.

\textbf{Respuesta.} Relación jurídica en virtud de la cual el acreedor tiene el derecho de exigir del deudor una prestación destinada a satisfacer un interés lícito y, ante el incumplimiento, obtener forzadamente su satisfacción. Elementos: sujeto activo o acreedor, sujeto pasivo o deudor y prestación (dar, hacer o no hacer).
\end{ejercicio}

\begin{ejercicio}{2}
\textbf{Consigna.} ¿Cuál es la diferencia esencial entre un derecho personal y uno real? Explique con un ejemplo.

\textbf{Respuesta.} La naturaleza del vínculo: en el personal, entre dos sujetos y relativo; en el real, directo e inmediato entre titular y cosa, y absoluto (\textit{erga omnes}). Si se presta dinero a Juan (derecho personal), solo Juan debe; si se es dueño de un inmueble (derecho real), toda la comunidad debe respetar esa propiedad.
\end{ejercicio}

\begin{ejercicio}{3}
\textbf{Consigna.} ¿Por qué los derechos personales son «relativos» y los reales «absolutos»?

\textbf{Respuesta.} Los personales son relativos porque el vínculo existe solo entre las partes y el acreedor solo puede exigir a su deudor. Los reales son absolutos porque son oponibles \textit{erga omnes}: toda la comunidad debe respetar el derecho del titular.
\end{ejercicio}

\begin{ejercicio}{4}
\textbf{Consigna.} Defina el derecho real según el art.~1882 e identifique sus elementos.

\textbf{Respuesta.} Poder jurídico de estructura legal que se ejerce directamente sobre su objeto en forma autónoma y que atribuye a su titular las facultades de persecución y preferencia y las demás previstas en el Código. Elementos: el titular y la cosa.
\end{ejercicio}

\begin{ejercicio}{5}
\textbf{Consigna.} ¿Qué integra el patrimonio de una persona según el CCyC?

\textbf{Respuesta.} Los bienes, que incluyen tanto los derechos personales como los reales.
\end{ejercicio}

\begin{ejercicio}{6}
\textbf{Consigna.} ¿Qué establece el art.~16 y por qué es fundamental para los derechos reales?

\textbf{Respuesta.} Que los derechos pueden recaer sobre bienes susceptibles de valor económico y que, cuando son materiales, se llaman cosas (también se aplica a la energía y fuerzas naturales). Define el objeto de los derechos reales: la cosa.
\end{ejercicio}

\section{Nivel intermedio}

\begin{ejercicio}{7}
\textbf{Consigna.} ¿Qué significa el \textit{numerus clausus} y cuál es su consecuencia jurídica?

\textbf{Respuesta.} Solo la ley puede crear derechos reales; los particulares no pueden crear nuevos ni modificar el contenido de los existentes. La consecuencia es la nulidad (art.~1884).
\end{ejercicio}

\begin{ejercicio}{8}
\textbf{Consigna.} Explique la convalidación (art.~1885) y su relación con el \textit{nemo plus iuris}.

\textbf{Respuesta.} El \textit{nemo plus iuris} impide transmitir un derecho mejor o más extenso que el que se tiene, de modo que quien transmite un derecho que no tiene realiza un acto nulo. La convalidación establece que si el transmitente adquiere el derecho con posterioridad, el acto queda convalidado.
\end{ejercicio}

\begin{ejercicio}{9}
\textbf{Consigna.} ¿Qué son el \textit{ius persequendi} y el \textit{ius preferendi} y por qué son características del derecho real?

\textbf{Respuesta.} Persecución: perseguir la cosa en poder de quien se encuentre. Preferencia: ser preferido frente a titulares de derechos reales o personales de fecha posterior. Derivan de la relación directa entre titular y cosa y de su carácter absoluto.
\end{ejercicio}

\begin{ejercicio}{10}
\textbf{Consigna.} Enumere los catorce derechos reales del art.~1887.

\textbf{Respuesta.} Dominio, condominio, propiedad horizontal, conjuntos inmobiliarios, tiempo compartido, cementerio privado, superficie, usufructo, uso, habitación, servidumbre, hipoteca, anticresis y prenda.
\end{ejercicio}

\begin{ejercicio}{11}
\textbf{Consigna.} Diferencie derechos reales sobre cosa propia y sobre cosa ajena, con un ejemplo de cada uno.

\textbf{Respuesta.} Sobre cosa propia, el titular ejerce el derecho sobre su propia cosa (dominio: ser dueño de una casa). Sobre cosa ajena, sobre la cosa de otro, como carga o gravamen para el dueño (hipoteca: el banco tiene un derecho real sobre el inmueble en garantía de un crédito).
\end{ejercicio}

\begin{ejercicio}{12}
\textbf{Consigna.} Diferencie los derechos reales principales de los accesorios, con ejemplos.

\textbf{Respuesta.} Los principales son independientes (dominio, usufructo, servidumbre). Los accesorios son de garantía y existen en función de un crédito (hipoteca, prenda, anticresis); si se extingue el crédito principal, se extingue el derecho de garantía.
\end{ejercicio}

\section{Nivel avanzado}

\begin{ejercicio}{13}
\textbf{Consigna.} ¿Por qué la normativa de derechos reales es de orden público y qué consecuencias tiene?

\textbf{Respuesta.} Porque los derechos reales trascienden el interés particular: su ejercicio afecta a toda la comunidad. Consecuencias: los particulares no pueden crear derechos reales ni modificar su contenido; la violación acarrea nulidad; no se aplica la autonomía de la voluntad como en los contratos.
\end{ejercicio}

\begin{ejercicio}{14}
\textbf{Consigna.} ¿Qué establece el art.~17 y por qué las partes del cuerpo no son objeto de derechos reales?

\textbf{Respuesta.} Que los derechos sobre el cuerpo humano o sus partes no tienen valor comercial sino afectivo, científico, terapéutico, humanitario o social, y se rigen por leyes especiales. No son objeto de derechos reales porque el art.~16 exige valor económico para que un bien sea cosa.
\end{ejercicio}

\begin{ejercicio}{15}
\textbf{Consigna.} Pedro vende un inmueble a María en enero, pero no era propietario. En marzo adquiere la propiedad. ¿Qué ocurre con la venta de enero?

\textbf{Respuesta.} Por el \textit{nemo plus iuris}, Pedro no podía transmitir un derecho que no tenía: la venta nació nula en cuanto a la transmisión del derecho real. Pero por el principio de convalidación (art.~1885), al adquirir Pedro la propiedad en marzo, la transmisión queda convalidada y María adquiere el dominio retroactivamente.
\end{ejercicio}

\end{document}
